% BEGIN INCLUDED SGA7I_English_current_buildable_source_20260823/components/93_expose_VIII_opening_and_introduction.tex
% SOURCE AUTHORITY:
% expose_VIII_body.tex lines 5--52; scan indices 229--231; source folios 218--220.
% Direct scan comparison: physical PDF pages 230--232 at 1100 dpi.
% Hard stop before section 1, source line 55.

\begin{center}
\SGARefAliasTarget{sga7.sga7i.eviii.nroot}SGA 7

\vspace{1em}
\underline{报告 VIII}

\vspace{1em}
\underline{双扩张的补充. 群概形的}

\underline{双扩张的一般性质}

\vspace{1em}
\underline{A. Grothendieck 著}
\end{center}

\subsection*{目录}

\begin{enumerate}[label=\arabic*.]
\setcounter{enumi}{-1}
\item 引言
\item 任意拓扑斯上的若干特殊情形
\item 双扩张所定义的配对
\item 群概形 $(P,Q)$ 以 $\underline G_m$ 为核的双扩张: 一般性质
\item 域上光滑连通群概形的扩张与双扩张
\item 以常值截断无挠群概形为核的扩张与双扩张
\item 以 $\underline G_m$ 的 N\'eron 模型为核的扩张与双扩张
\item 以 $\underline G_m$ 为核的扩张与双扩张的典范延拓
\end{enumerate}

0. \underline{引言}

\SGARefTarget{sga7.sga7i.eviii.n0-1}{0.1}. 本篇继续前一篇开始的双扩张 (biextension) 的一般研究. 首先对任意底拓扑斯的情形作一些补充, 尤其详细展开与双扩张相联系的配对形式体系 (第 2 节); 这些配对在阿贝尔概形理论中的重要性是众所周知的, 值得如此详述. 随后开始研究一般群概形的双扩张, 其中我们最关心的是以乘法群 $\underline G_m$ 为核的双扩张. 不过, 作为技术上的中间步骤, 我们还必须研究以其他类型的群概形为核的双扩张 (第 5、6 节).

对于后文最有用的结果见 \SGARefLink{sga7.sga7i.eviii.n7-1}{7.1}: 它给出一些条件, 使得当 $P$ 和 $Q$ 在一般点为 $\eta$ 的性状上光滑时, $(P_\eta,Q_\eta)$ 以 $\underline G_{m\eta}$ 为核的双扩张能够延拓成 $(P,Q)$ 以 $\underline G_m$ 为核的双扩张, 并且在唯一同构意义下唯一确定. 下一篇专门讨论阿贝尔簇的 N\'eron 模型理论, 这一结果将在其中十分方便. 或许那里的论述本可以更省力, 无须第 VII、VIII 篇这两块沉重的镇纸 (讨论班余下的部分也不再使用它们); 但我们仍觉得这两篇所展开的内容很可能在别处也有用处.

\SGARefTarget{sga7.sga7i.eviii.n0-2}{0.2}. 讨论概形 $S$ 上的群概形时, 约定在 $S$ 的 fppf 拓扑斯中工作. 这是由 $S$ 上局部有限表现的 (并且属于 $\underline U$ 的) 概形构成的 $\underline U$-景赋以 fppf 拓扑后所对应的拓扑斯 (SGA 3 IV 6.3). 读者容易验证, 我们将给出的大部分命题对 $S$ 上的其他拓扑也成立 (并且本质上等价于 fppf 情形的命题), 只要其底景至少包含 $S$ 上局部有限表现的概形; 例如 Zariski 拓扑、平展拓扑和 fpqc 拓扑. 原因之一在于, 由下降理论, $S$ 上局部有限表现的群概形 $P$ 以群概形 $\underline G_m$ 为核的扩张范畴, 在范畴等价意义下不依赖于所选的拓扑; 双扩张的概念也有同样的性质.

选用 fppf 拓扑 (或更细的拓扑, 如 fpqc 拓扑), 而不用 Zariski 拓扑、平展拓扑等较粗拓扑, 技术上的原因在于: 对每个整数 $n>0$, $\underline G_m$ 到自身的同态 $n\cdot\id$ 都是一个 \underline{fppf 满态射} (因为它忠实平坦且有限表现), 但一般既不是 Zariski 满态射, 也不是平展满态射. 注意, 它是平展满态射当且仅当 $n$ 与 $S$ 的剩余特征互素, 而它是 Zariski 满态射当且仅当 $S$ 为空.

\SGARefTarget{sga7.sga7i.eviii.n0-3}{0.3}. 最后, 将群 $\underline G_m$ 换成一个\underline{环面} (torus) $G$, 关于以它为核的扩张或双扩张的大部分命题仍然成立. 因为这样的 $G$ 局部地 (对于 fppf 拓扑, 甚至对于平展拓扑) 同构于 $\underline G_m^r$, 这一本质上是平凡的推广.

% END INCLUDED SGA7I_English_current_buildable_source_20260823/components/93_expose_VIII_opening_and_introduction.tex

% BEGIN INCLUDED SGA7I_English_current_buildable_source_20260823/components/94_expose_VIII_section_1.tex
% R3 explicit revisions R3-VIII-07; see qa/FIDELITY-R3-CHANGES.json.
% SOURCE AUTHORITY:
% expose_VIII_body.tex lines 55--276; scan indices 231--237; source folios 220--226.
% Direct scan comparison: physical PDF pages 232--238 at 1100 dpi.
% Transparent source dispositions:
% - the sentence before (1.1.3) is stated in standard truncation language rather
%   than repeating the printed impossible claim that all Tor_i vanish for i >= 1;
% - the fourth term of (1.1.4) uses Hom(Q,G), as required by the displayed
%   spectral sequence, rather than the printed Hom(Q,P);
% - broken parentheses in (1.3.5), the transition exponent in part (c), and the
%   references to the final member of (1.3.5) are restored from the formulas;
% - in 1.6 the parenthetical Biext argument and the terminal Ext coefficient are
%   corrected to Q and G respectively.
% Hard stop before section 2, source line 279.

\subsection*{\SGARefTarget{sga7.sga7i.eviii.n1}{1}. 任意拓扑斯上的若干特殊情形}

\SGARefTarget{sga7.sga7i.eviii.n1-1}{1.1}. 在本节和下一节中, 与第 VII 篇一样, $T$ 表示一个拓扑斯, $P,Q,G$ 表示 $T$ 的三个阿贝尔群. 因而有典范同构 (VII \SGARefLink{sga7.sga7i.evii.n3-6-5}{3.6.5})
\begin{equation}\tag{\SGARefTarget{sga7.sga7i.eviii.n1-1-1}{1.1.1}}
 \Biext^1(P,Q;G)
 \simeq
 \Ext^1\!\left(P\mathbin{\overset{\mathbb L}{\otimes}}Q,G\right).
\end{equation}
回顾如下典范同构 (称为 Cartan 同构)
\begin{equation}\tag{\SGARefTarget{sga7.sga7i.eviii.n1-1-2}{1.1.2}}
\begin{split}
 \Ext^1\!\left(P\mathbin{\overset{\mathbb L}{\otimes}}Q,G\right)
 &\simeq \Ext^1\!\left(P,R\uHom(Q,G)\right)\\
 &\simeq \Ext^1\!\left(Q,R\uHom(P,G)\right).
\end{split}
\end{equation}
这些同构对三个变元 $P,Q,G$ 都具有函子性.

导出张量积的同调对象为 $\Tor_i(P,Q)$; 对它作标准截断, 由 (\SGARefLink{sga7.sga7i.eviii.n1-1-1}{1.1.1}) 得到如下无限正合列:
\begin{equation}\tag{\SGARefTarget{sga7.sga7i.eviii.n1-1-3}{1.1.3}}
\begin{aligned}
0\longrightarrow{}&
 \Ext^1(P\otimes Q,G)
 \longrightarrow \Biext^1(P,Q;G)
 \longrightarrow \Hom(\Tor_1(P,Q),G)\\
 \longrightarrow{}&
 \Ext^2(P\otimes Q,G)
 \longrightarrow
 \Ext^2\!\left(P\mathbin{\overset{\mathbb L}{\otimes}}Q,G\right)
 \longrightarrow\cdots .
\end{aligned}
\end{equation}
另一方面, 因为 $R\uHom(Q,G)$ 的上同调对象为 $\uExt^i(Q,G)$, (\SGARefLink{sga7.sga7i.eviii.n1-1-2}{1.1.2}) 中的第一个等式给出五项正合列
\begin{equation}\tag{\SGARefTarget{sga7.sga7i.eviii.n1-1-4}{1.1.4}}
\begin{aligned}
0\longrightarrow{}&
 \Ext^1(P,\uHom(Q,G))
 \longrightarrow \Biext^1(P,Q;G)
 \longrightarrow \Hom(P,\uExt^1(Q,G))\\
 \longrightarrow{}&
 \Ext^2(P,\uHom(Q,G))
 \longrightarrow \Ext^2(P,R\uHom(Q,G)),
\end{aligned}
\end{equation}
它对应于通常的谱序列
\[
 \Ext^*(P,R\uHom(Q,G))
 \Longleftarrow
 \Ext^p(P,\uExt^q(Q,G)).
\]
如有需要, 读者可利用\underline{自然的全忠实函子}
\begin{equation}\tag{\SGARefTarget{sga7.sga7i.eviii.n1-1-5}{1.1.5}}
 \operatorname{EXT}(P\otimes Q,G)
 \longrightarrow
 \operatorname{BIEXT}(P,Q;G),
\end{equation}
以及
\begin{equation}\tag{\SGARefTarget{sga7.sga7i.eviii.n1-1-6}{1.1.6}}
 \operatorname{EXT}(P,\uHom(Q,G))
 \longrightarrow
 \operatorname{BIEXT}(P,Q;G),
\end{equation}
分别将 (\SGARefLink{sga7.sga7i.eviii.n1-1-3}{1.1.3}) 和 (\SGARefLink{sga7.sga7i.eviii.n1-1-4}{1.1.4}) 的第一个箭头具体写出. 这里 $\operatorname{EXT}$ 和 $\operatorname{BIEXT}$ 分别表示显然的扩张范畴和双扩张范畴.

当然, 同构
\[
 \Biext^1(P,Q;G)\simeq\Ext^1(Q,R\uHom(P,G))
\]
也给出与 (\SGARefLink{sga7.sga7i.eviii.n1-1-4}{1.1.4}) 类似的正合列, 其中交换了 $P$ 与 $Q$ 的地位. 它的第一个箭头可以用全忠实函子
\begin{equation}\tag{\SGARefTarget{sga7.sga7i.eviii.n1-1-7}{1.1.7}}
 \operatorname{EXT}(Q,\uHom(P,G))
 \longrightarrow
 \operatorname{BIEXT}(P,Q;G)
\end{equation}
具体写出, 此函子与 (\SGARefLink{sga7.sga7i.eviii.n1-1-6}{1.1.6}) 类似.

\SGARefTarget{sga7.sga7i.eviii.n1-2}{1.2}. 假设
\begin{equation}\tag{\SGARefTarget{sga7.sga7i.eviii.n1-2-1}{1.2.1}}
 \Tor_1(P,Q)=0.
\end{equation}
则 (\SGARefLink{sga7.sga7i.eviii.n1-1-3}{1.1.3}) 给出典范同构
\begin{equation}\tag{\SGARefTarget{sga7.sga7i.eviii.n1-2-2}{1.2.2}}
 \Biext^1(P,Q;G)\simeq\Ext^1(P\otimes Q,G).
\end{equation}
例如, 当 $P$ 或 $Q$ 是\underline{平坦} $\mathbb Z$-模时, 这就成立. 更确切地说, 在条件 (\SGARefLink{sga7.sga7i.eviii.n1-2-1}{1.2.1}) 下, 函子 (\SGARefLink{sga7.sga7i.eviii.n1-1-5}{1.1.5}) 是范畴等价.

\SGARefTarget{sga7.sga7i.eviii.n1-3}{1.3}. 假设
\begin{equation}\tag{\SGARefTarget{sga7.sga7i.eviii.n1-3-1}{1.3.1}}
 P\otimes Q=0.
\end{equation}
则 (\SGARefLink{sga7.sga7i.eviii.n1-1-3}{1.1.3}) 给出典范同构
\begin{equation}\tag{\SGARefTarget{sga7.sga7i.eviii.n1-3-2}{1.3.2}}
 \Biext^1(P,Q;G)\simeq\Hom(\Tor_1(P,Q),G).
\end{equation}
我们最关心的满足 (\SGARefLink{sga7.sga7i.eviii.n1-3-1}{1.3.1}) 的情形是存在整数 $n>0$, 使得
\begin{equation}\tag{\SGARefTarget{sga7.sga7i.eviii.n1-3-3}{1.3.3}}
 nP=P,\qquad
 Q=\varinjlim_i {}_{n^i}Q;
\end{equation}
也就是说, $P$ 是 $n$-可除的, 而 $Q$ 是由 $n$ 的幂所零化的子群的极限. 照例记
\[
 T_n(P)=({}_{n^i}P)_{i\geq0}\quad\text{(投射系)},\qquad
 {}_{n^\infty}P=\varinjlim_i{}_{n^i}P.
\]
由 (\SGARefLink{sga7.sga7i.eviii.n1-3-3}{1.3.3}) 得到比 (\SGARefLink{sga7.sga7i.eviii.n1-3-1}{1.3.1}) 更精确的下列公式:
\begin{equation}\tag{\SGARefTarget{sga7.sga7i.eviii.n1-3-4}{1.3.4}}
\begin{split}
 P\otimes Q&=0,\\
 \Tor_1(P,Q)
 &\simeq\varinjlim_i
   \bigl({}_{n^i}P\otimes{}_{n^i}Q\bigr)
 \simeq T_n(P)\otimes Q\\
 &\simeq{}_{n^\infty}P\otimes T_n(Q).
\end{split}
\end{equation}
因而, 同构 (\SGARefLink{sga7.sga7i.eviii.n1-3-2}{1.3.2}) 在这里有更具体的形式
\begin{equation}\tag{\SGARefTarget{sga7.sga7i.eviii.n1-3-5}{1.3.5}}
\begin{split}
 \Biext^1(P,Q;G)
 &\simeq\Hom(T_n(P),D(Q))\\
 &\simeq\Hom(Q,D(T_n(P)))\\
 &\simeq\Hom(T_n(P)\otimes T_n(Q),T_n(G)).
\end{split}
\end{equation}
下面说明这些公式的意义.

\begin{enumerate}[label=\alph*)]
\item 在 (\SGARefLink{sga7.sga7i.eviii.n1-3-4}{1.3.4}) 的归纳系中, 转移同态
\[
 {}_{n^i}P\otimes{}_{n^i}Q
 \longrightarrow
 {}_{n^j}P\otimes{}_{n^j}Q
 \qquad(i\leq j)
\]
可作如下描述: 先利用 $n^{j-i}\otimes\id$ 将第一项与 $\,{}_{n^j}P\otimes{}_{n^i}Q$ 等同, 再使用同态
\[
 \id\otimes\operatorname{inj}:
 {}_{n^j}P\otimes{}_{n^i}Q
 \longrightarrow
 {}_{n^j}P\otimes{}_{n^j}Q,
\]
其中 $\operatorname{inj}:{}_{n^i}Q\longrightarrow{}_{n^j}Q$ 表示包含映射.

\item 按定义, 令
\[
 T_n(P)\otimes Q
 \overset{\mathrm{def}}{=}
 \varinjlim_i\bigl(T_n(P)\otimes{}_{n^i}Q\bigr),
\]
其中
\[
 T_n(P)\otimes M
 \overset{\mathrm{def}}{=}
 {}_{n^i}P\otimes M
 \qquad\text{若 }n^iM=0.
\]
如果还假设 $nQ=Q$, 则可用同样的方法说明 (\SGARefLink{sga7.sga7i.eviii.n1-3-4}{1.3.4}) 中第三项的意义; 此时 $T_n(Q)$ 与 $T_n(P)$ 一样是一个严格投射系. (\SGARefLink{sga7.sga7i.eviii.n1-3-5}{1.3.5}) 的最后一项表示群的投射系之间的 Hom.

\item 字母 $D$ 表示 $\uHom(-,G)$. 对象 $D(Q)$ 定义为群的投射系
\[
 D(Q)=\bigl(D({}_{n^i}Q)\bigr)_{i\geq0}
 =\bigl(\uHom({}_{n^i}Q,G)\bigr)_{i\geq0},
\]
而 (\SGARefLink{sga7.sga7i.eviii.n1-3-5}{1.3.5}) 的第二项是群的投射系之间的同态群. 以下定义给出了下一项更简单的解释:
\[
 D(T_n(P))=\varinjlim_iD({}_{n^i}P),
\]
其中这个归纳系的转移态射由 $T_n(P)$ 的转移态射
\[
 n^{j-i}:{}_{n^j}P\longrightarrow{}_{n^i}P
\]
转置得到.
\end{enumerate}

读者可将 (\SGARefLink{sga7.sga7i.eviii.n1-3-4}{1.3.4}) 所示的 $\Tor_1(P,Q)$ 的计算具体写出, 并由此推出 (\SGARefLink{sga7.sga7i.eviii.n1-3-5}{1.3.5}). 我们只指出, 要用到由 (\SGARefLink{sga7.sga7i.eviii.n1-3-3}{1.3.3}) 的第二个假设推出的关系
\[
 \Tor_p(P,Q)\simeq
 \varinjlim_i\Tor_p(P,{}_{n^i}Q).
\]
当 $p=0$ 时, 这立刻给出 $P\otimes Q=0$, 并把 $\Tor_1(P,Q)$ 的计算归结为各个 $\Tor_1(P,{}_{n^i}Q)$ 及其转移态射的计算. 因此, 当 $Q$ 被某个满足 $mP=P$ 的整数 $m=n^i$ 零化时, 只须具体给出一个同构
\begin{equation}\tag{\SGARefTarget{sga7.sga7i.eviii.n1-3-6}{1.3.6}}
 \Tor_1(P,Q)
 \xleftarrow{\ \sim\ }
 {}_mP\otimes{}_mQ
 \qquad(mP=P,\ mQ=0).
\end{equation}
有几种方法可以做到这一点. 一种是利用正合列
\[
 0\longrightarrow{}_mP\longrightarrow P
 \xrightarrow{\ m\ }P\longrightarrow0.
\]
所对应的 Tor 正合列. 也可以使用下面的同态 (\SGARefLink{sga7.sga7i.eviii.n2-1-1}{2.1.1}); 转移态射正是 (a) 中指定的那些这一事实, 就是 \SGARefLink{sga7.sga7i.eviii.n2-1-11}{2.1.11}. 无论如何, 选取 (\SGARefLink{sga7.sga7i.eviii.n1-3-6}{1.3.6}) 就归结为选取一个符号; 参见 \SGARefLink{sga7.sga7i.eviii.n2-1-10}{2.1.10}.

\SGARefTarget{sga7.sga7i.eviii.n1-4}{1.4}. 假设
\begin{equation}\tag{\SGARefTarget{sga7.sga7i.eviii.n1-4-1}{1.4.1}}
 \uHom(Q,G)=0.
\end{equation}
则 (\SGARefLink{sga7.sga7i.eviii.n1-1-4}{1.1.4}) 给出典范同构
\begin{equation}\tag{\SGARefTarget{sga7.sga7i.eviii.n1-4-2}{1.4.2}}
 \Biext^1(P,Q;G)
 \xrightarrow{\ \sim\ }
 \Hom(P,\uExt^1(Q,G)).
\end{equation}
特别地, 当 $Q$ 是阿贝尔概形且 $G=\underline G_m$ 时就有这种情形; 参见 VII \SGARefLink{sga7.sga7i.evii.n1-3-8}{1.3.8}. 因此, 关于 $P$ 的函子
\[
 P\longmapsto\Biext^1(P,Q;G),
\]
由 $Q^*=\uExt^1(Q,G)$ \underline{可表}. 特别地, 存在 $(Q^*,Q)$ 以 $G$ 为核的典范双扩张 $E^Q$, 它在唯一同构意义下唯一确定, 因为其自同构群为
\[
 \Hom(Q^*,\uHom(Q,G))=0.
\]
这个双扩张, 或与之对称的 $(Q,Q^*)$ 以 $G$ 为核的双扩张 $\,{}^sE^Q$, 可分别称为 $(Q^*,Q)$ 或 $(Q,Q^*)$ 以 $G$ 为核的 \underline{Weil 双扩张}.

\SGARefTarget{sga7.sga7i.eviii.n1-5}{1.5}. 假设
\begin{equation}\tag{\SGARefTarget{sga7.sga7i.eviii.n1-5-1}{1.5.1}}
 \Hom(P,\uExt^1(Q,G))=0.
\end{equation}
特别地, 当 $\uExt^1(Q,G)=0$ 时这就成立; 例如, 若 $Q$ 是有限局部自由群概形或乘法型群 (SGA 3 IX), 并且 $G=\underline G_m$, 则属此情形; 参见下面的 \SGARefLink{sga7.sga7i.eviii.n3-4}{3.4}. 这时 (\SGARefLink{sga7.sga7i.eviii.n1-1-4}{1.1.4}) 给出典范同构
\begin{equation}\tag{\SGARefTarget{sga7.sga7i.eviii.n1-5-2}{1.5.2}}
 \Biext^1(P,Q;G)
 \xleftarrow{\ \sim\ }
 \Ext^1(P,\uHom(Q,G)),
\end{equation}
其中 $\uHom(Q,G)$ 起到 $Q$ (相对于 $G$) 的“对偶群”的作用. 更确切地说, 由 (\SGARefLink{sga7.sga7i.eviii.n1-5-2}{1.5.2}) 可知全忠实函子 (\SGARefLink{sga7.sga7i.eviii.n1-1-6}{1.1.6}) 是\underline{范畴等价}.

\SGARefTarget{sga7.sga7i.eviii.n1-6}{1.6}. 假设两个关系
\begin{equation}\tag{\SGARefTarget{sga7.sga7i.eviii.n1-6-1}{1.6.1}}
 \uHom(P,G)=0,\qquad
 \Hom(P,\uExt^1(Q,G))=0
\end{equation}
同时成立. 分别由 \SGARefLink{sga7.sga7i.eviii.n1-4}{1.4} 和 \SGARefLink{sga7.sga7i.eviii.n1-5}{1.5}, 这两个关系给出同构
\[
 \Biext^1(P,Q;G)
 \xrightarrow{\ \sim\ }
 \Hom(Q,\uExt^1(P,G)),
\]
\[
 \Biext^1(P,Q;G)
 \xleftarrow{\ \sim\ }
 \Ext^1(P,\uHom(Q,G)).
\]
其复合为典范同构
\begin{equation}\tag{\SGARefTarget{sga7.sga7i.eviii.n1-6-2}{1.6.2}}
 \Ext^1(P,\uHom(Q,G))
 \xrightarrow{\ \sim\ }
 \Hom(Q,\uExt^1(P,G))
 \quad\bigl(\simeq\Biext^1(P,Q;G)\bigr),
\end{equation}
它起到对偶公式的作用.

还应注意, (\SGARefLink{sga7.sga7i.eviii.n1-6-1}{1.6.1}) 的第一个关系显然蕴含 $\Hom(P\otimes Q,G)=0$. 因而 $P$ 以 $\uHom(Q,G)$ 为核的交换扩张所成的范畴 $\operatorname{EXT}(P,\uHom(Q,G))$ 是刚性的 (每个对象的每个自同构都是恒等态射), $\operatorname{BIEXT}(P,Q;G)$ 也是如此. 所以一旦知道对象的同构类群, 就知道了这些范畴的结构; 由 (\SGARefLink{sga7.sga7i.eviii.n1-6-2}{1.6.2}), 这个群典范同构于
\[
 \Hom(Q,\uExt^1(P,G)).
\]

映射 (\SGARefLink{sga7.sga7i.eviii.n1-6-2}{1.6.2}) 实际上在所有情形下都有定义, 无须假设 (\SGARefLink{sga7.sga7i.eviii.n1-6-1}{1.6.1}); 也可以不用经过 $\Biext^1(P,Q;G)$, 而用自然配对
\[
 Q\times Q'\longrightarrow G,
\qquad Q'=\uHom(Q,G).
\]
直接描述它. 将该配对解释为同态 $Q\longrightarrow\uHom(Q',G)$, 就得到复合映射
\begin{equation}\tag{\SGARefTarget{sga7.sga7i.eviii.n1-6-3}{1.6.3}}
\begin{split}
 \Ext^1(P,Q')\times Q
 &\longrightarrow
 \Ext^1(P,Q')\times\Hom(Q',G)\\
 &\longrightarrow
 \Ext^1(P,G),
\end{split}
\end{equation}
由此得到的正是 (\SGARefLink{sga7.sga7i.eviii.n1-6-2}{1.6.2}).

下面的 \SGARefLink{sga7.sga7i.eviii.n3-7}{3.7} 给出一个适用 \SGARefLink{sga7.sga7i.eviii.n1-6}{1.6} 的讨论的有用情形.

% END INCLUDED SGA7I_English_current_buildable_source_20260823/components/94_expose_VIII_section_1.tex

% BEGIN INCLUDED SGA7I_English_current_buildable_source_20260823/components/95_expose_VIII_section_2_1_opening_through_2_1_8.tex
% SOURCE AUTHORITY:
% expose_VIII_body.tex lines 279--382; scan indices 237--241; source folios 226--230.
% Direct scan comparison: physical PDF pages 238--242 at 1100 dpi.
% Transparent source dispositions:
% - typewritten letter-o homological indices are normalized to zero;
% - after (2.1.5), the second map is induced by the inclusion {}_nP -> P,
%   not the authority's impossible {}_nP -> Q;
% - the malformed two-argument Tor display in 2.1.8 is set with both arguments.
% Hard stop before 2.1.9, source line 384.

\subsection*{\SGARefTarget{sga7.sga7i.eviii.n2}{2}. 双扩张所定义的配对}

\SGARefTarget{sga7.sga7i.eviii.n2-1}{2.1}. 设 $P$ 和 $Q$ 是 $T$ 的两个阿贝尔群, $n>0$ 是一个整数. 回顾 $\,{}_nP$ 表示乘法 $n\cdot\id_P$ 的核. 我们将定义一个典范同态
\begin{equation}\tag{\SGARefTarget{sga7.sga7i.eviii.n2-1-1}{2.1.1}}
 \alpha=\alpha^n_{P,Q}:
 {}_nP\otimes{}_nQ\longrightarrow\uTor_1(P,Q),
\end{equation}
它将是下列复合映射
\begin{equation}\tag{\SGARefTarget{sga7.sga7i.eviii.n2-1-2}{2.1.2}}
 {}_nP\otimes{}_nQ
 \longrightarrow \uTor_1({}_nP,{}_nQ)
 \longrightarrow \uTor_1(P,Q),
\end{equation}
其中第一个箭头为 $\alpha^n_{{}_nP,{}_nQ}$, 第二个箭头由包含映射 $\,{}_nP\longrightarrow P$ 和 $\,{}_nQ\longrightarrow Q$ 诱导. 因此只须在 $nP=0$ 和 $nQ=0$ 时, 即 $P$ 和 $Q$ 为 $\mathbb Z/n\mathbb Z$-模时定义 (\SGARefLink{sga7.sga7i.eviii.n2-1-1}{2.1.1}).

仅假设 $nQ=0$. 取 $P$ (作为 $\mathbb Z$-模) 的平坦分解 $L_\bullet$, 就得到典范同构
\[
 P\mathbin{\overset{\mathbb L}{\otimes}_{\mathbb Z}}Q
 \simeq L_\bullet\otimes_{\mathbb Z}Q
 \simeq
 (L_\bullet\otimes_{\mathbb Z}\mathbb Z/n\mathbb Z)
 \otimes_{\mathbb Z/n\mathbb Z}Q,
\]
可将它写成导出范畴中的同构
\begin{equation}\tag{\SGARefTarget{sga7.sga7i.eviii.n2-1-3}{2.1.3}}
 P\mathbin{\overset{\mathbb L}{\otimes}_{\mathbb Z}}Q
 \simeq
 \left(P\mathbin{\overset{\mathbb L}{\otimes}_{\mathbb Z}}
       \mathbb Z/n\mathbb Z\right)
 \mathbin{\overset{\mathbb L}{\otimes}_{\mathbb Z/n\mathbb Z}}Q.
\end{equation}
由此得到同调的 K\"unneth 型谱序列 (参见 SGA 1 IV 5.2)
\begin{equation}\tag{\SGARefTarget{sga7.sga7i.eviii.n2-1-4}{2.1.4}}
 \uTor^{\mathbb Z}_*(P,Q)
 \Longleftarrow
 E^2_{p,q}
 =
 \uTor^{\mathbb Z/n\mathbb Z}_p
 \bigl(\uTor^{\mathbb Z}_q(P,\mathbb Z/n\mathbb Z),Q\bigr).
\end{equation}
还要回顾典范同构
\[
 \uTor^{\mathbb Z}_0(P,\mathbb Z/n\mathbb Z)
 =P_n\mathrel{\mathop:}=P/nP,
 \qquad
 \uTor^{\mathbb Z}_1(P,\mathbb Z/n\mathbb Z)\simeq{}_nP,
\]
以及
\[
 \uTor^{\mathbb Z}_i(P,\mathbb Z/n\mathbb Z)=0
 \qquad\text{对于 }i\geq2.
\]
它们可由 $(\mathbb Z/n\mathbb Z)_T$ 的典范平坦分解
\[
\begin{tikzcd}[column sep=large,row sep=large]
 0\arrow[r]&
 \mathbb Z_T\arrow[r,"n"]\arrow[d,equal]&
 \mathbb Z_T\arrow[r]\arrow[d,equal]&
 0\\
 &L_1&L_0&
\end{tikzcd}
\]
计算得出, 其中增广 $L_0\longrightarrow(\mathbb Z/n\mathbb Z)_T$ 将 $1$ 映至 $1$. 因此谱序列 (\SGARefLink{sga7.sga7i.eviii.n2-1-4}{2.1.4}) 本质上归结为正合列
\begin{equation}\tag{\SGARefTarget{sga7.sga7i.eviii.n2-1-5}{2.1.5}}
\begin{split}
0\longrightarrow{}&
 \uTor^{\mathbb Z/n\mathbb Z}_2(P_n,Q)
 \longrightarrow{}_nP\otimes Q
 \xrightarrow{\alpha}\uTor^{\mathbb Z}_1(P,Q)\\
\longrightarrow{}&
 \uTor^{\mathbb Z/n\mathbb Z}_1(P_n,Q)
 \longrightarrow0,
\end{split}
\end{equation}
其中 $P_n=P\otimes_{\mathbb Z}\mathbb Z/n\mathbb Z=P/nP$. 所求同态 $\alpha$ 就是这个正合列中的边缘同态 $E^2_{0,1}\longrightarrow\uTor^{\mathbb Z}_1$. 由于该正合列显然对 $P$ 和 $Q$ 具有函子性, 关于 $P$ 的函子性表明, 这个同态确实是复合映射
\[
 {}_nP\otimes Q
 \longrightarrow\uTor_1({}_nP,Q)
 \longrightarrow\uTor_1(P,Q),
\]
其中第一个同态是刚定义的 $\alpha$, 只是将 $(P,Q)$ 换成 $({}_nP,Q)$, 第二个同态由包含映射 $\,{}_nP\longrightarrow P$ 诱导.

\SGARefTarget{sga7.sga7i.eviii.n2-1-6}{2.1.6}. 由此对任意 $P$ 和 $Q$ 定义了满足分解条件 (\SGARefLink{sga7.sga7i.eviii.n2-1-2}{2.1.2}) 的同态 (\SGARefLink{sga7.sga7i.eviii.n2-1-1}{2.1.1}). 此外, 当 $nQ=0$ 时, 正合列 (\SGARefLink{sga7.sga7i.eviii.n2-1-5}{2.1.5}) 表明, 只要 $P_n$ 在 $\mathbb Z/n\mathbb Z$ 上平坦, $\alpha$ 就是同构. 特别地, 若 $P_n=0$, 即若 $nP=P$ 且 $P$ 是 $n$-可除的, 或者若 $P$ 本身是平坦 $\mathbb Z/n\mathbb Z$-模, 就有此性质. 若 $Q$ 是平坦 $\mathbb Z/n\mathbb Z$-模, 也同样成立.

\SGARefTarget{sga7.sga7i.eviii.n2-1-7}{2.1.7}. 下面具体计算 (\SGARefLink{sga7.sga7i.eviii.n2-1-5}{2.1.5}) 中的同态 $\alpha$. 为了确定它, 只须对 $Q$ 的每个截面 $q$ (在 $T$ 的某个对象 $S$ 上; 为简单起见取终对象), 确定相应的同态
\[
 {}_nP\longrightarrow\uTor^{\mathbb Z}_1(P,Q).
\]
截面 $q$ 可视为 $(\mathbb Z/n\mathbb Z)_T$ 的截面 $1$ 在唯一确定的同态
\[
 u_q:(\mathbb Z/n\mathbb Z)_T\longrightarrow Q.
\]
下的像. 因而 $\alpha$ 关于 $Q$ 的函子性把问题归结为 $Q=(\mathbb Z/n\mathbb Z)_T$ 且 $q$ 是其单位截面的情形. 此时, 按照 (\SGARefLink{sga7.sga7i.eviii.n2-1-4}{2.1.4}) 之后的说明, 将
\[
 \uTor^{\mathbb Z}_1(P,Q)
 =
 \uTor^{\mathbb Z}_1(P,\mathbb Z/n\mathbb Z)
\]
与 $\,{}_nP$ 等同, 则所讨论的同态就是恒等同态. 等价地, 边缘同态
\[
 E^2_{0,q}
 =
 \uTor^{\mathbb Z}_q(P,\mathbb Z/n\mathbb Z)
 \otimes_{\mathbb Z/n\mathbb Z}\mathbb Z/n\mathbb Z
 \longrightarrow
 \uTor^{\mathbb Z}_q(P,\mathbb Z/n\mathbb Z)
\]
是恒等同态. 这由定义立刻可得. 将 $\mathbb Z$ 及其商 $\mathbb Z/n\mathbb Z$ 换成 $T$ 上的任意环 $A$ 及任意商环 $A/I$, 类似的论断仍然成立.

\SGARefTarget{sga7.sga7i.eviii.n2-1-8}{2.1.8}. 还可以用下面的方法确定 (\SGARefLink{sga7.sga7i.eviii.n2-1-1}{2.1.1}) 的同态 $\alpha$: 对每一对截面, 即 $\,{}_nP$ 的截面 $p$ 和 $\,{}_nQ$ 的截面 $q$ (在 $T$ 的某个对象 $S$ 上, 为简化记号仍取终对象), 描述 $(p,q)$ 所确定的 $\uTor^{\mathbb Z}_1(P,Q)$ 的截面. 与上面一样, 由关于 $P$ 和 $Q$ 的函子性, 问题归结为
\[
 P=Q=(\mathbb Z/n\mathbb Z)_T,
\]
且 $p$ 与 $q$ 为单位截面的情形. 因而只须确定
\[
 \alpha:
 (\mathbb Z/n\mathbb Z)_T\otimes(\mathbb Z/n\mathbb Z)_T
 \longrightarrow
 \uTor^{\mathbb Z}_1
 \bigl((\mathbb Z/n\mathbb Z)_T,(\mathbb Z/n\mathbb Z)_T\bigr).
\]
\SGARefLink{sga7.sga7i.eviii.n2-1-7}{2.1.7} 中的描述表明, 这就是同构
\begin{equation}\tag{\SGARefTarget{sga7.sga7i.eviii.n2-1-8-1}{2.1.8.1}}
 (\mathbb Z/n\mathbb Z)_T
 \xrightarrow{\sim}
 \uTor^{\mathbb Z}_1
 \bigl((\mathbb Z/n\mathbb Z)_T,(\mathbb Z/n\mathbb Z)_T\bigr)
\end{equation}
它由 $\uTor_1$ 中\underline{第一个}因子 $P=(\mathbb Z/n\mathbb Z)_T$ 的典范平坦分解
\begin{equation}\tag{\SGARefTarget{sga7.sga7i.eviii.n2-1-8-2}{2.1.8.2}}
 0\longrightarrow\mathbb Z_T
 \xrightarrow{n}\mathbb Z_T
 \longrightarrow0
\end{equation}
得到.

% END INCLUDED SGA7I_English_current_buildable_source_20260823/components/95_expose_VIII_section_2_1_opening_through_2_1_8.tex

% BEGIN INCLUDED SGA7I_English_current_buildable_source_20260823/components/96_expose_VIII_section_2_1_9_through_2_1_12.tex
% SOURCE AUTHORITY:
% expose_VIII_body.tex lines 384--547; scan indices 241--244; source folios 230--233.
% Direct scan comparison: physical PDF pages 242--245 at 1100 dpi.
% Transparent source dispositions:
% - typewritten letter-o complex indices are normalized to zero;
% - broken internal references 2.11.1/2.11.2 are restored as 2.1.11.1/2.1.11.2;
% - the upper vertex in the proof of 2.1.11 is set as P/nP, as required by
%   the identification in the prose and the four maps.
% Hard stop before 2.1.13, source line 552.

\SGARefTarget{sga7.sga7i.eviii.n2-1-9}{2.1.9}. (\SGARefLink{sga7.sga7i.eviii.n2-1-1}{2.1.1}) 中的定义并不对称. 在上面的构造中交换 $P$ 和 $Q$ 的地位, 可以定义类似的同态
\begin{equation}\tag{\SGARefTarget{sga7.sga7i.eviii.n2-1-9-1}{2.1.9.1}}
 \beta^n_{P,Q}:{}_nP\otimes{}_nQ\longrightarrow\uTor_1(P,Q),
\end{equation}
先归结为 $nP=0$ 的情形, 再使用 $Q$ \underline{的}平坦分解. 等价地, 可以借助 $\alpha$ 定义 $\beta$, 要求下图交换:
\begin{equation}\tag{\SGARefTarget{sga7.sga7i.eviii.n2-1-9-2}{2.1.9.2}}
\begin{tikzcd}[column sep=huge,row sep=huge]
 {}_nP\otimes{}_nQ
   \arrow[r,"\sim"]
   \arrow[d,"{\beta^n_{P,Q}}"']&
 {}_nQ\otimes{}_nP
   \arrow[d,"{\alpha^n_{Q,P}}"]\\
 \uTor_1(P,Q)\arrow[r,"\sim"]&
 \uTor_1(Q,P).
\end{tikzcd}
\end{equation}
这里水平箭头是对称同构. 按照这个约定, $\alpha$ 与 $\beta$ 的关系为
\begin{equation}\tag{\SGARefTarget{sga7.sga7i.eviii.n2-1-9-3}{2.1.9.3}}
 \beta^n_{P,Q}=-\alpha^n_{P,Q}.
\end{equation}
换言之:

\medskip
\noindent\underline{引理 \SGARefTarget{sga7.sga7i.eviii.n2-1-10}{2.1.10}.}
\underline{设 $P$ 和 $Q$ 是两个交换群, $n>0$ 是一个整数. 则下图反交换, 其中 $\alpha$ 定义于 (\SGARefLink{sga7.sga7i.eviii.n2-1-1}{2.1.1}):}
\begin{equation}\tag{\SGARefTarget{sga7.sga7i.eviii.n2-1-10-1}{2.1.10.1}}
\begin{tikzcd}[column sep=huge,row sep=huge]
 {}_nP\otimes{}_nQ
   \arrow[r,"\sim"]
   \arrow[d,"{\alpha^n_{P,Q}}"']&
 {}_nQ\otimes{}_nP
   \arrow[d,"{\alpha^n_{Q,P}}"]\\
 \uTor_1(P,Q)
   \arrow[r,"\sim"]
   \arrow[ur,phantom,"-" description]&
 \uTor_1(Q,P).
\end{tikzcd}
\end{equation}
\underline{水平箭头为对称同构.}

\medskip
\noindent\underline{证明.}
\SGARefLink{sga7.sga7i.eviii.n2-1-8}{2.1.8} 的方法立刻把问题归结为 $P=Q=(\mathbb Z/n\mathbb Z)_T$ 的情形. 此时须证明: 分别用 (\SGARefLink{sga7.sga7i.eviii.n2-1-8-2}{2.1.8.2}) 分解 $\uTor_1$ 的第一个和第二个变元所得到的两个同构 (\SGARefLink{sga7.sga7i.eviii.n2-1-8-1}{2.1.8.1}) 相差一个符号 (等价地, $\uTor^{\mathbb Z}_1(\mathbb Z/n\mathbb Z,\mathbb Z/n\mathbb Z)$ 上的对称同构是 $-\id$). 显然可以假设 $T$ 是点拓扑斯, 于是所讨论的是通常的交换群. 设 $L_\bullet$ 是复形
\[
 L_\bullet=\bigl(0\longrightarrow\mathbb Z
 \xrightarrow{n}\mathbb Z\longrightarrow0\bigr),
\]
把它视为 $P=Q=\mathbb Z/n\mathbb Z$ 的平坦分解. 使用 $L_\bullet\otimes L_\bullet$ 来计算
\[
 \Tor_1(P,Q)=\Tor_1(\mathbb Z/n\mathbb Z,\mathbb Z/n\mathbb Z).
\]
与 $H_1(L_\bullet\otimes\mathbb Z/n\mathbb Z)$ 的同构 (它定义 $\alpha$) 以及与 $H_1(\mathbb Z/n\mathbb Z\otimes L_\bullet)$ 的同构 (它定义 $\beta$), 分别由 $L_\bullet\otimes L_\bullet$ 的第二个和第一个因子的增广 $L_\bullet\longrightarrow\mathbb Z/n\mathbb Z$ 诱导. 由此立即可见, $\alpha(1,1)$ 由闭链
\[
 -1_0\otimes1_1+1_1\otimes1_0
 \in(L_\bullet\otimes L_\bullet)_1
 =L_0\otimes L_1+L_1\otimes L_0,
\]
表示, 其中 $1_0$ 和 $1_1$ 分别表示 $L_0$ 和 $L_1$ 中的元素 $1$, 而 $\beta(1,1)$ 由 $1_0\otimes1_1-1_1\otimes1_0$ 表示. 这就证明了 \SGARefLink{sga7.sga7i.eviii.n2-1-10}{2.1.10}.

要结束关于配对 $\alpha$ 的这些一般讨论, 还须说明 $n$ 变化时各个同态 $\alpha^n_{P,Q}$ 之间的关系. 主要结果如下.

\medskip
\noindent\underline{引理 \SGARefTarget{sga7.sga7i.eviii.n2-1-11}{2.1.11}.}
\underline{设 $P,Q$ 是 $T$ 的两个交换群, $n,n'$ 是两个正整数, 使 $n'$ 是 $n$ 的倍数, 记 $n'=mn$. 则下图交换:}
\begin{equation}\tag{\SGARefTarget{sga7.sga7i.eviii.n2-1-11-1}{2.1.11.1}}
\begin{tikzcd}[column sep=large,row sep=large]
 &{}_{n'}P\otimes{}_nQ
   \arrow[ld,"P_m\otimes\id_{{}_nQ}"']
   \arrow[rd,"\id_{{}_{n'}P}\otimes i_m"]&\\
 {}_nP\otimes{}_nQ
   \arrow[rd,"{\alpha^n_{P,Q}}"']&&
 {}_{n'}P\otimes{}_{n'}Q
   \arrow[ld,"{\alpha^{n'}_{P,Q}}"]\\
 &\uTor_1(P,Q)&
\end{tikzcd}
\end{equation}
\underline{这里 $P_m:{}_{n'}P\longrightarrow{}_nP$ 是由 $m\cdot\id_P$ 诱导的同态, $i_m:{}_nQ\longrightarrow{}_{n'}Q$ 是包含映射. 类似地, 与 (\SGARefLink{sga7.sga7i.eviii.n2-1-11-1}{2.1.11.1}) 对称的图 (\SGARefTarget{sga7.sga7i.eviii.n2-1-11-2}{2.1.11.2}) 也是交换的; 除了将最上方的顶点 $\,{}_{n'}P\otimes{}_nQ$ 换为 $\,{}_nP\otimes{}_{n'}Q$ 以外, 它具有相同的顶点.}

\medskip
只须证明第一个断言; 根据 \SGARefLink{sga7.sga7i.eviii.n2-1-10}{2.1.10}, 第二个断言由对称性得出. 按照 \SGARefLink{sga7.sga7i.eviii.n2-1-8}{2.1.8} 的方法, 问题归结为点拓扑斯上的情形, 其中 $P$ 和 $Q$ 是通常的交换群, 且
\[
 P=\mathbb Z/n'\mathbb Z,\qquad Q=\mathbb Z/n'\mathbb Z.
\]
事实上, 对 $P$ 我们只用到假设 $n'P=0$. 用分解
\[
 0\longrightarrow\mathbb Z\xrightarrow{n'}\mathbb Z
 \longrightarrow0
\]
来分解 $Q$, 将 $\Tor_1(P,Q)$ 与 $\,{}_{n'}P=P$ 等同. 类似地, 用
\[
 0\longrightarrow\mathbb Z\xrightarrow n\mathbb Z
 \longrightarrow0
\]
分解 $\,{}_nQ\simeq\mathbb Z/n\mathbb Z$, 将 $\Tor_1({}_nP,{}_nQ)$ 与 $\,{}_nP$ 等同. 后一个分解的增广将 $1_0$ 映到 $m\cdot1_Q$, 其中 $1_Q$ 表示 $Q=\mathbb Z/n'\mathbb Z$ 的元素 $1$. 在这些等同下, 由包含映射 $\,{}_nP\hookrightarrow P$ 和 $\,{}_nQ\hookrightarrow Q$ 诱导的同态
\[
\begin{tikzcd}[column sep=large,row sep=large]
 \Tor_1({}_nP,{}_nQ)\arrow[r]\arrow[d,no head,"\wr"']&
 \Tor_1(P,Q)\arrow[d,no head,"\wr"']\\
 {}_nP&{}_{n'}P=P
\end{tikzcd}
\]
恰是自然包含映射 $\,{}_nP\hookrightarrow{}_{n'}P$.

因此, 利用 $\,{}_nQ$ 的生成元 $m\cdot1_Q$ 将 (\SGARefLink{sga7.sga7i.eviii.n2-1-11-1}{2.1.11.1}) 的上方顶点与 $P/nP$ 等同, 该生成元将 $\,{}_nQ$ 与 $\mathbb Z/n\mathbb Z$ 等同. 图就变为
\[
\begin{tikzcd}[column sep=large,row sep=large]
 &P/nP
   \arrow[ld,"x\longmapsto mx"']
   \arrow[rd,"x\longmapsto mx"]&\\
 {}_nP\arrow[rd,"x\longmapsto x"']&&
 P\arrow[ld,"x\longmapsto x"]\\
 &{}_{n'}P&
\end{tikzcd}
\]
它显然交换.

\medskip
\noindent\underline{推论 \SGARefTarget{sga7.sga7i.eviii.n2-1-12}{2.1.12}.}
\underline{沿用 \SGARefLink{sga7.sga7i.eviii.n2-1-11}{2.1.11} 的记号, 下图交换:}
\begin{equation}\tag{\SGARefTarget{sga7.sga7i.eviii.n2-1-12-1}{2.1.12.1}}
\begin{tikzcd}[column sep=huge,row sep=huge]
 {}_{n'}P\otimes{}_{n'}Q
   \arrow[r,"m\otimes m"]
   \arrow[d,"{\alpha^{n'}_{P,Q}}"']&
 {}_nP\otimes{}_nQ
   \arrow[d,"{\alpha^n_{P,Q}}"]\\
 {}_{n'}\uTor_1(P,Q)
   \arrow[r,"m"]&
 {}_n\uTor_1(P,Q).
\end{tikzcd}
\end{equation}
\underline{在水平箭头的描述中, $m$ 分别表示对于 $E=P,Q$ 和 $\uTor_1(P,Q)$, 由 $m\cdot\id_E$ 诱导的同态 $\,{}_{n'}E\longrightarrow{}_nE$.}

\medskip
事实上, 由 \SGARefLink{sga7.sga7i.eviii.n2-1-11}{2.1.11}, 对于 $x\in{}_{n'}P(S)$ 和 $y\in{}_{n'}Q(S)$, 有
\[
 \alpha^n(mx\otimes my)=\alpha^{n'}(x\otimes my).
\]
右端等于 $m\alpha^{n'}(x\otimes y)$, 因为 $x\otimes my=m(x\otimes y)$. 这就证明了 (\SGARefLink{sga7.sga7i.eviii.n2-1-12-1}{2.1.12.1}) 的交换性.

% END INCLUDED SGA7I_English_current_buildable_source_20260823/components/96_expose_VIII_section_2_1_9_through_2_1_12.tex

% BEGIN INCLUDED SGA7I_English_current_buildable_source_20260823/components/97_expose_VIII_section_2_1_13_and_remark_2_1_14.tex
% SOURCE AUTHORITY:
% expose_VIII_body.tex lines 552--675; scan indices 245--248; source folios 234--237.
% Direct scan comparison: physical PDF pages 246--249 at 1100 dpi.
% Diagram 2.1.14.6: six targeted vector crops at 5000 dpi, covering every
% node, arrow, label, attachment, and terminal punctuation.
% Transparent source dispositions:
% - the authority's self-contradictory phrase "injective isomorphism" in
%   2.1.14 is rendered as "injective homomorphism";
% - in the multiplication-by-pq map following (2.1.14.6), the authority
%   omits pq before the second ideal in the target numerator; it is restored
%   as pq a \cap pq b, as required by multiplication by pq and by the diagram.
% Hard stop before 2.2, source line 676.

\SGARefTarget{sga7.sga7i.eviii.n2-1-13}{2.1.13}. 设 $\ell$ 为素数. 对每个交换群 $E$, 用 $T_\ell(E)$ 表示投射系
\begin{equation}\tag{\SGARefTarget{sga7.sga7i.eviii.n2-1-13-1}{2.1.13.1}}
 T_\ell(E)=({}_{\ell^\nu}E)_{\nu\geq0},
\end{equation}
它显然函子性地依赖于 $E$. 对 $n,n'$ 形如 $\ell^\nu,\ell^{\nu'}$ 且 $\nu'\geq\nu$ 的情形应用 \SGARefLink{sga7.sga7i.eviii.n2-1-12}{2.1.12} 的相容性, 可知 $\nu$ 变化时的同态 $\alpha^{\ell^\nu}_{P,Q}$ 定义投射系的同态, 它对 $P,Q$ 具有函子性:
\begin{equation}\tag{\SGARefTarget{sga7.sga7i.eviii.n2-1-13-2}{2.1.13.2}}
 \alpha^{(\ell)}_{P,Q}:
 T_\ell(P)\otimes T_\ell(Q)
 \longrightarrow T_\ell\bigl(\uTor_1(P,Q)\bigr).
\end{equation}
显然, 对每个素数 $\ell$ 知道这些同态, 等价于对每个整数 $n$ 知道 $\alpha^n_{P,Q}$. 对每个 $\ell$, \SGARefLink{sga7.sga7i.eviii.n2-1-10}{2.1.10} 中的反对称关系此时变成类似的 $\ell$ 进反对称关系, 表现为下列图的反交换性
\begin{equation}\tag{\SGARefTarget{sga7.sga7i.eviii.n2-1-13-3}{2.1.13.3}}
\begin{tikzcd}[column sep=huge,row sep=huge]
 T_\ell(P)\otimes T_\ell(Q)
   \arrow[r]
   \arrow[d,"{\alpha^{(\ell)}_{P,Q}}"']
   \arrow[rd,phantom,"-" description]&
 T_\ell(Q)\otimes T_\ell(P)
   \arrow[d,"{\alpha^{(\ell)}_{Q,P}}"]\\
 T_\ell\bigl(\uTor_1(P,Q)\bigr)
   \arrow[r]&
 T_\ell\bigl(\uTor_1(Q,P)\bigr).
\end{tikzcd}
\end{equation}
这里水平箭头是对称箭头.

\medskip
\noindent\underline{注 \SGARefTarget{sga7.sga7i.eviii.n2-1-14}{2.1.14}.}
同态 (\SGARefLink{sga7.sga7i.eviii.n2-1-1}{2.1.1}) 的构造及关于它的结果 \SGARefLink{sga7.sga7i.eviii.n2-1-10}{2.1.10} 和 \SGARefLink{sga7.sga7i.eviii.n2-1-11}{2.1.11}, 可推广到给定 $T$ 上的交换环 $\underline A$ 以及 $\underline A$ 的两个理想 $\underline a$ 和 $\underline b$ 的情形. 在上面讨论的情形中, 两个理想都是 $T$ 上的常值环 $\mathbb Z_T$ 的理想 $n\mathbb Z_T$. 这时可定义关于 $\underline A$-模 $P,Q$ 具有函子性的同态
\begin{equation}\tag{\SGARefTarget{sga7.sga7i.eviii.n2-1-14-1}{2.1.14.1}}
 {}_{\underline a}P\otimes{}_{\underline b}Q
 \xrightarrow{\ \alpha\ }
 \uHom_{\underline A}
 \left(
   \frac{\underline a\cap\underline b}
        {\underline a\mathbin{\cdot}\underline b},
   \uTor^{\underline A}_1(P,Q)
 \right),
\end{equation}
它是同态
\[
\begin{split}
 {}_{\underline a}P\otimes{}_{\underline b}Q
 &\longrightarrow
 \uHom(\underline A/\underline a,P)
 \otimes
 \uHom(\underline A/\underline b,Q)\\
 &\longrightarrow
 \uHom\!\left(
   \uTor_1(\underline A/\underline a,\underline A/\underline b),
   \uTor_1(P,Q)
 \right),
\end{split}
\]
的复合, 这里利用了熟知的典范同构
\begin{equation}\tag{\SGARefTarget{sga7.sga7i.eviii.n2-1-14-2}{2.1.14.2}}
 \uTor_1(\underline A/\underline a,\underline A/\underline b)
 \simeq
 \frac{\underline a\cap\underline b}
      {\underline a\mathbin{\cdot}\underline b}.
\end{equation}
事实上, \SGARefLink{sga7.sga7i.eviii.n2-1-10}{2.1.10} 中的符号“隐藏”在这个同构的选择中, 取决于使用 $\underline A/\underline a$ 的分解 $0\longrightarrow\underline a\longrightarrow\underline A
\longrightarrow0$, 还是使用 $\underline A/\underline b$ 的分解 $0\longrightarrow\underline b\longrightarrow\underline A
\longrightarrow0$. 在这里它表现为图
\begin{equation}\tag{\SGARefTarget{sga7.sga7i.eviii.n2-1-14-3}{2.1.14.3}}
\begin{tikzcd}[column sep=huge,row sep=huge]
 {}_{\underline a}P\otimes{}_{\underline b}Q
   \arrow[r]
   \arrow[d]
   \arrow[rd,phantom,"-1" description]&
 {}_{\underline b}Q\otimes{}_{\underline a}P
   \arrow[d]\\
 \uHom\!\left(
   \dfrac{\underline a\cap\underline b}
         {\underline a\mathbin{\cdot}\underline b},
   \uTor_1(P,Q)
 \right)
   \arrow[r]&
 \uHom\!\left(
   \dfrac{\underline a\cap\underline b}
         {\underline a\mathbin{\cdot}\underline b},
   \uTor_1(Q,P)
 \right).
\end{tikzcd}
\end{equation}
的反交换性. 水平箭头为对称箭头. 当 $\underline a=\underline b$ 时, 这种反交换性也可以表述为: $\uTor_1(\underline A/\underline a,\underline A/\underline a)$ 上的对称自同构是 $-\id$.

在 $\underline a=\underline b=n\underline A$ 的特殊情形下, 其中 $n$ 是 $\underline A$ 的截面, 乘以 $n$ 诱导一个满态射
\[
 \underline A/\underline a
 \longrightarrow
 \frac{\underline a\cap\underline a}
      {\underline a\mathbin{\cdot}\underline a}
 =\underline a/\underline a^2,
\]
从而得到关于 $\underline A$-模 $T$ 具有函子性的单同态
\[
 \uHom(\underline a/\underline a^2,T)
 \longrightarrow{}_{\underline a}T.
\]
因此, 此时同态 (\SGARefLink{sga7.sga7i.eviii.n2-1-14-1}{2.1.14.1}) 可解释为同态
\begin{equation}\tag{\SGARefTarget{sga7.sga7i.eviii.n2-1-14-4}{2.1.14.4}}
 {}_nP\otimes{}_nQ
 \longrightarrow{}_n\uTor_1(P,Q),
\end{equation}
它在形式上更接近 (\SGARefLink{sga7.sga7i.eviii.n2-1-1}{2.1.1}).

继续假设 $\underline a=\underline b=n\underline A$, 并且 $n$ 是 $\underline A$ 的正则截面. 谱序列 (\SGARefLink{sga7.sga7i.eviii.n2-1-4}{2.1.4}) 在当前情形仍有意义, 并给出以下低次项正合列, 推广了 (\SGARefLink{sga7.sga7i.eviii.n2-1-5}{2.1.5}); 此时 $Q$ 是被 $\underline a$ 零化的 $\underline A$-模:
\begin{equation}\tag{\SGARefTarget{sga7.sga7i.eviii.n2-1-14-5}{2.1.14.5}}
\begin{split}
 \uTor^{\underline A}_2(P,Q)
 &\longrightarrow
 \uTor^{\underline A/n\underline A}_2(P_n,Q)
 \longrightarrow{}_nP\otimes Q
 \xrightarrow{\ \alpha\ }
 \uTor^{\underline A}_1(P,Q)\\
 &\longrightarrow
 \uTor^{\underline A/n\underline A}_1(P_n,Q)
 \longrightarrow0.
\end{split}
\end{equation}

现在回到任意 $\underline a$ 和 $\underline b$ 的情形, 设 $p$ 和 $q$ 是 $\underline A$ 的两个截面. 便得到下列交换图; (\SGARefLink{sga7.sga7i.eviii.n2-1-11-1}{2.1.11.1}) 本质上就是取 $q=1$ 所得到的特殊情形:
\begin{equation}\tag{\SGARefTarget{sga7.sga7i.eviii.n2-1-14-6}{2.1.14.6}}
{
\begin{tikzcd}[row sep=large,column sep=tiny]
 &{}_{p\underline a}P\otimes{}_{q\underline b}Q
   \arrow[dl,"p\otimes q"']
   \arrow[dr,"i\otimes j"]&\\
 {}_{\underline a}P\otimes{}_{\underline b}Q
   \arrow[d,"\alpha"']&&
 {}_{pq\underline a}P\otimes{}_{pq\underline b}Q
   \arrow[d,"\alpha"]\\
 \uHom\!\left(
   \dfrac{\underline a\cap\underline b}
         {\underline a\mathbin{\cdot}\underline b},
   \uTor_1(P,Q)
 \right)
   \arrow[dr,"\text{包含}"']&&
 \uHom\!\left(
   \dfrac{pq\underline a\cap pq\underline b}
         {pq\underline a\mathbin{\cdot}pq\underline b},
   \uTor_1(P,Q)
 \right)
   \arrow[dl,"(pq)^*"]\\
 &\uHom\!\left(
   \dfrac{\underline a\cap\underline b}
         {pq\underline a\mathbin{\cdot}\underline b},
   \uTor_1(P,Q)
 \right)&
\end{tikzcd}}
\end{equation}
这里
\[
 i:{}_{p\underline a}P\longrightarrow{}_{pq\underline a}P,
 \qquad
 j:{}_{q\underline b}Q\longrightarrow{}_{pq\underline b}Q
\]
是典范包含映射, 而 $(pq)^*$ 表示从乘以 $pq$ 所诱导的满同态
\[
 \frac{\underline a\cap\underline b}
      {pq\underline a\mathbin{\cdot}\underline b}
 \longrightarrow
 \frac{pq\underline a\cap pq\underline b}
      {pq\underline a\mathbin{\cdot}pq\underline b}
\]
通过取 $\uHom$ 得到的单同态. 当 $\underline a=\underline b=n\underline A$ 时, 可以把此图写成较简单的形式; 当 $q=1$ 时这一形式见 (\SGARefLink{sga7.sga7i.eviii.n2-1-11-1}{2.1.11.1}).

% END INCLUDED SGA7I_English_current_buildable_source_20260823/components/97_expose_VIII_section_2_1_13_and_remark_2_1_14.tex

% BEGIN INCLUDED SGA7I_English_current_buildable_source_20260823/components/98_expose_VIII_section_2_2.tex
% R3 explicit revisions R3-VIII-03; see qa/FIDELITY-R3-CHANGES.json.
% SOURCE AUTHORITY:
% expose_VIII_body.tex lines 676--834; scan indices 248--252; source folios 237--241.
% Direct scan comparison: physical PDF pages 249--253 at 1100 dpi.
% Transparent source dispositions:
% - in 2.2.7, the restriction is to the image of alpha^n_{P,Q} in
%   Tor_1(P,Q), not the authority's impossible phi-map to Tor_1(P,Q;G);
% - the cited pairings are (2.2.5) and (2.2.6), not (2.2.4) and (2.2.6);
% - typewritten letter-o complex subscripts are normalized to zero;
% - the missing closing parenthesis in the lower-left node of (2.2.11.2)
%   is restored.
% Hard stop before 2.3, source line 836.

\SGARefTarget{sga7.sga7i.eviii.n2-2}{2.2}. 现在除 $P$ 和 $Q$ 外, 再考虑一个交换群 $G$. 在 (\SGARefLink{sga7.sga7i.eviii.n1-1-3}{1.1.3}) 中, 我们定义了对 $P,Q$ 和 $G$ 具有函子性的典范同态
\begin{equation}\tag{\SGARefTarget{sga7.sga7i.eviii.n2-2-1}{2.2.1}}
 \mathrm{Biext}^1(P,Q;G)
 \longrightarrow
 \operatorname{Hom}\bigl(\operatorname{Tor}_1(P,Q),G\bigr).
\end{equation}
另一方面, 若 $n$ 为正整数 (分别地, 若 $\ell$ 为素数), 则 (\SGARefLink{sga7.sga7i.eviii.n2-1-1}{2.1.1}) 的同态 $\alpha^n_{P,Q}$ (分别地, (\SGARefLink{sga7.sga7i.eviii.n2-1-13-2}{2.1.13.2}) 的同态 $\alpha^{(\ell)}_{P,Q}$) 给出同态
\[
 \operatorname{Hom}\bigl(\operatorname{Tor}_1(P,Q),G\bigr)
 \longrightarrow
 \operatorname{Hom}({}_nP\otimes{}_nQ,G)
\]
(分别地, 同态
\[
 \operatorname{Hom}\bigl(\operatorname{Tor}_1(P,Q),G\bigr)
 \longrightarrow
 \operatorname{Hom}\bigl(
   T_\ell(P)\otimes T_\ell(Q),T_\ell(G)
 \bigr)).
\]
这些同态也对 $P,Q,G$ 具有函子性. 将它们与 (\SGARefLink{sga7.sga7i.eviii.n2-2-1}{2.2.1}) 复合, 得到对 $P,Q,G$ 具有函子性的典范同态
\begin{equation}\tag{\SGARefTarget{sga7.sga7i.eviii.n2-2-2}{2.2.2}}
 \varphi^n=\varphi^n_{P,Q}:
 \mathrm{Biext}^1(P,Q;G)
 \longrightarrow
 \operatorname{Hom}({}_nP\otimes{}_nQ,G),
\end{equation}
\begin{equation}\tag{\SGARefTarget{sga7.sga7i.eviii.n2-2-3}{2.2.3}}
 \varphi^{(\ell)}=\varphi^{(\ell)}_{P,Q}:
 \mathrm{Biext}^1(P,Q;G)
 \longrightarrow
 \operatorname{Hom}\bigl(
   T_\ell(P)\otimes T_\ell(Q),T_\ell(G)
 \bigr).
\end{equation}
对于
\begin{equation}\tag{\SGARefTarget{sga7.sga7i.eviii.n2-2-4}{2.2.4}}
 \xi\in\mathrm{Biext}^1(P,Q;G),
\end{equation}
用 $\varphi^n_\xi$ (分别地, $\varphi^{(\ell)}_\xi$) 表示其在 (\SGARefLink{sga7.sga7i.eviii.n2-2-2}{2.2.2}) (分别地, (\SGARefLink{sga7.sga7i.eviii.n2-2-3}{2.2.3})) 下的像:
\begin{equation}\tag{\SGARefTarget{sga7.sga7i.eviii.n2-2-5}{2.2.5}}
 \varphi^n_\xi:{}_nP\otimes{}_nQ
 \longrightarrow G
 \qquad\text{(或等价地, }{}_nG\text{)},
\end{equation}
\begin{equation}\tag{\SGARefTarget{sga7.sga7i.eviii.n2-2-6}{2.2.6}}
 \varphi^{(\ell)}_\xi:
 T_\ell(P)\otimes T_\ell(Q)
 \longrightarrow T_\ell(G).
\end{equation}

\SGARefTarget{sga7.sga7i.eviii.n2-2-7}{2.2.7}. 显然, 知道 $\varphi^{(\ell)}_\xi$ 等价于对所有形如 $n=\ell^\nu$ ($\nu\geq0$) 的指数知道 $\varphi^n_\xi$. 对于固定的 $n$, 知道 $\varphi^n_\xi$ 等价于在
\[
 \alpha^n_{P,Q}:{}_nP\otimes{}_nQ\longrightarrow\uTor_1(P,Q),
\]
的像上知道同态
\begin{equation}\tag{\SGARefTarget{sga7.sga7i.eviii.n2-2-8}{2.2.8}}
 \uTor_1(P,Q)\longrightarrow G
\end{equation}
它是 $\xi$ 在 (\SGARefLink{sga7.sga7i.eviii.n2-2-1}{2.2.1}) 下的像.

注意, 若 $P=nP$ 且 $Q=nQ$, 则上述像为 $\,{}_n\uTor_1(P,Q)$. 考虑正合列 (\SGARefLink{sga7.sga7i.eviii.n1-3-5}{1.3.5}) 及关于 $Q$ 的类似正合列, 即容易得出这一点. 此外, 若 $P$ 和 $Q$ 都是可除的, 即对每个正整数 $n$ 都有 $nP=P$ 和 $nQ=Q$, 则容易验证 $\uTor_1(P,Q)$ 是挠群, 也就是说
\[
 \uTor_1(P,Q)=\varinjlim_n{}_{n}\uTor_1(P,Q).
\]
因此, 在这种情形下, 对所有 $n$ 知道 $\varphi^n_\xi$ 等价于知道 (\SGARefLink{sga7.sga7i.eviii.n2-2-8}{2.2.8}), 而对给定素数 $\ell$ 知道 $\varphi^{(\ell)}_\xi$ 则等价于知道 (\SGARefLink{sga7.sga7i.eviii.n2-2-8}{2.2.8}) 在 $\uTor_1(P,Q)$ 的 $\ell$-准素分量上的限制.

下面具体说明双扩张所对应的配对 (\SGARefLink{sga7.sga7i.eviii.n2-2-5}{2.2.5}) 和 (\SGARefLink{sga7.sga7i.eviii.n2-2-6}{2.2.6}) 的对称性质. 首先注意有交换图
\begin{equation}\tag{\SGARefTarget{sga7.sga7i.eviii.n2-2-9}{2.2.9}}
\begin{tikzcd}[column sep=large,row sep=huge]
 \mathrm{Biext}^1(P,Q;G)
   \arrow[r,"\sim"]
   \arrow[d]&
 \mathrm{Biext}^1(Q,P;G)
   \arrow[d]\\
 \operatorname{Hom}\bigl(\uTor_1(P,Q),G\bigr)
   \arrow[r,"\sim"]&
 \operatorname{Hom}\bigl(\uTor_1(Q,P),G\bigr),
\end{tikzcd}
\end{equation}
其中水平箭头为对称箭头 (第一个定义于 VII \SGARefLink{sga7.sga7i.evii.n2-7}{2.7}), 竖直箭头为 (\SGARefLink{sga7.sga7i.eviii.n2-2-1}{2.2.1}). 它的交换性来自下列较完整的图中两个组成方块的交换性:
\begin{equation}\tag{\SGARefTarget{sga7.sga7i.eviii.n2-2-10}{2.2.10}}
\begin{tikzcd}[column sep=large,row sep=large]
 \mathrm{Biext}^1(P,Q;G)
   \arrow[r,"\sim"]
   \arrow[d]&
 \mathrm{Biext}^1(Q,P;G)
   \arrow[d]\\
 \operatorname{Ext}^1\!\left(
   P\mathbin{\overset{\mathbb L}{\otimes}}Q,G
 \right)
   \arrow[r,"\sim"]
   \arrow[d]&
 \operatorname{Ext}^1\!\left(
   Q\mathbin{\overset{\mathbb L}{\otimes}}P,G
 \right)
   \arrow[d]\\
 \operatorname{Hom}\bigl(\uTor_1(P,Q),G\bigr)
   \arrow[r,"\sim"]&
 \operatorname{Hom}\bigl(\uTor_1(Q,P),G\bigr).
\end{tikzcd}
\end{equation}
每列第一个竖直箭头是 VII \SGARefLink{sga7.sga7i.evii.n3-6-5}{3.6.5} 的典范同构, 第二个则是显然的同态. (\SGARefLink{sga7.sga7i.eviii.n2-2-10}{2.2.10}) 的第二个方块显然交换: 根据定义, 对称同构
\[
 P\mathbin{\overset{\mathbb L}{\otimes}}Q
 \xrightarrow{\ \sim\ }
 Q\mathbin{\overset{\mathbb L}{\otimes}}P
\]
诱导对称同构 $\uTor_1(P,Q)\simeq\uTor_1(Q,P)$, 而对于可变复形 $L_\bullet$, 典范同态
\[
 \operatorname{Ext}^1(L_\bullet,G)
 \longrightarrow
 \operatorname{Hom}\bigl(H_1(L_\bullet),G\bigr)
\]
对 $L_\bullet$ 具有函子性.

至于 (\SGARefLink{sga7.sga7i.eviii.n2-2-10}{2.2.10}) 的第一个方块的交换性, 只须同时具体写出对称同构及 VII \SGARefLink{sga7.sga7i.evii.n3-6-5}{3.6.5} 所定义的竖直同构即可得出. 验证的细节留给读者. 这里用到的事实是: 对称同构
\[
 L_\bullet(P)\otimes L_\bullet(Q)
 \simeq
 L_\bullet(Q)\otimes L_\bullet(P)
\]
在次数 $1$ 上诱导同构
\[
 L_0(P)\otimes L_1(Q)+L_1(P)\otimes L_0(Q)
 \simeq
 L_0(Q)\otimes L_1(P)+L_1(Q)\otimes L_0(P),
\]
它在相应的直和项上由底层无分次模的通常对称同构给出 (因而不引入符号).

根据 (\SGARefLink{sga7.sga7i.eviii.n2-2-9}{2.2.9}) 的交换性, (\SGARefLink{sga7.sga7i.eviii.n2-1-10-1}{2.1.10.1}) 的反交换性, 用 $\varphi^n$ 或 $\varphi^{(\ell)}$ 表示, 就变成如下关系.

\medskip
\noindent\underline{命题 \SGARefTarget{sga7.sga7i.eviii.n2-2-11}{2.2.11}.}
\underline{下列图反交换:}
\begin{equation}\tag{\SGARefTarget{sga7.sga7i.eviii.n2-2-11-1}{2.2.11.1}}
\begin{tikzcd}[row sep=large,column sep=huge]
 \mathrm{Biext}^1(P,Q;G)
   \arrow[rr,"\sim"]
   \arrow[dd,"{\varphi^n_{P,Q}}"']&&
 \mathrm{Biext}^1(Q,P;G)
   \arrow[dd,"{\varphi^n_{Q,P}}"]\\
 &- &\\
 \operatorname{Hom}({}_nP\otimes{}_nQ,G)
   \arrow[rr,"\sim"]&&
 \operatorname{Hom}({}_nQ\otimes{}_nP,G),
\end{tikzcd}
\end{equation}
\begin{equation}\tag{\SGARefTarget{sga7.sga7i.eviii.n2-2-11-2}{2.2.11.2}}
\begin{tikzcd}[row sep=large,column sep=huge]
 \mathrm{Biext}^1(P,Q;G)
   \arrow[rr,"\sim"]
   \arrow[dd,"{\varphi^{(\ell)}_{P,Q}}"']&&
 \mathrm{Biext}^1(Q,P;G)
   \arrow[dd,"{\varphi^{(\ell)}_{Q,P}}"]\\
 &- &\\
 \operatorname{Hom}\bigl(
   T_\ell(P)\otimes T_\ell(Q),T_\ell(G)
 \bigr)
   \arrow[rr,"\sim"]&&
 \operatorname{Hom}\bigl(
   T_\ell(Q)\otimes T_\ell(P),T_\ell(G)
 \bigr),
\end{tikzcd}
\end{equation}
\underline{其中水平箭头为对称箭头.}

记下以下重要的推论.

\medskip
\noindent\underline{推论 \SGARefTarget{sga7.sga7i.eviii.n2-2-12}{2.2.12}.}
\underline{假设 $P=Q$, 并假设 $\xi$ 是 $(P,P)$ 以 $G$ 为核的对称 (分别地, 反对称) 双扩张的类; 即假设 ${}^s\xi=\xi$ (分别地, ${}^s\xi=-\xi$), 其中 ${}^s\xi$ 表示 $\xi$ 在对称映射下的像. 则 ${}_nP\times{}_nP$ 上的双线性映射 $\varphi^n_\xi$ 以及 $T_\ell(P)\times T_\ell(P)$ 上的双线性映射 $\varphi^{(\ell)}_\xi$ 都是反对称的 (分别地, 对称的).}

\medskip
\noindent\underline{注 \SGARefTarget{sga7.sga7i.eviii.n2-2-13}{2.2.13}.}
假设 $\xi$ 对称. 若 $n$ (分别地, $\ell$) 为奇数, 则 \SGARefLink{sga7.sga7i.eviii.n2-2-12}{2.2.12} 表明 $\varphi^n_\xi$ (分别地, $\varphi^{(\ell)}_\xi$) 实际上是交错形式. 若不要求 $n$ (分别地, $\ell$) 为奇数, 这个结论一般不成立; 即使当 $P$ 为通常的群 $\mathbb Z/2\mathbb Z$ 且 $G=\mathbb Z/2\mathbb Z$ 时, 就已可看出这一点.

不过, 假设 $2P=P$. 特别地, 如通常那样在 fppf 或 fpqc 拓扑中工作时, 若 $P$ 为阿贝尔概形, 就有这种情形. 此时前述形式仍是交错的. 只须证明 $\varphi^{2^\nu}_\xi(x,x)=0$. 因为 $2P=P$, 给定元素 $x\in{}_{2^\nu}P(S)$ 可以写成 $x=2y$, 其中 $y\in{}_{2^{\nu+1}}P(S)$. 从而
\[
 \varphi^{2^\nu}_\xi(x,x)
 =2\varphi^{2^{\nu+1}}_\xi(y,y)=0,
\]
因为 $\varphi^{2^{\nu+1}}_\xi$ 是反对称的.

% END INCLUDED SGA7I_English_current_buildable_source_20260823/components/98_expose_VIII_section_2_2.tex

% BEGIN INCLUDED SGA7I_English_current_buildable_source_20260823/components/99_expose_VIII_section_2_3_opening_through_proposition_2_3_11.tex
% SOURCE AUTHORITY:
% expose_VIII_body.tex lines 836--923; scan indices 252--254; source folios 241--243.
% Direct scan comparison: physical PDF pages 253--255 at 1100 dpi.
% Transparent source disposition:
% - the right-hand vertex of (2.3.10) is restored as
%   Hom({}_nP,{}_n Ext^1(nQ,G)); the printed/transcribed "nQnG" is a
%   comma-key slip and is incompatible with (2.3.5), (2.3.9), and the prose.
% The authority's numbering jump from (2.3.2) to (2.3.4) is preserved.
% Hard stop before the proof of 2.3.11, source line 924.

\SGARefTarget{sga7.sga7i.eviii.n2-3}{2.3}. \underline{与“Cartier 对偶”的关系.}
\SGARefLink{sga7.sga7i.eviii.n2-2}{2.2} 利用同构定理 VII \SGARefLink{sga7.sga7i.evii.n3-6-5}{3.6.5} 及同态 (\SGARefLink{sga7.sga7i.eviii.n2-2-1}{2.2.1}) 来描述双扩张所定义的配对; 这显然不是早先用于阿贝尔概形的定义, 因为那时甚至还没有双扩张这一概念. 早先的定义通过 Cartier 对偶来给出, 下面说明它与当前定义之间的关系.

再考虑由正合列 (\SGARefLink{sga7.sga7i.eviii.n1-4}{1.4}) 得到的同态
\begin{equation}\tag{\SGARefTarget{sga7.sga7i.eviii.n2-3-1}{2.3.1}}
 \mathrm{Biext}^1(P,Q;G)
 \longrightarrow
 \operatorname{Hom}\bigl(P,\uExt^1(Q,G)\bigr)
\end{equation}
当
\begin{equation}\tag{\SGARefTarget{sga7.sga7i.eviii.n2-3-2}{2.3.2}}
 \uHom(Q,G)=0.
\end{equation}
时, 它实际上是同构. 特别地, 当 $Q$ 是概形 $S$ 上的阿贝尔概形, 而 $G$ 是乘法群 $\underline G_m$ 时, 这个条件成立 (VII \SGARefLink{sga7.sga7i.evii.n1-3-8}{1.3.8}).

此外还有显然的“限制到 $nQ$”同态
\[
 \uExt^1(Q,G)\longrightarrow\uExt^1(nQ,G),
\]
从而有复合同态
\begin{equation}\tag{\SGARefTarget{sga7.sga7i.eviii.n2-3-4}{2.3.4}}
\begin{split}
 \operatorname{Hom}\bigl(P,\uExt^1(Q,G)\bigr)
 &\longrightarrow
 \operatorname{Hom}\bigl(P,\uExt^1(nQ,G)\bigr)\\
 &\longrightarrow
 \operatorname{Hom}\bigl({}_nP,{}_n\uExt^1(nQ,G)\bigr),
\end{split}
\end{equation}
其中第二个箭头是限制到 $\,{}_nP$. 与 (\SGARefLink{sga7.sga7i.eviii.n2-3-1}{2.3.1}) 复合, 得到同态
\begin{equation}\tag{\SGARefTarget{sga7.sga7i.eviii.n2-3-5}{2.3.5}}
 \mathrm{Biext}^1(P,Q;G)
 \longrightarrow
 \operatorname{Hom}\bigl({}_nP,{}_n\uExt^1(nQ,G)\bigr).
\end{equation}

另一方面, 正合列
\begin{equation}\tag{\SGARefTarget{sga7.sga7i.eviii.n2-3-6}{2.3.6}}
 0\longrightarrow{}_nQ\longrightarrow Q
 \xrightarrow{\ n\ }nQ\longrightarrow0
\end{equation}
给出 $\uExt^i(-,G)$ 的长正合列, 从而给出短正合列
\begin{equation}\tag{\SGARefTarget{sga7.sga7i.eviii.n2-3-7}{2.3.7}}
 0\longrightarrow\uHom(Q,G)_n
 \longrightarrow\uHom({}_nQ,G)
 \xrightarrow{\ \partial\ }
 {}_n\uExt^1(nQ,G)
 \longrightarrow0,
\end{equation}
其中照例用 $E_n$ 表示 $E/nE$. 从这个正合列中取出典范满同态
\begin{equation}\tag{\SGARefTarget{sga7.sga7i.eviii.n2-3-8}{2.3.8}}
 \uHom({}_nQ,G)
 \xrightarrow{\hspace{3em}\partial\hspace{3em}}
 {}_n\uExt^1(nQ,G),
\end{equation}
若 $\uHom(Q,G)_n=0$, 它就是同构; 特别地, 若 (\SGARefLink{sga7.sga7i.eviii.n2-3-2}{2.3.2}) 成立, 亦如此. 由 (\SGARefLink{sga7.sga7i.eviii.n2-3-8}{2.3.8}) 推出同态
\begin{equation}\tag{\SGARefTarget{sga7.sga7i.eviii.n2-3-9}{2.3.9}}
\begin{split}
 \operatorname{Hom}\bigl({}_nP,\uHom({}_nQ,G)\bigr)
 &\simeq
 \operatorname{Hom}({}_nP\otimes{}_nQ,G)\\
 &\longrightarrow
 \operatorname{Hom}\bigl({}_nP,{}_n\uExt^1(nQ,G)\bigr),
\end{split}
\end{equation}
当 (\SGARefLink{sga7.sga7i.eviii.n2-3-2}{2.3.2}) 成立时, 这也是同构.

所考虑的一切同态都对 $P,Q,G$ 具有函子性. 同态 (\SGARefLink{sga7.sga7i.eviii.n2-3-5}{2.3.5})、(\SGARefLink{sga7.sga7i.eviii.n2-3-9}{2.3.9}) 和 (\SGARefLink{sga7.sga7i.eviii.n2-2-2}{2.2.2}) 构成图
\begin{equation}\tag{\SGARefTarget{sga7.sga7i.eviii.n2-3-10}{2.3.10}}
\begin{tikzcd}[column sep=large,row sep=large]
 &\mathrm{Biext}^1(P,Q;G)
   \arrow[ddl,"\varphi^n"']
   \arrow[ddr,"(2.3.5)"]&\\
 &-&\\
 \operatorname{Hom}({}_nP\otimes{}_nQ,G)
   \arrow[rr,"(2.3.9)"]&&
 \operatorname{Hom}\bigl({}_nP,{}_n\uExt^1(nQ,G)\bigr),
\end{tikzcd}
\end{equation}
其中若 (\SGARefLink{sga7.sga7i.eviii.n2-3-2}{2.3.2}) 成立, 水平箭头 (\SGARefLink{sga7.sga7i.eviii.n2-3-9}{2.3.9}) 就是同构.

\medskip
\noindent\underline{命题 \SGARefTarget{sga7.sga7i.eviii.n2-3-11}{2.3.11}.}
\underline{图 (\SGARefLink{sga7.sga7i.eviii.n2-3-10}{2.3.10}) 反交换.}

当 $\uHom(Q,G)=0$ 时, 这一命题就用同态 (\SGARefLink{sga7.sga7i.eviii.n2-3-5}{2.3.5}) 和 (\SGARefLink{sga7.sga7i.eviii.n2-3-9}{2.3.9}) 重新给出了对 (\SGARefLink{sga7.sga7i.eviii.n2-2-2}{2.2.2}) 的同态 $\varphi^n$ 的刻画. 这正是早先在阿贝尔概形理论中所用的刻画; 在那里还有 $nQ=Q$.

% END INCLUDED SGA7I_English_current_buildable_source_20260823/components/99_expose_VIII_section_2_3_opening_through_proposition_2_3_11.tex

% BEGIN INCLUDED SGA7I_English_current_buildable_source_20260823/components/100_expose_VIII_proof_2_3_11_parts_a_and_b.tex
% SOURCE AUTHORITY:
% expose_VIII_body.tex lines 924--1021; scan indices 254--257; source folios 243--246.
% Direct scan comparison: physical PDF pages 255--258 at 1100 dpi.
% Transparent source dispositions:
% - the proof heading's printed 2.5.11 is restored as 2.3.11;
% - typewritten letter-o indices and cochain degrees are normalized to zero;
% - the arrow-key slip inside the defining formula (2.3.14) is restored as
%   the minus sign already printed in (2.3.13) and required by the cocycle.
% Hard stop before part c), source line 1022.

\medskip
\noindent\underline{\SGARefLink{sga7.sga7i.eviii.n2-3-11}{2.3.11} 的证明.}\footnote{更简洁的证明见下面的 \SGARefLink{sga7.sga7i.eviii.n2-3-19}{2.3.19}.}
设 $\xi$ 为 $\mathrm{Biext}^1(P,Q;G)$ 的一个元素,
\[
 \varphi_\xi:{}_nP\otimes{}_nQ\longrightarrow G
\]
为所对应的配对, 并设
\[
 u_\xi:{}_nP\longrightarrow{}_n\uExt^1(nQ,G)
\]
为 $\xi$ 在 (\SGARefLink{sga7.sga7i.eviii.n2-3-5}{2.3.5}) 下的像. 须证明 $-u_\xi$ 是 $\varphi_\xi$ 在箭头 (\SGARefLink{sga7.sga7i.eviii.n2-3-9}{2.3.9}) 下的像, 即它是复合映射
\begin{equation}\tag{*}
 {}_nP
 \xrightarrow{\ \varphi_\xi\ }
 \uHom({}_nQ,G)
 \xrightarrow{\ (\SGARefLink{sga7.sga7i.eviii.n2-3-8}{2.3.8})\ }
 {}_n\uExt^1(nQ,G).
\end{equation}
为此, 可以取 $\,{}_nP$ 的一个截面 $x$, 证明它在 $-u_\xi$ 下的像等于它在复合映射 (*) 下的像. 如果愿意, 这样就可以归结为 $P=(\mathbb Z/n\mathbb Z)_T$ 的情形, 不过我们不需要作这一归结.

\medskip
\noindent a) 分别取 $P$ 和 $Q$ 的平坦分解 $L_\bullet$ 和 $M_\bullet$,
\[
 0\longrightarrow L_1\longrightarrow L_0\longrightarrow0,
 \qquad
 0\longrightarrow M_1\longrightarrow M_0\longrightarrow0,
\]
例如, 将 $P$ 写成平坦 $\mathbb Z_T$-模 $L_0$ 的商, 并取 $L_1$ 为 $L_0\longrightarrow P$ 的核, 就得到 $P$ 的平坦分解; 这个核无挠, 所以平坦. 得到复形 $L_\bullet\otimes M_\bullet$
\[
 0\longrightarrow L_1\otimes M_1
 \xrightarrow{d_2}
 L_1\otimes M_0+L_0\otimes M_1
 \xrightarrow{d_1}
 L_0\otimes M_0
 \longrightarrow0,
\]
其层 $\underline H_1$ 将用作 $\uTor_1(P,Q)$.

设 $L'_0$ (分别地, $M'_0$) 为 $L_0$ (分别地, $M_0$) 中 $\,{}_nP$ (分别地, $\,{}_nQ$) 的逆像子层. 则
\[
 0\longrightarrow L_1\longrightarrow L'_0\longrightarrow0
 \qquad
 \left(\text{分别地, }
 0\longrightarrow M_1\longrightarrow M'_0\longrightarrow0\right)
\]
是 $\,{}_nP$ (分别地, $\,{}_nQ$) 的平坦分解. 同态
\[
 {}_nP\otimes{}_nQ
 \longrightarrow
 \uTor_1(P,Q)=\underline H_1(L_\bullet\otimes M_\bullet)
\]
于是可描述为由同态
\[
 L'_0\otimes M'_0
 \longrightarrow
 \underline H_1(L_\bullet\otimes M_\bullet)
\]
通过取商得到, 后一同态由
\[
 x_0\otimes y_0
 \longmapsto
 \text{$1$-闭链的类 }
 [nx_0]\otimes y_0-x_0\otimes[ny_0].
\]
给出. 这里方括号表示把所讨论的元素视为 $L_0$ 的子层 $L_1$ (分别地, $M_0$ 的子层 $M_1$) 的截面. 这个描述由 \SGARefLink{sga7.sga7i.eviii.n2-1-10}{2.1.10} 的证明中的计算得出.

\medskip
\noindent b) 取 $G$ 的内射分解
\[
 C^\bullet=\bigl(0\longrightarrow C^0
 \xrightarrow{d^0}C^1\longrightarrow\cdots\bigr)
\]
并作如下等同
\[
 \mathrm{Biext}^1(P,Q;G)
 \simeq
 \operatorname{Ext}^1\!\left(
   P\mathbin{\overset{\mathbb L}{\otimes}}Q,G
 \right)
 =H^1\!\left(
   \operatorname{Hom}^{\boldsymbol\cdot}
   (L_\bullet\otimes M_\bullet,C^{\boldsymbol\cdot})
 \right).
\]
于是元素 $\xi$ 可由一对同态 $f_1,f_0$ 表示, 给出反交换图
\begin{equation}\tag{\SGARefTarget{sga7.sga7i.eviii.n2-3-12}{2.3.12}}
\begin{tikzcd}[column sep=huge,row sep=huge]
 L_1\otimes M_0+L_0\otimes M_1
   \arrow[r,"d_1"]
   \arrow[d,"f_1"']
   \arrow[rd,phantom,"-" description]&
 L_0\otimes M_0
   \arrow[d,"f_0"']\\
 C^0\arrow[r,"d^0"]&C^1.
\end{tikzcd}
\end{equation}
按定义, 配对
\[
 \varphi_\xi:{}_nP\otimes{}_nQ\longrightarrow G
\]
由同态
\[
 \varphi:L'_0\otimes M'_0\longrightarrow G
\]
取商得到, 而后者由
\begin{equation}\tag{\SGARefTarget{sga7.sga7i.eviii.n2-3-13}{2.3.13}}
 \varphi(x_0\otimes y_0)
 =f_1\bigl([nx_0]\otimes y_0-x_0\otimes[ny_0]\bigr)
 \in G=\ker(C^0\longrightarrow C^1).
\end{equation}
给出. 因为断言 \SGARefLink{sga7.sga7i.eviii.n2-3-11}{2.3.11} 在 $T$ 上是局部的, 可以假设 $\,{}_nP$ 的给定截面 $x$ 提升成 $L'_0$ 的截面 $x_0$. 把 $y_0$ 视为变元, 上式就定义同态
\begin{equation}\tag{\SGARefTarget{sga7.sga7i.eviii.n2-3-14}{2.3.14}}
 \varphi':M'_0\longrightarrow G,
 \qquad
 y_0\longmapsto
 f_1\bigl([nx_0]\otimes y_0-x_0\otimes[ny_0]\bigr),
\end{equation}
取商后给出
\[
 {}_nQ\longrightarrow G,
 \qquad
 y\longmapsto\varphi_\xi(x,y),
\]
即对应于 $\varphi_\xi$ 和 $x$ 的同态.

% END INCLUDED SGA7I_English_current_buildable_source_20260823/components/100_expose_VIII_proof_2_3_11_parts_a_and_b.tex

% BEGIN INCLUDED SGA7I_English_current_buildable_source_20260823/components/101_expose_VIII_proof_2_3_11_parts_c_through_e_and_remark_2_3_16.tex
% SOURCE AUTHORITY:
% expose_VIII_body.tex lines 1022--1221; scan indices 257--260;
% source folios 246--249.
% Direct scan comparison: physical PDF pages 258--261 in overlapping
% 1100-dpi bands.  The ambiguous right-hand component in the first diagram
% of part d was additionally examined at 5000-dpi equivalent magnification.
% Transparent source dispositions:
% - typewritten letter-o indices and cochain degrees are normalized to zero;
% - the source prints f_1 on the right vertical arrow of the degree-one
%   diagram in part d; its codomain is C^1, and the immediately following
%   formula reads g_0=n f_0, so the English diagram restores the type-correct
%   f_0;
% - the uppercase X in Remark 2.3.16 is restored to the section x used in
%   the preceding and following clauses.
% Hard stop before 2.3.17, source line 1222.

\medskip
\noindent c) 须取这个同态在连接算子 $\partial$ 下的像; 此算子来自 (\SGARefLink{sga7.sga7i.eviii.n2-3-7}{2.3.7}), 对应于正合列 (\SGARefLink{sga7.sga7i.eviii.n2-3-6}{2.3.6}). 现在 $nQ\simeq Q/{}_nQ$ 显然也同构于 $M_0/M'_0$, 因而有正合列的同态
\[
\begin{tikzcd}[column sep=large,row sep=large]
 0\arrow[r]&M'_0\arrow[r]\arrow[d]&M_0\arrow[r]\arrow[d]
   &M_0/M'_0\arrow[r]\arrow[d,"\wr"']&0\\
 0\arrow[r]&{}_nQ\arrow[r]&Q\arrow[r]&nQ\arrow[r]&0&{,}
\end{tikzcd}
\]
其中最后一个竖直箭头是同构. 这表明所求 $\uExt^1(nQ,G)$ 的截面也可以计算为 (\SGARefLink{sga7.sga7i.eviii.n2-3-14}{2.3.14}) 中 $\varphi'$ 在连接算子
\[
 \uHom(M'_0,G)\xrightarrow{\ \partial\ }
 \uExt^1(M_0/M'_0,G),
\]
下的像; 等价地, 它是第二项的一个截面, 由如下整体 $\mathrm{Ext}^1$ 的元素定义: 取 $\varphi'$ 在连接算子
\[
 \mathrm{Hom}(M'_0,G)\longrightarrow
 \mathrm{Ext}^1(M_0/M'_0,G).
\]
下的像. 连接算子的通常构造给出如下做法. 将复合映射 (\SGARefLink{sga7.sga7i.eviii.n2-3-13}{2.3.13})
\[
 M'_0\xrightarrow{\ \varphi'\ }G\longrightarrow C^0
\]
延拓为同态
\[
 \psi:M_0\longrightarrow C^0,
\]
这是可能的, 因为 $C^0$ 内射; 再考虑唯一同态
\begin{equation}\tag{\SGARefTarget{sga7.sga7i.eviii.n2-3-15}{2.3.15}}
 \lambda:M_0/M'_0\longrightarrow C^1
\end{equation}
它使图
\[
\begin{tikzcd}[column sep=large,row sep=large]
 0\arrow[r]&M'_0\arrow[r]\arrow[d,"\varphi'"']
   &M_0\arrow[r]\arrow[d,"\psi"']
   &M_0/M'_0\arrow[r]\arrow[d,"\lambda"']&0\\
 0\arrow[r]&G\arrow[r]&C^0\arrow[r]&C^1&&{.}
\end{tikzcd}
\]
交换. 显然, (\SGARefLink{sga7.sga7i.eviii.n2-3-15}{2.3.15}) 的值落在 $1$-闭上链中, 即它是一个 $1$-闭上链, 而且这个闭上链表示所求的
\[
 \mathrm{Ext}^1(nQ,G)=\mathrm{Ext}^1(M_0/M'_0,G).
\]
的元素.

\medskip
\noindent d) 还须具体描述由 $\xi$ 和 $x$ 经图 (\SGARefLink{sga7.sga7i.eviii.n2-3-10}{2.3.10}) 的箭头 (\SGARefLink{sga7.sga7i.eviii.n2-3-5}{2.3.5}) 得到的 $\uExt^1(nQ,G)$ 的截面. 为此, 可以暂不考虑 $P$ 的截面 $x$ 被 $n$ 零化这一事实. 用所选的 $x$ 在 $L_0$ 中的截面提升 $x_0$, 可作如下描述. 这个提升定义增广同态
\[
 u:M_\bullet\longrightarrow L_\bullet\otimes M_\bullet,
 \qquad
 y_0\longmapsto x_0\otimes y_0,
 \quad
 y_1\longmapsto x_0\otimes y_1.
\]
所以由 (\SGARefLink{sga7.sga7i.eviii.n2-3-1}{2.3.1}) 得到的 $\uExt^1(Q,G)$ 的截面, 就是由 $\mathrm{Ext}^1(L_\bullet\otimes M_\bullet,G)$ 的元素 $\xi$ 与前述同态复合所得到的截面. 因而它由次数为 1 的同态
\[
\begin{tikzcd}[column sep=large,row sep=huge]
 0\arrow[r]&M_1\arrow[r,"d_1"]
   \arrow[d,"{y_1\mapsto f_1(x_0\otimes y_1)}"']
   \arrow[rd,phantom,"-" description]
   &M_0\arrow[r]\arrow[d,"{y_0\mapsto f_0(x_0\otimes y_0)}"]&0\\
 0\arrow[r]&C^0\arrow[r,"d^0"]&C^1&\cdots&{.}
\end{tikzcd}
\]
表示. 为将这个类限制到 $nQ\hookrightarrow Q$, 利用同构 $nQ\simeq M_0/M'_0$, 它给出 $nQ$ 的平坦分解
\[
 0\longrightarrow M'_0\longrightarrow M_0\longrightarrow0
\]
包含映射 $nQ\hookrightarrow Q$ 可提升为下列分解之间的同态:
\[
\begin{tikzcd}[column sep=large,row sep=huge]
 0\arrow[r]&M'_0\arrow[r]
   \arrow[d,"{y'_0\mapsto[ny'_0]}"']
   \arrow[rd,phantom,"+" description]
   &M_0\arrow[r]\arrow[d,"{y_0\mapsto ny_0}"]&0\\
 0\arrow[r]&M_1\arrow[r]&M_0\arrow[r]&0&{.}
\end{tikzcd}
\]
取复合后, 所求的 $\uExt^1(nQ,G)$ 的截面就对应于由复形的次数为 1 的同态
\[
\begin{tikzcd}[column sep=large,row sep=large]
 0\arrow[r,dash]&M'_0\arrow[r,dash,"d_1"]\arrow[d,"g_1"']
   &M_0\arrow[r,dash]\arrow[d,"g_0"]&0&{}\\
 0\arrow[r,dash]&C^0\arrow[r,dash]&C^1\arrow[r,dash]&\dots&{,}
\end{tikzcd}
\]
表示的 $\mathrm{Ext}^1(nQ,G)$ 的元素, 其中
\[
 g_1(y'_0)=f_1\bigl(x_0\otimes[ny'_0]\bigr),
 \qquad
 g_0(y_0)=n f_0(x_0\otimes y_0).
\]

\medskip
\noindent e) 现在利用 $x$ 是 $\,{}_nP$ 的截面这一事实, 证明前述 $\mathrm{Ext}^1(nQ,G)$ 的元素是 (\SGARefLink{sga7.sga7i.eviii.n2-3-15}{2.3.15}) 所定义元素的相反数. 对任意同态
\[
 h:M_0\longrightarrow C^0,
\]
若将 $g_1,g_0$ 分别换为
\[
 g'_1=g_1+h d_1,
 \qquad
 g'_0=g_0-d^0h.
\]
则 d) 中考虑的 $\mathrm{Ext}^1$ 的元素不变. 我们可以利用 c) 中构造的同态 $\psi:M_0\longrightarrow C^0$; 按定义, 对 $M'_0$ 的每个截面 $y'_0$ 它满足
\[
 \psi(y'_0)=
 f_1\bigl([nx_0]\otimes y'_0-x_0\otimes[ny'_0]\bigr)
\]
等价地,
\[
 f_1\bigl(x_0\otimes[ny'_0]\bigr)
 \bigl(=g_1(y'_0)\bigr)
 =f_1\bigl([nx_0]\otimes y'_0\bigr)-\psi(y'_0).
\]
定义 $h:M_0\longrightarrow C^0$ 如下:
\[
 h(y_0)=-f_1\bigl([nx_0]\otimes y_0\bigr)+\psi(y_0).
\]
则根据构造,
\[
 g'_1=0,
 \qquad
 g'_0(y_0)=n f_0(x_0\otimes y_0)
 +d^0f_1\bigl([nx_0]\otimes y_0\bigr)-d^0\psi(y_0).
\]
因为 $d^0f_1=-f_0d_1$, 有
\[
 d^0f_1\bigl([nx_0]\otimes y_0\bigr)
 =-f_0(nx_0\otimes y_0),
\]
从而
\[
 g'_1=0,
 \qquad
 g'_0=-d^0\psi.
\]
因此, 所讨论的元素由同态
\[
 M_0/M'_0\simeq nQ\longrightarrow C^1
\]
定义, 这个同态由 $-d^0\psi$ 通过取商诱导. 所以得到的是 (\SGARefLink{sga7.sga7i.eviii.n2-3-15}{2.3.15}) 所定义元素的相反数.

\medskip
\noindent\underline{注 \SGARefTarget{sga7.sga7i.eviii.n2-3-16}{2.3.16}.} 上述证明实际上给出比 \SGARefLink{sga7.sga7i.eviii.n2-3-11}{2.3.11} 所述稍精确的相容性, 方法是选取 $\,{}_nP$ 的一个截面 $x$. 这个截面定义同态
\[
 {}_nQ\longrightarrow G,
 \qquad
 y\longmapsto\varphi_\xi(x,y),
\]
于是经由 (\SGARefLink{sga7.sga7i.eviii.n2-3-6}{2.3.6}) 对应的连接同态得到 $\mathrm{Ext}^1(nQ,G)$ 的一个元素. 另一方面, 如果忘掉 $x$ 被 $n$ 零化这一点, 就可以考虑由 $x$ 定义的同态 $\mathbb Z_T\longrightarrow P$. 它的逆像是 $(\mathbb Z_T,Q)$ 以 $G$ 为核的双扩张, 其类可解释为 $\mathrm{Ext}^1(Q,G)$ 的一个元素, 因而诱导 $\mathrm{Ext}^1(nQ,G)$ 的一个元素. 按这些约定, 在 $\mathrm{Ext}^1(nQ,G)$ 中得到的两个元素\underline{互为相反数}.

% END INCLUDED SGA7I_English_current_buildable_source_20260823/components/101_expose_VIII_proof_2_3_11_parts_c_through_e_and_remark_2_3_16.tex

% BEGIN INCLUDED SGA7I_English_current_buildable_source_20260823/components/102_expose_VIII_section_2_3_17.tex
% SOURCE AUTHORITY:
% expose_VIII_body.tex lines 1222--1241; scan indices 260--262;
% source folios 249--251.
% Direct scan comparison: physical PDF pages 261--263 in overlapping
% 1100-dpi bands.
% Transparent source dispositions:
% - typewritten letter-o indices are normalized to zero;
% - the final reference to phi^n is rendered as phi_xi^n, the pairing that
%   is explicitly being identified throughout 2.3.17.
% Hard stop before 2.3.19, source line 1242.

\medskip
\noindent \SGARefTarget{sga7.sga7i.eviii.n2-3-17}{2.3.17}. 利用 (\SGARefLink{sga7.sga7i.eviii.n2-3-12}{2.3.12}) 所给配对
\[
 \varphi_\xi^n:{}_nP\otimes{}_nQ\longrightarrow G
\]
的显式闭上链描述, 可以得到它的另一个表达式. 首先假设类为 $\xi$ 的双扩张 $E$ 在 $P\times Q$ 上有截面 $s$. 如 VII \SGARefLink{sga7.sga7i.evii.n3-8-13}{3.8.13}, 取 $P$ 的平坦分解, 其中 $L_0=\mathbb Z[P]$, $L_1$ 为典范同态 $\mathbb Z[P]\longrightarrow P$ 的核, 并类似地定义 $M_\bullet$. 此时 $E$ 在 $(L_0,M_0)$ 上的逆像有唯一的双扩张平凡化, 其在 $P\times Q$ 上的限制由 $s$ 给出; 仍用 $s$ 表示这个平凡化 (参见 VII \SGARefLink{sga7.sga7i.evii.n3-8-2}{3.8.2}). 根据 VII \SGARefLink{sga7.sga7i.evii.n3-8-12}{3.8.12}, $\xi$ 在
\[
 \mathrm{Ext}^1\!\left(
 P\mathbin{\overset{\mathbb L}{\otimes}}Q,G\right)
\]
中的类由具有分量 $(f,g)$ 的闭上链表示, 这些分量由
\[
 s(x_1,y_0)=f(x_1\otimes y_0)e_{E/Q}(y_0),
 \qquad
 s(x_0,y_1)=g(x_0\otimes y_1)e_{E/P}(x_0).
\]
刻画. 若 $x$ 和 $y$ 是 $\,{}_nP$ 和 $\,{}_nQ$ 的截面, 将它们提升为 $L_0$ 和 $M_0$ 的相应截面 $x_0=[x]$ 和 $y_0=[y]$. 应用公式 (\SGARefLink{sga7.sga7i.eviii.n2-3-13}{2.3.13}), 并令 $z=s(x,y)$, 得到
\begin{equation}\tag{\SGARefTarget{sga7.sga7i.eviii.n2-3-18}{2.3.18}}
 \varphi_\xi^n(x,y)
 ={}^n z/e_{E/Q}(y)-z^n/e_{E/P}(x).
\end{equation}
这里写在 $z$ 左侧或右侧的指数 $n$, 分别表示对 $E$ 的左侧或右侧乘法结构取 $n$ 次幂. 因而它给出 $E$ 在 $(0,y)$ 或 $(x,0)$ 上的截面; 这样的元素除以 $e_{E/Q}(y)$ 或 $e_{E/P}(x)$ 所得“商”的意义是清楚的.

只要 $z$ 是 $E$ 的整体截面, 且提升 $\,{}_nP,\,{}_nQ$ 的截面 $x,y$, 前述公式 (\SGARefLink{sga7.sga7i.eviii.n2-3-18}{2.3.18}) 就有意义. 我断言, (\SGARefLink{sga7.sga7i.eviii.n2-3-18}{2.3.18}) 对每个这样的截面都成立, 无须假设 $E$ 在 $P\times Q$ 上有截面. $P\times Q$ 的截面 $(x,y)$ 总能局部提升为 $E$ 的截面, 因为 $E\longrightarrow P\times Q$ 是满态射; 因此这个公式能在所有情形下计算 $\varphi_\xi^n(x,y)$. 此外, 若把 $z$ 换成另一个提升 $(x,y)$ 的 $E$ 的截面 $z'$, 则有 $z'=z+g$, 其中 $g$ 是 $G$ 的截面. (\SGARefLink{sga7.sga7i.eviii.n2-3-18}{2.3.18}) 右侧对应于 $z'$ 的两项, 就是对应于 $z$ 的两项各加上 $ng$. 因此它们的差不依赖于 $z$, 从而 (\SGARefLink{sga7.sga7i.eviii.n2-3-18}{2.3.18}) 定义了一个完全确定的同态
\[
 {}_nP\otimes{}_nQ\longrightarrow G.
\]

还须证明这个同态确实就是 $\varphi_\xi^n$. 为此, 立即归结为
\[
 P=Q=(\mathbb Z/n\mathbb Z)_T.
\]
的情形. 这时 $P\times Q$ 具有 $I_T$ 的形式, 其中 $I=\mathbb Z/n\mathbb Z\times\mathbb Z/n\mathbb Z$ 为有限集. 此时 $G_{P\times Q}$ 下的挠子 $E$ 可与一族 $G$ 下的挠子 $E_i$ ($i\in I$) 等同, 它显然在 $T$ 上局部平凡. 因此归结为 $E$ 在 $P\times Q$ 上有整体截面的已讨论情形.

% END INCLUDED SGA7I_English_current_buildable_source_20260823/components/102_expose_VIII_section_2_3_17.tex

% BEGIN INCLUDED SGA7I_English_current_buildable_source_20260823/components/103_expose_VIII_alternative_proof_2_3_19.tex
% SOURCE AUTHORITY:
% expose_VIII_body.tex lines 1242--1342; scan indices 262--263;
% source folios 251--252.
% Direct scan comparison: physical PDF pages 263--264 in overlapping
% 1100-dpi bands.  The three-square homotopy display and the full 2.3.10
% derived-Hom diagram were additionally inspected in five targeted crops at
% 5000-dpi equivalent magnification.
% Transparent source dispositions:
% - the source's typewritten [o] is normalized to the cohomological [0];
% - the circled source mark "- ?" in the large diagram is retained exactly
%   as an unresolved sign annotation, rather than silently adjudicated.
% Hard stop before section 2.4, source line 1346 / scan index 264.

\medskip
\noindent \SGARefTarget{sga7.sga7i.eviii.n2-3-19}{2.3.19}. \underline{\SGARefLink{sga7.sga7i.eviii.n2-3-11}{2.3.11} 的另一证明} (Deligne 给出).

\medskip
\noindent a) 归结为 $P=\mathbb Z/n\mathbb Z$ 及 $T=(\mathrm{Sets})$ 的情形.

\medskip
\noindent b) 设 $\partial$ 为从 $nQ$ 到 $\,{}_nQ[1]$ 的如下箭头:
\[
 nQ\xleftarrow[\text{``$n$''}]{\ \sim\ }
 [\,{}_nQ\longrightarrow Q\,]
 \longrightarrow{}_nQ[1].
\]
下图反交换:
\[
\begin{tikzcd}[row sep=huge,column sep=huge]
 nQ[0]\arrow[r,"1"]\arrow[d,"\partial"']
   &{[Q\xrightarrow{n}Q]}\\
 {}_nQ[1]\arrow[ru,"1"']&{.}
\end{tikzcd}
\]
事实上, 箭头
\[
\begin{tikzcd}[column sep=small,row sep=large]
 {[{}_nQ}\arrow[r]\arrow[dd,"0"']&{Q]}\arrow[dd,"n"]
   &&{[{}_nQ}\arrow[r]\arrow[dd,"1"']&{Q]}\arrow[dd,"0"']
   &\text{互为相反数}
   &{[{}_nQ}\arrow[r]&{Q]}\arrow[ddl,"1"]\\
 &&\text{以及}&&&\text{借助}\quad H\,:&&\\
 {[Q}\arrow[r,"n"]&{Q]}&&{[Q}\arrow[r,"n"]&{Q]}
   &&{[Q}\arrow[r,"n"]&{Q]}
\end{tikzcd}
\]

\medskip
\noindent c) 图 \SGARefLink{sga7.sga7i.eviii.n2-3-10}{2.3.10} 为
\[
\begin{tikzcd}[column sep=small,row sep=small]
 &&&H^1\mathrm{RHom}
   (\mathbb Z\mathbin{\overset{\mathbb L}{\otimes}}Q,G)
   \arrow[ddd]\\
 &&&+\\
 H^1\mathrm{RHom}
   (\mathbb Z/n\mathbb Z\mathbin{\overset{\mathbb L}{\otimes}}Q,G)
   \arrow[rr]\arrow[dd]\arrow[uurrr]
 &&H^1\mathrm{RHom}
   (\mathbb Z/n\mathbb Z,\mathrm{RHom}(Q,G))
   \arrow[dd]&\\
 &&&\mathrm{Ext}^1(Q,G)\arrow[ddl]\\
 \mathrm{Hom}(\mathrm{Tor}_1(\mathbb Z/n\mathbb Z,Q),G)\arrow[d]
 &|[draw,circle,inner sep=1pt]|{-\,?}
 &\mathrm{Hom}(\mathbb Z/n\mathbb Z,\mathrm{Ext}^1(Q,G))
   \arrow[d]&\\
 \mathrm{Hom}({}_nQ,G)\arrow[rr]
 &&\mathrm{Ext}^1(nQ,G)&
\end{tikzcd}
\]
而且:

\smallskip
\noindent c1) 有同构
\[
 \mathbb Z/n\mathbb Z
 \mathbin{\overset{\mathbb L}{\otimes}}Q
 \simeq[Q\xrightarrow{n}Q],
\]
在此同构下, 第一个复合竖直箭头由
\[
 {}_nQ[1]\longrightarrow[Q\xrightarrow{n}Q].
\]
诱导.

\smallskip
\noindent c2) 第二个水平箭头是上面 b) 所给箭头 $\partial$ 的转置.

\smallskip
\noindent c3) 若沿顺时针方向经过此图, 显然得到由
\[
 nQ[0]\longrightarrow[Q\xrightarrow{n}Q].
\]
诱导的箭头. 这就证明了 \SGARefLink{sga7.sga7i.eviii.n2-3-11}{2.3.11}.

% END INCLUDED SGA7I_English_current_buildable_source_20260823/components/103_expose_VIII_alternative_proof_2_3_19.tex

% BEGIN INCLUDED SGA7I_English_current_buildable_source_20260823/components/104_expose_VIII_section_2_4.tex
% SOURCE AUTHORITY:
% expose_VIII_body.tex lines 1346--1380; scan index 264; source folio 253.
% Direct scan comparison: physical PDF page 265 in four overlapping
% 1100-dpi bands.
% Hard stop before section 3, source line 1386 / scan index 265.

\medskip
\noindent \SGARefTarget{sga7.sga7i.eviii.n2-4}{2.4}. 设
\[
 f:\underline T'\longrightarrow\underline T
\]
为拓扑斯的态射. 则有典范同态 (VII \SGARefLink{sga7.sga7i.evii.n2-8}{2.8})
\[
 \Biext^1(P,Q;G)\longrightarrow\Biext^1(P',Q';G'),
\]
其中撇号表示逆像群. 因为对 $\underline T$ 的每个交换群 $E$, 显然都有典范同构
\[
 ({}_nE)'\simeq{}_n(E'),
\]
所以 $\underline T$ 和 $\underline T'$ 上的同态 (\SGARefLink{sga7.sga7i.eviii.n2-2-2}{2.2.2}) 给出典范图
\begin{equation}\tag{\SGARefTarget{sga7.sga7i.eviii.n2-4-1}{2.4.1}}
\begin{tikzcd}[column sep=huge,row sep=huge]
 \Biext^1(P,Q;G)\arrow[r]\arrow[d,"{\varphi^n_{P,Q}}"']
   &\Biext^1(P',Q';G')\arrow[d,"{\varphi^n_{P',Q'}}"]\\
 \Hom({}_nP\otimes{}_nQ,G)\arrow[r]
   &\Hom({}_nP'\otimes{}_nQ',G').
\end{tikzcd}
\end{equation}
同样, 同态 (\SGARefLink{sga7.sga7i.eviii.n2-2-3}{2.2.3}) 给出用 $T_\ell$ 表示的类似图. 我断言这些图都是\underline{交换的}. 这一本质上直接由定义可得, 验证留给读者. 更一般地, 注意本节和前一节所考虑的\underline{所有}典范同态, 都在显然的意义下\underline{与拓扑斯的变换相容}. 此后将使用这类相容性而不再另作说明.


% END INCLUDED SGA7I_English_current_buildable_source_20260823/components/104_expose_VIII_section_2_4.tex

% BEGIN INCLUDED SGA7I_English_current_buildable_source_20260823/components/105_expose_VIII_sections_3_1_and_3_2.tex
% SOURCE AUTHORITY:
% expose_VIII_body.tex lines 1386--1427; scan indices 265--266;
% source folios 254--255.
% Direct scan comparison: physical PDF pages 266--267 in four overlapping
% 1100-dpi bands per page.
% Hard stop before section 3.3, source line 1432 / scan index 267.

\subsection*{\SGARefTarget{sga7.sga7i.eviii.n3}{3}. 群概形 $(P,Q)$ 以 $\underline G_m$ 为核的双扩张: 一般性质}

本节始终固定一个底概形 $S$, 并按第 0 节所说, 在 $S$ 的 fppf 拓扑斯中工作.

\medskip
\noindent \SGARefTarget{sga7.sga7i.eviii.n3-1}{3.1}. 首先回顾群 $\Biext^1(P,Q;\underline G_m)$ 最重要的例子: $P$ 和 $Q$ 是 $S$ 上的两个阿贝尔概形, 而 $\Biext^1$ 解释为 $(P,Q)$ 上的除子对应群的情形 (VII \SGARefLink{sga7.sga7i.evii.n2-9-4}{2.9.4}). 此时前一节的 $\ell$ 进配对就是阿贝尔概形理论中熟知的配对 [1].

\medskip
\noindent \SGARefTarget{sga7.sga7i.eviii.n3-2}{3.2}. 更一般地, 仅假设 $Q$ 为阿贝尔概形. 众所周知 (VII \SGARefLink{sga7.sga7i.evii.n1-3-8}{1.3.8}),
\begin{equation}\tag{\SGARefTarget{sga7.sga7i.eviii.n3-2-1}{3.2.1}}
 \uHom(Q,\underline G_m)=0,
\end{equation}
因而 (\SGARefLink{sga7.sga7i.eviii.n1-4}{1.4}) 有典范同构
\begin{equation}\tag{\SGARefTarget{sga7.sga7i.eviii.n3-2-2}{3.2.2}}
 \Biext^1(P,Q;\underline G_m)\simeq\Hom(P,Q^*),
\end{equation}
其中
\begin{equation}\tag{\SGARefTarget{sga7.sga7i.eviii.n3-2-3}{3.2.3}}
 Q^*=\uExt^1(Q,\underline G_m)
\end{equation}
是 $Q$ 的对偶阿贝尔概形. ($Q^*$ 的可表性在一般情形下由 M. Raynaud 证明, 参见 [1].)

当 $P$ 也为阿贝尔概形时, 与 (\SGARefLink{sga7.sga7i.eviii.n3-2-2}{3.2.2}) 对称地, 有典范同构
\begin{equation}\tag{\SGARefTarget{sga7.sga7i.eviii.n3-2-4}{3.2.4}}
 \Biext^1(P,Q;\underline G_m)\simeq\Hom(Q,P^*).
\end{equation}
因而, 给定 $(P,Q)$ 以 $\underline G_m$ 为核的双扩张 $\xi$, 等价于给定同态
\[
 u:P\longrightarrow Q^*
\]
也等价于给定同态
\[
 v:Q\longrightarrow P^*.
\]

回顾阿贝尔概形的同态 $w:A\longrightarrow B$ 为同构, 当且仅当对每个整数 $n>0$, 所诱导的同态 ${}_nw:{}_nA\longrightarrow{}_nB$ 都是同构. 在上述 $u$ 和 $v$ 的情形, ${}_nu:{}_nP\longrightarrow{}_nQ^*$ 可与配对
\[
 \varphi_\xi^n:{}_nP\times{}_nQ\longrightarrow\underline G_m
\]
所定义的同态 ${}_nP\longrightarrow D({}_nQ)$ 等同 (参见 \SGARefLink{sga7.sga7i.eviii.n2-3-11}{2.3.11}), 而 ${}_nv:{}_nQ\longrightarrow{}_nP^*$ 则与配对
\[
 \psi_\xi^n:{}_nP\times{}_nQ\longrightarrow\underline G_m.
\]
所定义的同态 ${}_nQ\longrightarrow D({}_nP)$ 等同. 现在 $\varphi_\xi^n=-\psi_\xi^n$ (\SGARefLink{sga7.sga7i.eviii.n2-1-10}{2.1.10}), 因此在相差一个符号的意义下, ${}_nu$ 和 ${}_nv$ 互为转置. 因为 ${}_nP$ 和 ${}_nQ$ 在 $S$ 上有限局部自由, 从而满足 Cartier 对偶, 可知 ${}_nu$ 是同构当且仅当 ${}_nv$ 是同构. 因此 $u$ 是同构当且仅当 $v$ 是同构.

一个等价而稍欠对称的表述如下. 考虑 $(P,P^*)$ 以 $\underline G_m$ 为核的典范双扩张 (Weil 双扩张), 它对应于恒等同态 $v:P^*\longrightarrow P^*$. 则相应的典范同态
\[
 u:P\longrightarrow P^{**}
\]
是一个\underline{同构}. 这就是 Barsotti 和 Cartier 的\underline{双重对偶定理}.


% END INCLUDED SGA7I_English_current_buildable_source_20260823/components/105_expose_VIII_sections_3_1_and_3_2.tex

% BEGIN INCLUDED SGA7I_English_current_buildable_source_20260823/components/106_expose_VIII_section_3_3.tex
% R3 explicit revisions R3-VIII-04, R3-VIII-05; see qa/FIDELITY-R3-CHANGES.json.
% SOURCE AUTHORITY:
% expose_VIII_body.tex lines 1432--1481; scan indices 267--270;
% source folios 256--259.
% Direct scan comparison: physical PDF pages 268--271 in four overlapping
% 1100-dpi bands per page.
% Transparent source dispositions:
% - dualizing 0 -> T -> M -> L -> 0 gives the type-correct terminal D(T),
%   not the repeated D(L) in the working transcription;
% - the visibly blank citation after "SGA 3 X 4.9 et" is not guessed; the
%   dangling conjunction is omitted in English.
% Hard stop before Proposition 3.4, source line 1483 / scan index 270.

\medskip
\noindent \SGARefTarget{sga7.sga7i.eviii.n3-3}{3.3}. 现在转向 $Q$ 在 $S$ 上有限局部自由, 或者是 $S$ 上的乘法型群概形的情形. 在此情形中, 关键引理是下列熟知的结果.

\medskip
\noindent\underline{命题 \SGARefTarget{sga7.sga7i.eviii.n3-3-1}{3.3.1}.} \underline{设 $Q$ 是 $S$ 上的交换群概形, 并假设 $Q$ 在 $S$ 上有限局部自由, 或者 $Q$ 在 $S$ 上为有限型且为乘法型 (SGA 3 IX). 则}
\[
 \uExt^1(Q,\underline G_m)=0.
\]

\noindent\underline{证明.} 分两种情形.

\medskip
\noindent a) 假设 $Q$ 在 $S$ 上有限局部自由. 须证明 $Q$ 以 $\underline G_m$ 为核的每个扩张 $E$ 都在 fppf 拓扑下局部分裂. 显然可以假设 $S$ 仿射; 于是存在整数 $n>0$ 使 $nQ=0$ (SGA 3 IX 6.1). 此外, $n\cdot\id$ 在 fppf 拓扑下是 $\underline G_m$ 的满态射, 因而 Kummer 正合列
\[
 0\longrightarrow\mu_n\longrightarrow\underline G_m
 \xrightarrow{\ n\ }\underline G_m\longrightarrow0
\]
给出 $\Ext^i(Q,-)$ 的正合列, 特别地给出\underline{满态射}
\[
 \Ext^1(Q,\mu_n)\longrightarrow
 {}_n\Ext^1(Q,\underline G_m)=\Ext^1(Q,\underline G_m).
\]
最后一个等式成立, 是因为乘以 $n$ 在 $Q$ 上为零, 从而在 $\Ext^1(Q,\underline G_m)$ 上也为零. 因此扩张 $E$ 来自 $Q$ 以 $\mu_n$ 为核的扩张 $F$.

群 $F$ 由 $S$ 上的有限局部自由概形表示 (fpqc 下降理论, SGA 1 VII). 由有限局部自由群的 Cartier 对偶的已知性质 (SGA 3 VII 3),
\[
 D(F)=\uHom(F,\underline G_m)
\]
是 $D(\mu_n)$ 以 $D(Q)$ 为核的扩张. 这里用到如下事实: 若 $A\longrightarrow B$ 是 $S$ 上交换有限局部自由群的单态射, 则 $D(B)\longrightarrow D(A)$ 忠实平坦. 由于 $D(B)$ 和 $D(A)$ 本身在 $S$ 上有限局部自由, 忠实平坦的纤维判别法把问题归结为 $S$ 是域 $k$ 的谱的情形. 此时断言即为仿射坐标环的同态
\[
 k(D(A))\longrightarrow k(D(B))
\]
是单射 (SGA 3 $\mathrm{VI}_B$ 11.14). 这个同态是 $k(B)\longrightarrow k(A)$ 的转置; 后者是满射, 因为 $A\longrightarrow B$ 是闭浸入 (SGA 3 $\mathrm{VI}_B$ 1.4.2). 因而 $k(D(A))\longrightarrow k(D(B))$ 的确是单射.

包含映射 $\mu_n\longrightarrow\underline G_m$ 可与
\[
 D(\mu_n)=\uHom(\mu_n,\underline G_m),
\]
的一个截面等同, 因而只须证明这个包含映射在 $S$ 上对于 fppf 拓扑局部地可延拓为同态 $F\longrightarrow\underline G_m$; 等价地, 只须证明刚才考虑的 $D(\mu_n)$ 的截面在 fppf 拓扑下局部可提升为 $D(F)$ 的截面. 这一点是清楚的, 因为 $D(F)\longrightarrow D(\mu_n)$ 忠实平坦且有限表现.

\medskip
\noindent b) 由 SGA 3 X 4.5, 将 $S$ 换为某个平展覆盖后, 可以假设 $Q$ 可对角化, 记为 $Q=D(M)$, 其中 $M$ 是有限生成的 $\mathbb Z$-模. 设 $T$ 是 $M$ 的挠子模, 令 $L=M/T$, 从而 $L\simeq\mathbb Z^r$. 则正合列
\[
 0\longrightarrow T\longrightarrow M\longrightarrow L
 \longrightarrow0
\]
给出正合列 (SGA 3 IX 3.1)
\[
 0\longrightarrow D(L)\longrightarrow D(M)\longrightarrow D(T)
 \longrightarrow0.
\]
因此只须分别证明 $D(L)$ 和 $D(T)$ 以 $\underline G_m$ 为系数的 $\uExt^1$ 都为零. $D(T)$ 的情形由 a) 得出, 而 $D(L)$ 的情形归结为 $L=\mathbb Z$. 因而还须验证
\[
 \uExt^1(\underline G_m,\underline G_m)=0,
\]
即 $\underline G_m$ 以 $\underline G_m$ 为核的每个扩张 $E$ 都在 fppf 拓扑下局部平凡.

这样的扩张由 $S$ 上的仿射平坦群概形表示 (fpqc 下降理论, SGA 1 VII), 其纤维均为环面; 因而它是一个环面 (SGA 3 X 4.9). 由 SGA 3 X 4.5, 可以假设 $E$ 是分裂环面
\[
 E=D(R)\simeq\underline G_m^2.
\]
作 Zariski 局部化后, 还可以假设单态射 $\underline G_m\hookrightarrow E=D(R)$ 来自一个常值态射 $R\longrightarrow\mathbb Z_S$ (SGA 3 VIII 1.4); 该态射必为满态射 (SGA 3 VIII 3.2 b)). 选取位于 $1\in\mathbb Z$ 上方的 $r\in R$, 就得到分裂该扩张的同态 $E\longrightarrow\underline G_m$, 从而证明了断言.

\medskip
\noindent\underline{推论 \SGARefTarget{sga7.sga7i.eviii.n3-3-2}{3.3.2}.} \underline{设 $Q$ 如 \SGARefLink{sga7.sga7i.eviii.n3-3-1}{3.3.1} 所述. 则对每个 $P$, 有典范同构}
\begin{equation}\tag{\SGARefTarget{sga7.sga7i.eviii.n3-3-2-1}{3.3.2.1}}
 \Biext^1(P,Q;\underline G_m)
 \xleftarrow{\ \sim\ }\Ext^1(P,D(Q)),
\end{equation}
\underline{其中 $D(Q)$ 表示 $Q$ 的 Cartier 对偶 $\uHom(Q,\underline G_m)$. 更确切地说, 典范函子 (\SGARefLink{sga7.sga7i.eviii.n1-1-6}{1.1.6}) 是从 $P$ 以 $D(Q)$ 为核的扩张范畴到 $(P,Q)$ 以 $\underline G_m$ 为核的双扩张范畴的等价.}

这是应用 \SGARefLink{sga7.sga7i.eviii.n3-3-1}{3.3.1} 所得到的 \SGARefLink{sga7.sga7i.eviii.n1-5}{1.5} 的特殊情形.

\medskip
\noindent\underline{注 \SGARefTarget{sga7.sga7i.eviii.n3-3-3}{3.3.3}.} 在 \SGARefLink{sga7.sga7i.eviii.n3-3-1}{3.3.1} 的假设下, $D(Q)$ 由 $S$ 上的分离、局部有限表现且局部拟有限的群概形表示. 更确切地说, 若 $Q$ 在 $S$ 上有限局部自由, 则 $D(Q)$ 也是如此, 这是熟知的 (SGA 3 VII 3). 若 $Q$ 为乘法型, 则 $D(Q)$ 是\underline{有限生成的扭曲常值群} (SGA 3 X 5.1), 这一点已见于 SGA 3 IX 4.5 和 X 5.7. 因此若 $P$ 为群概形, 则 $P$ 以 $D(Q)$ 为核的每个扩张都可表, 这由 fpqc 下降理论的一个熟知引理得出 (SGA 3 X 5.4).

\medskip
\noindent \SGARefTarget{sga7.sga7i.eviii.n3-3-4}{3.3.4}. 对 $S$ 上每个有限表现的仿射平坦群概形 $Q$, 都有 $\uExt^1(Q,\underline G_m)=0$, 这一说法是错误的. 它在
\[
 Q=\underline G_{a_S},\qquad
 S=\Spec k[t]/(t^2),
\]
时就已经不成立, 其中 $k$ 是特征为 $p>0$ 的域.


% END INCLUDED SGA7I_English_current_buildable_source_20260823/components/106_expose_VIII_section_3_3.tex

% BEGIN INCLUDED SGA7I_English_current_buildable_source_20260823/components/107_expose_VIII_section_3_4.tex
% R3 explicit revisions R3-VIII-06; see qa/FIDELITY-R3-CHANGES.json.
% SOURCE AUTHORITY:
% expose_VIII_body.tex lines 1483--1508; scan indices 270--273;
% source folios 259--262.
% Direct scan comparison: physical PDF pages 271--274 in four overlapping
% 1100-dpi bands per page.  The pages are clear at this resolution and no
% higher-resolution ambiguity pass was required.
% Transparent source disposition:
% - in the proof of Proposition 3.4, the printed phrase "une section f de P
%   sur P" is rendered as the mathematically required section of E over P;
%   the following displayed cocycle condition and Lemma 3.4.2 make the
%   intended object unambiguous.
% Hard stop before Corollary 3.5, source line 1510 / scan index 273.

\medskip
\noindent\underline{命题 \SGARefTarget{sga7.sga7i.eviii.n3-4}{3.4}.} \underline{设 $P$ 是 $S$ 上的光滑有限表现群概形, 且纤维连通, $Q$ 是 $S$ 上的环面. 则 $(P,Q)$ 以 $\underline G_m$ 为核的双扩张范畴等价于单点范畴; 也就是说,}
\begin{equation}\tag{\SGARefTarget{sga7.sga7i.eviii.n3-4-1}{3.4.1}}
 \Biext^0(P,Q;\underline G_m)=
 \Biext^1(P,Q;\underline G_m)=0.
\end{equation}

\noindent\underline{证明.} 令 $M=D(Q)$. 这是 $S$ 上的扭曲常值群, 其纤维同构于 $\mathbb Z^r$ 形式的群 (\SGARefLink{sga7.sga7i.eviii.n3-3-3}{3.3.3}). 由 \SGARefLink{sga7.sga7i.eviii.n3-3-2}{3.3.2}, 所讨论的范畴等价于 $P$ 以 $M$ 为核的扩张范畴. 这个范畴是\underline{刚性的}: 每个对象的自同构群都是平凡群. 这就是 (\SGARefLink{sga7.sga7i.eviii.n3-4-1}{3.4.1}) 的第一个等式, 也可写成
\[
 \Hom(P,M)=0.
\]
它立刻由如下事实得出: $M$ 在 $S$ 上平展, 而 $P$ 的纤维连通. 还须证明
\[
 \Ext^1(P,M)=0,
\]
即 $P$ 以 $M$ 为核的每个扩张 $E$ 都平凡.

由存在时的分裂唯一性, 只须证明 $E$ 局部平凡. 因而可以归结为 $M=\mathbb Z_S^r$, 再归结为 $M=\mathbb Z_S$. 当 $S$ 为域的谱时, 结论由下面的 \SGARefLink{sga7.sga7i.eviii.n5-5}{5.5} (i) 得出. (注意在 IX 6 以前我们不会用到 \SGARefLink{sga7.sga7i.eviii.n3-4}{3.4}, 所以尤其不存在循环论证.) 下面将一般情形归结到这一特殊情形.

只须找出 $E$ 在 $P$ 上的截面 $f$ 使 $f(0)=0$, 因为这样的截面必然是乘法性的. 事实上,
\[
 f(x+y)-f(x)-f(y)=0:
\]
左端定义态射 $P\times_S P\longrightarrow\mathbb Z_S$, 它在 $P\times_S P$ 的零截面上为零. 因而它恒为零, 因为 $P\times_S P$ 的纤维连通 ($P$ 的几何纤维连通). 所需结论现由下列引理得出.

\medskip
\noindent\underline{引理 \SGARefTarget{sga7.sga7i.eviii.n3-4-2}{3.4.2}.} \underline{设 $S$ 为概形, $P$ 为光滑有限表现的 $S$-概形且具有几何连通的纤维, $Z$ 为群, $E$ 为 $Z_P$-挠子 (对于 fppf、平展或 fpqc 拓扑; 这些概念在此处等价, 参见 SGA 3 X 5.4). 设 $e$ 为 $P$ 在 $S$ 上的截面, $e'$ 为 $E$ 在 $S$ 上位于 $e$ 上方的截面 (即 $e^{-1}(E)$ 在 $S$ 上的截面).}

\underline{挠子 $E$ 平凡当且仅当对每个 $s\in S$, $P_s$ 上所诱导的挠子 $E_s$ 平凡. 此时 $E$ 在 $P$ 上存在唯一截面 $f$ 使 $f\circ e=e'$.}

\noindent\underline{证明.} 唯一性是清楚的. 两个这样的截面 $f$ 和 $g$ 之“差”为态射
\[
 h:P\longrightarrow Z_S,
 \qquad h\circ e=1,
\]
而这样的态射恒等于 $1$, 因为 $P$ 的纤维连通. 此外, 若 $E$ 在 $P$ 上有截面 $f$, 总可以用一个来自态射 $z:S\longrightarrow
Z_S$ 的态射 $P\longrightarrow Z_S$ 对它作修正, 得到满足 $f'\circ e=e'$ 的截面 $f'$. 只须选取 $z$ 使得
\[
 f(e(t))=e'(t)z(t).
\]

逐纤维应用上述观察, 对每个 $s\in S$, 得到 $E_s$ 在 $P_s$ 上唯一确定的截面 $f_s$, 它将 $e_s$ 映到 $e'_s$. 还须证明这些 $f_s$ 来自 $E$ 在 $P$ 上的一个截面 $f$. 等价地 (EGA IV 17.9.3), 须证明
\[
 U=\bigcup_{s\in S}f_s(P_s)\subset E
\]
在 $E$ 中是开的.

首先假设 $S$ 为诺特概形. 则 $U$ 可构造. 为此, 对每个 $s\in S$, 设 $T$ 为 $s$ 的闭包, 配以诱导的约化结构. 在 $T$ 中存在 $s$ 的开邻域 $T'$, 以及 $E_{T'}$ 在 $P_{T'}$ 上的一个截面, 它延拓 $f_s$, 且与指定的截面 $e$ 和 $e'$ 相容 (这直接由 EGA IV 8.8.2 得出). 由 EGA IV 1.10.1 的判别法, 只须再证明 $U$ 对一般化稳定. 利用 EGA II 7.1.9 的通常论证将问题归结为 $S$ 是离散赋值环的谱的情形.

因为 $P$ 在 $S$ 上光滑, 所以它正则, 从而正规; 又因为它的纤维连通, 所以它连通. 因此, 所给的 $E$ 在 $P$ 关于 $S$ 的一般纤维上的截面, 可唯一延拓为 $E$ 在 $P$ 上的截面 $f$ (EGA IV 18.10.19 和 18.10.20), 它必与 $e$ 和 $e'$ 相容. 因而在这种情形下, $U=f(P)$ 是开的. 这就证明了 $S$ 为诺特概形时的 \SGARefLink{sga7.sga7i.eviii.n3-4}{3.4}. 利用 EGA IV 第 8、9 节中通常的极限论证, 一般情形容易归结于此; 细节留给读者.

\medskip
\noindent\underline{注 \SGARefTarget{sga7.sga7i.eviii.n3-4-3}{3.4.3}.} 引理 3.4.2 可以用博引文献来“一举压倒”, 具体而言, 可引用 SGA 4 XV 1.15 ($n=1$ 的情形) 以及 \SGAExternalRef{sga7.sga7i.eviii.n4-1}{4.1}{SGA4 XV, not VIII}.


% END INCLUDED SGA7I_English_current_buildable_source_20260823/components/107_expose_VIII_section_3_4.tex

% BEGIN INCLUDED SGA7I_English_current_buildable_source_20260823/components/108_expose_VIII_sections_3_5_through_3_7.tex
% SOURCE AUTHORITY:
% expose_VIII_body.tex lines 1510--1553; scan indices 273--275;
% source folios 262--264.
% Direct scan comparison: physical PDF pages 274--276, with pages 275--276
% inspected in four overlapping 1100-dpi bands.  The typography is clear at
% that resolution; no 5000--9000-dpi ambiguity escalation was needed.
% Transparent source dispositions:
% - formula (3.6.2) has Biext^i on both sides.  The printed right-hand
%   exponent 1 is incompatible with the following qualifier i=0,1 and with
%   the asserted equivalence of categories;
% - in the proof of 3.7 the second cyclic case is M=Z/nZ, not the printed
%   phrase "Z=Z/nZ".
% Hard stop before section 4, source line 1559 / scan index 276.

\medskip
\noindent\underline{推论 \SGARefTarget{sga7.sga7i.eviii.n3-5}{3.5}.} \underline{设 $P$ 为 $S$ 上的光滑有限表现群概形, 且纤维连通, 并设}
\[
 0\longrightarrow Q'\longrightarrow Q\longrightarrow Q''
 \longrightarrow0
\]
\underline{为 $S$ 上交换群的正合列, 其中 $Q'$ 为环面. 则双扩张的逆像函子}
\begin{equation}\tag{\SGARefTarget{sga7.sga7i.eviii.n3-5-1}{3.5.1}}
 \operatorname{BIEXT}(P,Q'';\underline G_m)
 \xrightarrow{\ \sim\ }
 \operatorname{BIEXT}(P,Q;\underline G_m)
\end{equation}
\underline{为范畴等价. 特别地, 下列自然同态为同构:}
\begin{equation}\tag{\SGARefTarget{sga7.sga7i.eviii.n3-5-2}{3.5.2}}
 \Biext^i(P,Q'';\underline G_m)
 \xrightarrow{\ \sim\ }
 \Biext^i(P,Q;\underline G_m),
 \qquad i=0,1.
\end{equation}

事实上, 在范畴等价意义下, $\operatorname{BIEXT}(P,Q'';\underline
G_m)$ 可与如下范畴等同: 对象为 $(P,Q)$ 以 $\underline G_m$ 为核的双扩张 $E$, \underline{配以}所诱导的 $(P,Q')$ 以 $\underline G_m$ 为核的双扩张 $E'$ 的一个平凡化 (VII \SGARefLink{sga7.sga7i.evii.n3-7-6}{3.7.6}). 根据 \SGARefLink{sga7.sga7i.eviii.n3-4}{3.4}, 对每个给定的双扩张 $E$, $E'$ 恰有一个这样的平凡化. 当然, 也可以利用与
\[
 0\longrightarrow Q'\longrightarrow Q\longrightarrow Q''
 \longrightarrow0.
\]
相联系的 VII \SGARefLink{sga7.sga7i.evii.n3-7-5}{3.7.5} 中的 $\Biext^i$ 正合列, 直接证明 (\SGARefLink{sga7.sga7i.eviii.n3-5-2}{3.5.2}) 对 $i=0,1$ 都为同构; 这等价于 \SGARefLink{sga7.sga7i.eviii.n3-5}{3.5} 的结论.

\medskip
\noindent\underline{注 \SGARefTarget{sga7.sga7i.eviii.n3-6}{3.6}.} 从 \SGARefLink{sga7.sga7i.eviii.n3-5}{3.5} 交换左右得到的对称命题当然也成立; 通过取对称双扩张 (VII \SGARefLink{sga7.sga7i.evii.n2-7}{2.7}), 它也形式地由 \SGARefLink{sga7.sga7i.eviii.n3-5}{3.5} 得出. 在 \SGARefLink{sga7.sga7i.eviii.n3-5}{3.5} 的数据和假设以外, 再假设给定对称的数据和假设: 即 $Q$ 的纤维连通 (因为 $Q'$ 的纤维连通, 这等价于 $Q''$ 的纤维连通), 且有正合列
\[
 0\longrightarrow P'\longrightarrow P\longrightarrow P''
 \longrightarrow0
\]
其中 $P'$ 为环面. 先应用 \SGARefLink{sga7.sga7i.eviii.n3-5}{3.5}, 再对后一个正合列及 $Q''$ 应用它的对称形式, 可知函子
\begin{equation}\tag{\SGARefTarget{sga7.sga7i.eviii.n3-6-1}{3.6.1}}
 \operatorname{BIEXT}(P'',Q'';\underline G_m)
 \xrightarrow{\ \sim\ }
 \operatorname{BIEXT}(P,Q;\underline G_m)
\end{equation}
为范畴等价. 特别地, 下列同态为同构:
\begin{equation}\tag{\SGARefTarget{sga7.sga7i.eviii.n3-6-2}{3.6.2}}
 \Biext^i(P'',Q'';\underline G_m)
 \xrightarrow{\ \sim\ }
 \Biext^i(P,Q;\underline G_m),
 \qquad i=0,1.
\end{equation}

\medskip
\noindent\underline{命题 \SGARefTarget{sga7.sga7i.eviii.n3-7}{3.7}.} \underline{设 $A$ 是 $S$ 上的阿贝尔概形, $T$ 是 $S$ 上的有限型乘法型群概形 (SGA 3 IX), 并设}
\[
 \underline M=\uHom(T,\underline G_m)
\]
\underline{为 $T$ 的 Cartier 对偶 (由 \SGARefLink{sga7.sga7i.eviii.n3-3-3}{3.3.3}, 它是 $S$ 上有限生成的扭曲常值群). 则分别由 $A$ 以 $T$ 为核的交换扩张和 $(A,\underline M)$ 以 $\underline G_m$ 为核的双扩张构成的范畴 $\operatorname{EXT}(A,T)$ 和 $\operatorname{BIEXT}(A,\underline M;\underline G_m)$ 都是刚性的, 并且典范等价. 它们的同构类群都典范同构于 $\Hom(\underline M,A^*)$, 其中}
\[
 A^*=\uExt^1(A,\underline G_m)
\]
\underline{是 $A$ 的对偶阿贝尔概形 (\SGARefLink{sga7.sga7i.eviii.n3-2}{3.2}). 特别地, 有典范同构}
\begin{equation}\tag{\SGARefTarget{sga7.sga7i.eviii.n3-7-1}{3.7.1}}
 \Ext^1(A,T)\xrightarrow{\ \sim\ }
 \Hom(\underline M,A^*).
\end{equation}

这是取 $P=A$ 和 $Q=\underline M$ 所得到的 \SGARefLink{sga7.sga7i.eviii.n1-6}{1.6} 的特殊情形. 根据双重对偶定理 (SGA 3 X 5.7), 的确有
\[
 T\simeq\uHom(\underline M,\underline G_m),
\]
只须验证条件 (\SGARefLink{sga7.sga7i.eviii.n1-6-1}{1.6.1}), 在此处即为
\[
 \uHom(A,\underline G_m)=0,
 \qquad
 \Hom\bigl(A,\uExt^1(\underline M,\underline G_m)\bigr)=0.
\]
第一个等式已在 (\SGARefLink{sga7.sga7i.eviii.n3-2-1}{3.2.1}) 中回顾; 第二个可加强为
\[
 \uExt^1(\underline M,\underline G_m)=0.
\]
这个断言在 $S$ 上是局部的, 因而可以假设 $\underline
M$ 为常值群, $\underline M=M_S$. 把 $M$ 分解为循环 $\mathbb Z$-模的直和, 就归结为 $M=\mathbb Z$ 和 $M=\mathbb Z/n\mathbb Z$; 前者显然, 后者由 $\underline G_m$ 的 $n$-可除性得出 (若要引经据典, 也可由 \SGARefLink{sga7.sga7i.eviii.n3-3-1}{3.3.1} 得出).

当 $T=\underline G_m$, 从而 $\underline M=\mathbb Z_S$ 时, (\SGARefLink{sga7.sga7i.eviii.n3-7-1}{3.7.1}) 就归结为恒等同构 (\SGARefLink{sga7.sga7i.eviii.n3-2-3}{3.2.3}). 当然, 无须引入群
\[
 \Ext^1\bigl(A\mathbin{\mathop{\otimes}^{\mathbb L}}\underline M,
              \underline G_m\bigr)
 =\Biext^1(A,\underline M;\underline G_m).
\]
就可以直接证明 (\SGARefLink{sga7.sga7i.eviii.n3-7-1}{3.7.1}), 这一本质上是平凡的.


% END INCLUDED SGA7I_English_current_buildable_source_20260823/components/108_expose_VIII_sections_3_5_through_3_7.tex

% BEGIN INCLUDED SGA7I_English_current_buildable_source_20260823/components/109_expose_VIII_sections_4_1_through_4_4.tex
% SOURCE AUTHORITY:
% expose_VIII_body.tex lines 1559--1575; scan indices 276--277;
% source folios 265--266.
% Direct scan comparison: physical PDF pages 277--278 in four overlapping
% 1100-dpi bands.  A targeted crop of the formula in Proposition 4.2 confirms
% that the printed argument is P in all three pullbacks, although the sentence
% has just introduced the line bundle L.  The mathematically type-correct L is
% therefore used below and the product is written explicitly as a tensor
% product.  No other high-resolution ambiguity remained.
% Hard stop before Lemma 4.5, source line 1581 / scan index 278.

\subsection*{\SGARefTarget{sga7.sga7i.eviii.n4}{4}. 域上光滑连通群概形的扩张与双扩张}

回顾下面这个归于 Rosenlicht 的引理 [3 bis, 命题 3].

\medskip
\noindent\underline{引理 \SGARefTarget{sga7.sga7i.eviii.n4-1}{4.1}.} \underline{设 $k$ 为域, $X$ 和 $Y$ 是两个有限型、可分且几何连通的 $k$-概形, 分别配有 $k$-有理点 $e_X$ 和 $e_Y$. 则每个态射}
\[
 X\times Y\longrightarrow\underline G_m
\]
\underline{若在两个子概形 $X\times e_Y$ 和 $e_X\times Y$ 上平凡, 就是平凡的.}

我们只在 $X$ 和 $Y$ 是 $k$ 上的群概形时使用这一结果; 此时假设蕴含 $X$ 和 $Y$ 实际上在 $k$ 上光滑. 通过归纳, \SGARefLink{sga7.sga7i.eviii.n4-1}{4.1} 还蕴含任意多个因子的类似结果. 因而可以应用 VII \SGARefLink{sga7.sga7i.evii.n1-3}{1.3} 和 VII \SGARefLink{sga7.sga7i.evii.n2-9}{2.9} 的结果, 得到以下两个命题.

\medskip
\noindent\underline{命题 \SGARefTarget{sga7.sga7i.eviii.n4-2}{4.2}.} \underline{设 $P$ 是域 $k$ 上的光滑连通群. 从 $P$ 以 $\underline G_m$ 为核的群扩张范畴到 $P$ 上在单位元 $e_P$ 处刚化的线丛范畴的自然函子是全忠实的. 它的本质像由 $P$ 上满足下式的刚化线丛 $\mathbb L$ 构成:}
\[
 \pi^*\mathbb L\simeq
 \pr_1^*\mathbb L\otimes\pr_2^*\mathbb L,
\]
\underline{其中 $\pi:P\times P\longrightarrow P$ 是 $P$ 的群运算, $\pr_1,\pr_2$ 是两个投影.}

\medskip
\noindent\underline{命题 \SGARefTarget{sga7.sga7i.eviii.n4-3}{4.3}.} \underline{设 $P$ 和 $Q$ 是域 $k$ 上的两个光滑连通代数群. 从 $(P,Q)$ 以 $\underline G_m$ 为核的双扩张范畴, 到 $P\times Q$ 上相对于 $P$ 和 $Q$ 的单位截面 $e_P$ 和 $e_Q$ 双刚化的线丛范畴的自然函子是全忠实的; 而且这两个范畴都是刚性的. $P\times Q$ 上的双刚化线丛 $\mathbb L$ 来自双扩张, 当且仅当它沿 $\pi_P\times\id_Q$ 和 $\id_P\times\pi_Q$ 在 $P\times P\times Q$ 和 $P\times Q\times Q$ 上的拉回满足“立方体定理” ([1] 或 [2]).}

最后这个条件就是将 VII \SGARefLink{sga7.sga7i.evii.n2-9-3}{2.9.3} b) 的条件按照 VII \SGARefLink{sga7.sga7i.evii.n3-5}{3.5} b) 所给的显式形式翻译过来.

\medskip
\noindent \SGARefTarget{sga7.sga7i.eviii.n4-4}{4.4}. 我不知道 \SGARefLink{sga7.sga7i.eviii.n4-3}{4.3} 的立方体条件是否自动成立. 至少已知 [2, IV 2.6], 对给定的 $\mathbb L$, 总有一个整数 $n\geq1$, 为 $k$ 的特征指数的幂, 使得 $\mathbb L^{\otimes n}$ 满足该条件, 从而来自 $(P,Q)$ 以 $\underline G_m$ 为核的双扩张. 也总能找到 $k$ 的一个纯不可分扩张 $k'$, 使得所需条件在作基扩张 $k\longrightarrow k'$ 之后成立; 因而在这个扩张之后, $\mathbb L$ 是来自 $(P,Q)$ 以 $\underline G_m$ 为核的双扩张的双刚化线丛.

特别地, 当 $k$ 为完美域时, \SGARefLink{sga7.sga7i.eviii.n4-3}{4.3} 的条件总是成立. 最后这个事实通过迹的论证形式地蕴含前述命题, 也可以应用下面的 \SGARefLink{sga7.sga7i.eviii.n4-7}{4.7} 来证明; 参见 \SGARefLink{sga7.sga7i.eviii.n4-9}{4.9}.


% END INCLUDED SGA7I_English_current_buildable_source_20260823/components/109_expose_VIII_sections_4_1_through_4_4.tex

% BEGIN INCLUDED SGA7I_English_current_buildable_source_20260823/components/110_expose_VIII_sections_4_5_and_4_6.tex
% SOURCE AUTHORITY:
% expose_VIII_body.tex lines 1581--1605; scan indices 278--280;
% source folios 267--269.
% Direct scan comparison: physical PDF pages 279--281 in four overlapping
% 1100-dpi bands per page.  Prose and formulas are fully legible at that
% resolution.
% Transparent source dispositions in the second case of Proposition 4.6:
% - the descent base is Y=P x B, not the printed Y=P x A;
% - after reducing to the case in which Q is abelian, the rigidifying
%   restriction is P x e_Q, not the printed P x e_A.
% Both readings are forced by the displayed morphism to Pic_{Q/k} and by the
% immediately preceding extension 1 -> M -> Q -> B -> 1.
% Hard stop before Proposition 4.7, source line 1607 / scan index 280.

将 EGA $\mathrm{Err}_{\mathrm{IV}}$ 53, 21.4.13 中的以下结果写成引理的形式.

\medskip
\noindent\underline{引理 \SGARefTarget{sga7.sga7i.eviii.n4-5}{4.5}.} \underline{设 $f:X\longrightarrow Y$ 是纤维整的光滑满态射, 其中 $Y$ 正规且整, 一般点为 $y$. 设 $\mathbb L$ 是 $X$ 上的线丛, 其在 $X_y$ 上的限制平凡. 则存在 $Y$ 上的线丛 $\mathbb M$, 使得}
\[
 \mathbb L\simeq f^*\mathbb M.
\]

\medskip
\noindent\underline{命题 \SGARefTarget{sga7.sga7i.eviii.n4-6}{4.6}.} \underline{设 $P,Q$ 是域 $k$ 上的两个光滑连通代数群, 且 $P$ 仿射. 再假设以下两条件之一成立: 或者 $P$ 具有一条合成列, 其每个因子都是环面或同构于加法群 $\underline G_a$ (特别地, 当 $k$ 为完美域时就满足这个条件 [SGA 3, XVII 7.2.1, 4.1.1 (i)$\Longleftrightarrow$(ii), 和 4.1.5]); 或者 $Q$ 是阿贝尔簇 $B$ 以满足前一条件的光滑连通仿射代数群为核的扩张. 在任一假设下, $P\times Q$ 上的每个双刚化线丛 $\mathbb L$ 都平凡; 更有甚者, 由 \SGARefLink{sga7.sga7i.eviii.n4-3}{4.3}, $(P,Q)$ 以 $\underline G_m$ 为核的每个双扩张都平凡.}

\noindent\underline{证明.} 首先注意, 双刚化线丛的每个自同构都是恒等同构 (\SGARefLink{sga7.sga7i.eviii.n4-3}{4.3}). 因而下降理论表明, 只须证明 $\mathbb L$ 在 $k$ 的某个有限可分扩张之后变得平凡. 取极限便立刻归结为 $k$ 可分闭的情形.

这样的域上的每个环面都分裂. 因此若 $P$ 有 \SGARefLink{sga7.sga7i.eviii.n4-6}{4.6} 所述的合成列, 则它也有商群同构于 $\underline G_a$ 或 $\underline G_m$ 的合成列. 反复应用 \SGARefLink{sga7.sga7i.eviii.n4-5}{4.5}, 并利用对每个域 $K$ 都成立的
\[
 \Pic(\underline G_{a_K})=\Pic(\underline G_{m_K})=0
\]
可得
\[
 \Pic(P)=0.
\]
这个等式在底域 $k$ 的每个扩张后仍成立. 对
\[
 \pr_2:P\times Q\longrightarrow Q
\]
及 $\mathbb L$ 应用 \SGARefLink{sga7.sga7i.eviii.n4-5}{4.5}, 找到 $Q$ 上的线丛 $\mathbb M$, 使得 $\mathbb L\simeq\pr_2^*\mathbb M$. 因为 $\mathbb L$ 在 $e_P\times Q$ 上的限制平凡, $\mathbb M$ 平凡, 从而 $\mathbb L$ 也平凡. 这就证明了第一种情形.

在第二种情形, 设 $Q$ 是阿贝尔概形 $B$ 以群 $M$ 为核的扩张, 且该群具有商群同构于 $\underline G_m$ 或 $\underline G_a$ 的合成列. 对
\[
 f:X=P\times Q\longrightarrow Y=P\times B.
\]
应用 \SGARefLink{sga7.sga7i.eviii.n4-5}{4.5}. $P\times B$ 的一般点 $y$ 上的纤维是一个 $M_y$-挠子. 由 $M$ 的特殊结构及对每个域 $K$ 都成立的经典等式
\[
 H^1(K,\underline G_a)=0,
 \qquad
 H^1(K,\underline G_m)=0,
\]
可知它必平凡; 这里取该域为 $k(y)$. 因为 $\Pic(M)=0$, 有 $\Pic(X_y)=0$; 因而 \SGARefLink{sga7.sga7i.eviii.n4-5}{4.5} 适用, 表明 $\mathbb L$ 是 $Y=P\times B$ 上某个线丛 $\mathbb N$ 的拉回.

设 $\mathbb N_1$ 和 $\mathbb N_2$ 分别为 $\mathbb
N$ 在 $P\times e_B$ 和 $e_P\times B$ 上的限制. 将 $\mathbb N$ 换为
\[
 \mathbb N\otimes\pr_1^*(\mathbb N_1)^{-1}
 \otimes\pr_2^*(\mathbb N_2)^{-1},
\]
即可假设 $\mathbb N$ 是双刚化的, 而 $\mathbb L$ 是它的双刚化拉回. 由此归结为 $Q$ 本身就是阿贝尔概形的情形.

此时 $\mathbb L$ 定义典范态射
\begin{equation}\tag{*}
 P\longrightarrow\underline{\Pic}_{Q/k}
\end{equation}
它把单位元映到单位元. 因为 $P$ 连通, 这个态射通过对偶阿贝尔概形
\[
 Q'=\underline{\Pic}^{,0}_{Q/k}.
\]
分解. 众所周知, 从光滑连通仿射代数群到阿贝尔概形的每个概形态射都是常值态射. 这是一个几何性的断言, 因而可假设 $k$ 为代数闭域; 再作逐层分解便立刻归结为 $P=\underline
G_a$ 或 $P=\underline G_m$ 的情形, 而从 $\mathbb P^1_k$ 到 $k$ 上阿贝尔概形的每个有理映射都是常值的.

(*) 为常值, 从而为零, 这意味着 $\mathbb L$ 是 $P$ 上某个线丛 $\mathbb M$ 的拉回. 因为 $\mathbb L$ 在 $P\times e_Q$ 上平凡, $\mathbb M$ 平凡, 从而 $\mathbb L$ 也平凡. 证明完毕.


% END INCLUDED SGA7I_English_current_buildable_source_20260823/components/110_expose_VIII_sections_4_5_and_4_6.tex

% BEGIN INCLUDED SGA7I_English_current_buildable_source_20260823/components/111_expose_VIII_sections_4_7_through_4_10.tex
% R3 explicit revisions R3-VIII-08; see qa/FIDELITY-R3-CHANGES.json.
% SOURCE AUTHORITY:
% expose_VIII_body.tex lines 1607--1696; scan indices 280--284;
% source folios 269--273.
% Direct scan comparison: physical PDF pages 281--284 in overlapping
% 1100-dpi bands.  The prose and displayed formulas are fully legible at
% that resolution; no 5000--9000-dpi ambiguity escalation was needed.
% Source disposition in the proof of Corollary 4.8: the print displays
% Biext^i(P,M;G_m).  This is retained.  Together with (4.7.1), applied with
% Q=M, those vanishings give the Biext^i(A,M;G_m) vanishings required by
% the exact sequence for 0 -> M -> Q -> B -> 0.
% Hard stop before 4.11, source line 1701 / scan index 284.

\medskip
\noindent\underline{命题 \SGARefTarget{sga7.sga7i.eviii.n4-7}{4.7}.} \underline{设 $P$ 和 $Q$ 是域 $k$ 上的两个光滑连通代数群. 假设 $P$ 是阿贝尔概形 $A$ 以光滑连通仿射代数群 $L$ 为核的扩张, 且后者具有合成列, 每个因子都同构于环面或 $\underline
G_a$ (当 $k$ 为完美域时总满足这个条件). 则 \SGARefLink{sga7.sga7i.eviii.n4-3}{4.3} 所考虑的函子为范畴等价; 即 $P\times Q$ 上的每个双刚化线丛 $\mathbb L$ 都来自 $(P,Q)$ 的双扩张. 此外, 刚性范畴之间的函子}
\[
 \operatorname{BIEXT}(A,Q;\underline G_m)
 \longrightarrow
 \operatorname{BIEXT}(P,Q;\underline G_m)
\]
\underline{为范畴等价; 换言之, 它诱导群的同构}
\begin{equation}\tag{\SGARefTarget{sga7.sga7i.eviii.n4-7-1}{4.7.1}}
 \Biext^1(A,Q;\underline G_m)
 \xrightarrow{\ \sim\ }
 \Biext^1(P,Q;\underline G_m).
\end{equation}

对于第一个断言, 下降及 \SGARefLink{sga7.sga7i.eviii.n4-3}{4.3} 再次将问题归结为 $k$ 可分闭的情形. 按 \SGARefLink{sga7.sga7i.eviii.n4-6}{4.6} 的证明中的方法, 由 \SGARefLink{sga7.sga7i.eviii.n4-5}{4.5} 及关于 $L$ 的假设, 可推知 $\mathbb L$ 同构于 $A\times Q$ 上某个双刚化线丛 $\mathbb M$ 的拉回. 因而对 \SGARefLink{sga7.sga7i.eviii.n4-7}{4.7} 的第一个断言, 可归结为 $A=P$ 的情形. 此时线丛 $\mathbb L$ 定义带基点概形的态射
\[
 f:Q\longrightarrow\underline{\Pic}_{P/k},
\]
它必通过 $P$ 的对偶阿贝尔概形 $P'$ 分解. 设 $\mathbb L'$ 为 $P\times P'$ 上的 Weil 层沿
\[
 \id_P\times f:P\times Q\longrightarrow P\times P'.
\]
的拉回. 这是一个双刚化线丛; 根据 $f$ 的定义, 它与 $\mathbb L$ 仅相差一个 $P$ 上的线丛 $\mathbb M$ 的拉回. 因为 $\mathbb L'$ 和 $\mathbb L$ 在 $P\times e_Q$ 上的限制都平凡, 可知 $\mathbb M$ 平凡, 因而 $\mathbb L\simeq\mathbb L'$. 我们已在 VII \SGARefLink{sga7.sga7i.evii.n2-9-4}{2.9.4} 指出, 两个阿贝尔概形的积上的除子对应来自双扩张. 因此 $\mathbb L'$, 从而 $\mathbb L$, 都来自双扩张.

回到 $P$ 为 $A$ 以 $L$ 为核的扩张的一般情形, 还须证明 (\SGARefLink{sga7.sga7i.eviii.n4-7-1}{4.7.1}) 为同构. 这由与
\[
 0\longrightarrow L\longrightarrow P\longrightarrow A
 \longrightarrow0,
\]
相联系的 $\Biext^i$ ($i=0,1$) 正合列以及关系
\[
 \Biext^0(L,Q;\underline G_m)=0,
 \qquad
 \Biext^1(L,Q;\underline G_m)=0.
\]
得出. 第一个关系是 \SGARefLink{sga7.sga7i.eviii.n4-3}{4.3} 中刚性断言的特殊情形, 第二个恰是 \SGARefLink{sga7.sga7i.eviii.n4-6}{4.6} 的第一种情形.

\medskip
\noindent\underline{推论 \SGARefTarget{sga7.sga7i.eviii.n4-8}{4.8}.} \underline{在 \SGARefLink{sga7.sga7i.eviii.n4-7}{4.7} 的假设下, 假设 $Q$ 是代数群 $B$ (它必光滑连通) 以光滑连通仿射代数群 $M$ 为核的扩张. 则刚性范畴之间的函子}
\[
 \operatorname{BIEXT}(A,B;\underline G_m)
 \longrightarrow
 \operatorname{BIEXT}(P,Q;\underline G_m)
\]
\underline{为范畴等价; 换言之, 它诱导群的同构}
\begin{equation}\tag{\SGARefTarget{sga7.sga7i.eviii.n4-8-1}{4.8.1}}
 \Biext^1(A,B;\underline G_m)
 \xrightarrow{\ \sim\ }
 \Biext^1(P,Q;\underline G_m).
\end{equation}

根据 (\SGARefLink{sga7.sga7i.eviii.n4-7-1}{4.7.1}), 还须证明
\[
 \Biext^1(A,B;\underline G_m)
 \xrightarrow{\ \sim\ }
 \Biext^1(A,Q;\underline G_m)
\]
为同构. 与 (\SGARefLink{sga7.sga7i.eviii.n4-7-1}{4.7.1}) 的证明一样, 这由与
\[
 0\longrightarrow M\longrightarrow Q\longrightarrow B
 \longrightarrow0,
\]
相联系的 $\Biext^i$ 正合列以及关系
\[
 \Biext^0(P,M;\underline G_m)=0,
 \qquad
 \Biext^1(P,M;\underline G_m)=0.
\]
得出. 事实上, 以 $M$ 代替 $Q$ 应用 (\SGARefLink{sga7.sga7i.eviii.n4-7-1}{4.7.1}), 就将这些群与第一变元为 $A$ 的相应群等同, 正是前述正合列所需的等同. 上面所示的第二个消失关系, 是交换两个因子后 \SGARefLink{sga7.sga7i.eviii.n4-6}{4.6} 的第二种情形的特殊情形.

\medskip
\noindent\underline{注 \SGARefTarget{sga7.sga7i.eviii.n4-9}{4.9}.} 设 $P$ 和 $Q$ 是\underline{完美}域 $k$ 上的两个光滑连通代数群. 此时已知 [3], 每一个都是阿贝尔概形以光滑连通仿射群为核的扩张, 后者为环面与若干 $\underline G_a$ 的迭代扩张的积 (参见 \SGARefLink{sga7.sga7i.eviii.n4-6}{4.6} 所给参考文献):
\[
 0\longrightarrow L\longrightarrow P\longrightarrow A
 \longrightarrow0,
 \qquad
 0\longrightarrow M\longrightarrow Q\longrightarrow B
 \longrightarrow0.
\]
因此可以应用 \SGARefLink{sga7.sga7i.eviii.n4-8}{4.8}, 表明 $(P,Q)$ 以 $\underline G_m$ 为核的双扩张理论归结为阿贝尔簇对 $(A,B)$ 的类似理论 (后者要经典得多!).

\medskip
\noindent\underline{推论 \SGARefTarget{sga7.sga7i.eviii.n4-10}{4.10}.} \underline{设 $P$ 和 $Q$ 是 $k$ 上的两个光滑连通代数群, $L$ 是 $P$ 的光滑连通仿射代数子群, $E$ 是 $(P,Q)$ 以 $\underline G_m$ 为核的双扩张. 由 \SGARefLink{sga7.sga7i.eviii.n2-2}{2.2}, $E$ 诱导配对}
\[
 T_\ell(P)\times T_\ell(Q)\longrightarrow T_\ell(\underline G_m).
\]
\underline{则所诱导的配对}
\[
 T_\ell(L)\times T_\ell(Q)\longrightarrow T_\ell(\underline G_m)
\]
\underline{为零.}

事实上, 显然可以假设 $k$ 为代数闭域. 此时满足 \SGARefLink{sga7.sga7i.eviii.n4-6}{4.6} 的假设, 从而 $E$ 在 $(L,Q)$ 上的限制为平凡双扩张. 所述结论随即得出, 因为所讨论的诱导配对就是由这一限制所定义的配对.

% END INCLUDED SGA7I_English_current_buildable_source_20260823/components/111_expose_VIII_sections_4_7_through_4_10.tex

% BEGIN INCLUDED SGA7I_English_current_buildable_source_20260823/components/112_expose_VIII_section_4_11.tex
% SOURCE AUTHORITY:
% expose_VIII_body.tex lines 1701--1743; scan indices 284--286;
% source folios 273--275.
% Direct scan comparison: physical PDF pages 285--287 in overlapping
% 1100-dpi bands.  All formulas, primes, kernels, and arrow directions are
% legible at that resolution; no higher-resolution escalation was needed.
% Transparent source disposition: the reference printed as (1.11.1) in the
% middle of the discussion is the pairing (4.11.1), as the local argument
% and every surrounding reference require.
% Hard stop before 4.12, source line 1745 / scan index 286.

\medskip
\noindent \SGARefTarget{sga7.sga7i.eviii.n4-11}{4.11}. 设 $P$ 和 $Q$ 是 $k$ 上的光滑代数群, $E$ 是 $(P,Q)$ 以 $\underline G_m$ 为核的双扩张, $T$ 是 $P$ 的最大子环面. 下面确定配对
\begin{equation}\tag{\SGARefTarget{sga7.sga7i.eviii.n4-11-1}{4.11.1}}
 T_\ell(P)\times T_\ell(Q)
 \longrightarrow T_\ell(\underline G_{m,k}),
\end{equation}
的左零化子, 其中 $\ell$ 仍为某个固定素数. 由 \SGARefLink{sga7.sga7i.eviii.n4-10}{4.10}, 这个零化子总包含 $T_\ell(P)$ 的子群 $T_\ell(T)$.

选取 $k$ 的一个扩张 $k'$, 使 $P$ 和 $Q$ 各自都可写成阿贝尔概形以仿射群概形为核的扩张, 而该仿射群概形是由环面及若干 $\underline G_a$ 构成的迭代扩张. 若 $P$ 和 $Q$ 的幺幂秩为零, 可以取 $k'=k$; 一般地, 可以取 $k'$ 为 $k$ 的完美闭包 (或一个有限纯不可分扩张). 设 $A'$ 和 $B'$ 分别为 $P_{k'}$ 和 $Q_{k'}$ 的阿贝尔部分. 由 \SGARefLink{sga7.sga7i.eviii.n4-7}{4.7}, $(P_{k'},Q_{k'})$ 的双扩张 $E_{k'}$ 来自 $(A',B')$ 的一个双扩张 $F'$, 后者在唯一同构意义下唯一确定. 因而由 (\SGARefLink{sga7.sga7i.eviii.n4-11-1}{4.11.1}) 作底域扩张 $k\to k'$ 得到的配对
\begin{equation}\tag{\SGARefTarget{sga7.sga7i.eviii.n4-11-2}{4.11.2}}
 T_\ell(P_{k'})\times T_\ell(Q_{k'})
 \longrightarrow T_\ell(\underline G_{m,k'}),
\end{equation}
也就是由双扩张 $E_{k'}$ 定义的配对, 可由 $F'$ 所定义的配对
\begin{equation}\tag{\SGARefTarget{sga7.sga7i.eviii.n4-11-3}{4.11.3}}
 T_\ell(A')\times T_\ell(B')
 \longrightarrow T_\ell(\underline G_{m,k'})
\end{equation}
与下列典范同态复合得到 (当 $P$ 和 $Q$ 为 $\ell$-可除时, 它们实际上是满态射):
\begin{equation}\tag{\SGARefTarget{sga7.sga7i.eviii.n4-11-4}{4.11.4}}
 T_\ell(P_{k'})\longrightarrow T_\ell(A'),
 \qquad
 T_\ell(Q_{k'})\longrightarrow T_\ell(B').
\end{equation}

于是 (\SGARefLink{sga7.sga7i.eviii.n4-11-2}{4.11.2}) 的左零化子由 (\SGARefLink{sga7.sga7i.eviii.n4-11-1}{4.11.1}) 的左零化子经基变换得到, 它是在 (\SGARefLink{sga7.sga7i.eviii.n4-11-4}{4.11.4}) 的第一个同态下 (\SGARefLink{sga7.sga7i.eviii.n4-11-3}{4.11.3}) 的左零化子的逆像. 设 $C'$ 为 $A'$ 的最大阿贝尔子概形, 使 $F'$ 在 $(C',B')$ 上的限制为平凡双扩张. 等价地, $C'$ 是 $F'$ 所定义的同态
\begin{equation}\tag{\SGARefTarget{sga7.sga7i.eviii.n4-11-5}{4.11.5}}
 A'\longrightarrow (B')^*
\end{equation}
的核的约化单位连通分支, 其中 $(B')^*$ 是 $B'$ 的对偶阿贝尔概形. (\SGARefLink{sga7.sga7i.eviii.n4-11-3}{4.11.3}) 的左零化子可与 $T_\ell(C')$ 等同; 因而 (\SGARefLink{sga7.sga7i.eviii.n4-11-1}{4.11.1}) 的左零化子恰是
\[
 T_\ell(C')\lhook\joinrel\longrightarrow T_\ell(A')
\]
在 (\SGARefLink{sga7.sga7i.eviii.n4-11-4}{4.11.4}) 的第一个态射下的逆像.

特别地, 说 (\SGARefLink{sga7.sga7i.eviii.n4-11-1}{4.11.1}) 的左零化子恰为 $T_\ell(T)$, 等价地说 (\SGARefLink{sga7.sga7i.eviii.n4-11-2}{4.11.2}) 的左零化子是 (\SGARefLink{sga7.sga7i.eviii.n4-11-4}{4.11.4}) 第一个态射的核 $T_\ell(T')$, 就是说 (\SGARefLink{sga7.sga7i.eviii.n4-11-3}{4.11.3}) 左非退化, 即 $C'=0$, 或等价地, (\SGARefLink{sga7.sga7i.eviii.n4-11-5}{4.11.5}) 具有有限核. 在这最后的表述下, 条件显然不依赖于 $\ell$; 而在任何表述下, 它都显然不依赖于 $k'$ 的选择. 若此条件成立, 就简称所给的 $(P,Q)$ 以 $\underline G_m$ 为核的双扩张 $E$ \underline{左非退化}. 由对偶性, 这也意味着由 $F'$ 所定义的同态
\begin{equation}\tag{\SGARefTarget{sga7.sga7i.eviii.n4-11-6}{4.11.6}}
 B'\longrightarrow (A')^*
\end{equation}
为满态射, 即为满射.

\underline{右非退化}双扩张的概念可对称地定义: (\SGARefLink{sga7.sga7i.eviii.n4-11-6}{4.11.6}) 必须具有有限核, 或等价地, (\SGARefLink{sga7.sga7i.eviii.n4-11-5}{4.11.5}) 必须为满态射. 称 $(P,Q)$ 以 $\underline G_{m,k}$ 为核的双扩张 $E$ 为\underline{非退化的}, 如果它同时左非退化和右非退化. 这等价于说 (\SGARefLink{sga7.sga7i.eviii.n4-11-5}{4.11.5}), 或等价地 (\SGARefLink{sga7.sga7i.eviii.n4-11-6}{4.11.6}), 是同源映射.

由定义立即可见, 这些概念在底域 $k$ 的任意扩张下保持不变.

% END INCLUDED SGA7I_English_current_buildable_source_20260823/components/112_expose_VIII_section_4_11.tex

% BEGIN INCLUDED SGA7I_English_current_buildable_source_20260823/components/113_expose_VIII_sections_4_12_through_4_14.tex
% SOURCE AUTHORITY:
% expose_VIII_body.tex lines 1745--1799; scan indices 286--289;
% source folios 275--278.
% Direct scan comparison: physical PDF pages 287--290 in overlapping
% 1100-dpi bands.  The formulas and prose are fully legible at that
% resolution; no 5000--9000-dpi ambiguity escalation was needed.
% Transparent source dispositions in 4.12:
% - the canonical quotient is f:P -> A, not the printed f:A -> P; this is
%   forced by the displayed pullback Pic(A) -> Pic(P);
% - the descended relation is L \simeq f^*M, not the printed L^* relation;
% - M is a line bundle on A, not the printed "on M".
% In the last sentence of 4.14, the delta map on Pic(P) is (4.14.2), not
% the printed reference (4.14.1).
% Hard stop before section 5, source line 1800 / scan index 289.

\medskip
\noindent \SGARefTarget{sga7.sga7i.eviii.n4-12}{4.12}. 设 $P$ 是光滑连通群概形, $\mathbb L$ 为 $P$ 上的线丛, 并沿用 \SGARefLink{sga7.sga7i.eviii.n4-2}{4.2} 的记号, 令
\begin{equation}\tag{\SGARefTarget{sga7.sga7i.eviii.n4-12-1}{4.12.1}}
 \delta(\mathbb L)=
 \pi^*\mathbb L\otimes
 \pr_1^*(\mathbb L)^{-1}\otimes
 \pr_2^*(\mathbb L)^{-1},
\end{equation}
如 VII \SGARefLink{sga7.sga7i.evii.n1-3-4}{1.3.4} 中已经指出的, 这是 $P\times P$ 上的双刚化线丛. 根据 \SGARefLink{sga7.sga7i.eviii.n4-3}{4.3}, 若 $\delta(\mathbb L)$ 来自某个双扩张, 则该双扩张在唯一同构意义下唯一确定. 特别地, 根据 \SGARefLink{sga7.sga7i.eviii.n4-7}{4.7}, 当 $P$ 是阿贝尔概形 $A$ 以由环面和若干 $\underline G_a$ 构成的迭代扩张为核的扩张时, 这个双扩张有定义. 只要有定义, 这个双扩张显然就是对称的.

在上述有利情形下, 再假设 $P$ 的最大环面分裂, 则 $\delta(\mathbb L)$ 的构造可归结为 $P$ 是阿贝尔概形的情形. 事实上, 该构造对 $P$ 具有函子性, 而由典范商同态 $f:P\longrightarrow A$ 所诱导的同态
\begin{equation}\tag{\SGARefTarget{sga7.sga7i.eviii.n4-12-2}{4.12.2}}
 \Pic(A)\longrightarrow\Pic(P)
\end{equation}
根据 \SGARefLink{sga7.sga7i.eviii.n4-5}{4.5} 为满射 (参见 \SGARefLink{sga7.sga7i.eviii.n4-6}{4.6} 的证明开头). 因而可以找到同构
\[
 \mathbb L\simeq f^*\mathbb M,
\]
其中 $\mathbb M$ 是 $A$ 上的线丛. 则 $\delta(\mathbb L)$ 是 $(A,A)$ 以 $\underline G_{m,k}$ 为核的双扩张 $\delta(\mathbb M)$ 的拉回.

\medskip
\noindent\underline{命题 \SGARefTarget{sga7.sga7i.eviii.n4-13}{4.13}.} \underline{在 \SGARefLink{sga7.sga7i.eviii.n4-12}{4.12} 的假设下, 假设 $\mathbb L$ 为 $P$ 上的丰沛线丛. 若 $(P,P)$ 以 $\underline G_{m,k}$ 为核的双扩张 $\delta(\mathbb L)$ 有定义, 则它在 \SGARefLink{sga7.sga7i.eviii.n4-11}{4.11} 的意义下非退化. 若 $P$ 是阿贝尔概形 $A$ 以满足 \SGARefLink{sga7.sga7i.eviii.n4-7}{4.7} 条件的光滑连通仿射群 $L$ 为核的扩张, 则由 \SGARefLink{sga7.sga7i.eviii.n4-8}{4.8} 给出的、$\delta(\mathbb L)$ 所来自的 $(A,A)$ 的双扩张具有 $\delta(\mathbb M)$ 的形式, 其中 $\mathbb M$ 是 $A$ 上的丰沛线丛.}

归结为 $k$ 代数闭的情形. 根据定义, 沿用 \SGARefLink{sga7.sga7i.eviii.n4-12}{4.12} 的记号, 此时只须证明 $(A,A)$ 的双扩张 $\delta(\mathbb M)$ 非退化. M. Raynaud 在他的学位论文 [2, XI 1.11] 中证明, 若 $\mathbb L=f^*\mathbb M$ 丰沛, 则 $\mathbb M$ 也丰沛. 于是 A. Weil 的理论给出熟知的结论: $\delta(\mathbb M)$ 非退化 [1, 2], 即从 $A$ 到其对偶阿贝尔概形的同态是同源映射.

\medskip
\noindent\underline{注 \SGARefTarget{sga7.sga7i.eviii.n4-14}{4.14}.} 假设处在有利情形, 即 $P$ 是阿贝尔概形 $A$ 以群 $L$ 为核的扩张, 而该群是由环面和若干 $\underline G_a$ 构成的迭代扩张, 并考虑典范复合映射
\begin{equation}\tag{\SGARefTarget{sga7.sga7i.eviii.n4-14-1}{4.14.1}}
 \operatorname{NS}(A)
 \xrightarrow{\ \delta\ }
 \Biext^1(A,A;\underline G_m)
 \xrightarrow{\ \sim\ }
 \Biext^1(P,P;\underline G_m).
\end{equation}
由 \SGARefLink{sga7.sga7i.eviii.n4-8}{4.8}, 第二个箭头为双射. 另一方面, 阿贝尔簇理论表明, 第一个箭头将 $A$ 的 N\'eron--Severi 群 $\operatorname{NS}(A)$ 与
\[
 \Biext^1(A,A;\underline G_m)=\operatorname{Corr}(A,A)
\]
中由对称双扩张的类 (等价地, 对称除子对应的类) 构成的子群等同. 因而 (\SGARefLink{sga7.sga7i.eviii.n4-14-1}{4.14.1}) 将 $\operatorname{NS}(A)$ 与 $(P,P)$ 以 $\underline G_m$ 为核的对称双扩张群等同. 由 \SGARefLink{sga7.sga7i.eviii.n4-7}{4.7}, 这个群也可以描述为 $(P,P)$ 上的对称除子对应的类所构成的群.

当 $(P,P)$ 上的除子对应来自 $\operatorname{NS}(A)$ 的一个丰沛元素时, 自然称之为\underline{丰沛的}; 也同样称相应的 $(P,P)$ 以 $\underline G_m$ 为核的双扩张为\underline{丰沛的}. 这个条件显然在底域扩张下保持不变; 因而即使没有上述关于 $P$ 的限制性假设, 只假设 $P$ 光滑连通, 这个概念也仍然有意义.

进一步假设 $L$ “在 $k$ 上分裂可解”, 即它可由若干 $\underline G_a$ 和 $\underline G_m$ 逐层分解. 上述观察及 (\SGARefLink{sga7.sga7i.eviii.n4-12-2}{4.12.2}) 的满射性表明,
\begin{equation}\tag{\SGARefTarget{sga7.sga7i.eviii.n4-14-2}{4.14.2}}
 \Pic(P)\xrightarrow{\ \delta\ }
 \Biext^1(P,P;\underline G_m)
\end{equation}
的像由对称双扩张构成, 而其核是 $\Pic(P)$ 的子群 $\Pic^0(P)$, 即 $\Pic^0(A)$ 的像. 因此若令
\[
 \operatorname{NS}(P)=\Pic(P)/\Pic^0(P)
 \simeq\operatorname{NS}(A),
\]
也可将 (\SGARefLink{sga7.sga7i.eviii.n4-14-1}{4.14.1}) 解释为典范同构
\begin{equation}\tag{\SGARefTarget{sga7.sga7i.eviii.n4-14-3}{4.14.3}}
 \operatorname{NS}(P)
 \xrightarrow{\ \sim\ }
 \Biext^1(P,P;\underline G_m)^{\mathrm{sym}}.
\end{equation}

最后, \SGARefLink{sga7.sga7i.eviii.n4-13}{4.13} 中引用的 Raynaud 定理说明, $\Pic(P)$ 的元素 $\xi$ 丰沛当且仅当它在 (\SGARefLink{sga7.sga7i.eviii.n4-14-2}{4.14.2}) 下的像 $\delta(\xi)$ 是一个丰沛双扩张的类.

% END INCLUDED SGA7I_English_current_buildable_source_20260823/components/113_expose_VIII_sections_4_12_through_4_14.tex

% BEGIN INCLUDED SGA7I_English_current_buildable_source_20260823/components/114_expose_VIII_section_5_opening_through_5_6.tex
% R3 explicit revisions R3-VIII-01, R3-VIII-02; see qa/FIDELITY-R3-CHANGES.json.
% SOURCE AUTHORITY:
% expose_VIII_body.tex lines 1800--1916; scan indices 289--295;
% source folios 278--284.
% Direct scan comparison: physical PDF pages 290--296 in overlapping
% 1100-dpi bands.  The mathematical text and displayed formulas are fully
% legible at that resolution; no higher-resolution escalation was needed.
% Transparent source dispositions:
% - in the proof of 5.1, G_S and the final E are corrected to Z_S and Z;
% - the source repeats the target category in Corollary 5.2; the first
%   category is restored to torsors on P under Z_P;
% - the repeated P_i at the end of that proof is read as P_{i\eta};
% - in 5.5(ii), the second finite-type example is P, and every stray G_k or
%   G in the statement and proof is corrected to the stated coefficient Z;
% - in 5.6, Phi and Psi are groups of connected components (not identity
%   components), and the stray G_k in (5.6.1) is corrected to Z_k.
% Hard stop before 5.7, source line 1917 / scan index 295.

\medskip
\noindent\underline{5. 以无挠截断常值群概形为核的扩张与双扩张.}

本节和下一节给出一些辅助结果, 将在第 7 节中用于 $Z=\mathbb Z$ 的情形.

\medskip
\noindent\underline{命题 \SGARefTarget{sga7.sga7i.eviii.n5-1}{5.1}.} \underline{设 $S$ 是不可约且几何单支的概形, 一般点为 $\eta$, $Z$ 是无挠群, 从而给出 $S$ 上的常值群概形 $Z_S$. 则 $Z_S$ 下的每个挠子 $E$ 都平凡, 且函子}
\[
 E\longmapsto E(\eta)
\]
\underline{是从 $Z_S$ 下的挠子范畴到 $Z=Z_S(\eta)$ 下的挠子范畴的等价.}

下降理论的一个熟知引理 (SGA 3 X 5.4) 表明, $Z_S$ 下的每个挠子 $E$ 都可表. 事实上, 这一结果适用于 $S$ 上每个分离、局部有限表现且在 $S$ 上局部拟有限的群概形. 因为 $Z_S$ 在 $S$ 上平展, fpqc 下降表明 $Z_S$ 下的每个挠子 $E$ 也如此; 因而这样的挠子在\underline{平展拓扑}下局部平凡 (不仅仅在 fppf 拓扑下如此). 因此, 为证明 $E$ 平凡, 只须证明
\[
 H^1(S_{\mathrm{\acute{e}t}},Z_S)=e.
\]
类似地, \SGARefLink{sga7.sga7i.eviii.n5-1}{5.1} 的全忠实性断言归结为关系
\[
 Z\xrightarrow{\ \sim\ }H^0(S_{\mathrm{\acute{e}t}},Z_S).
\]
后者直接可得. 为证明前者, 考虑包含映射
\[
 i:\eta\longrightarrow S.
\]
由于 $S$ 几何单支, 有
\[
 Z_S\xrightarrow{\ \sim\ }i_*(Z_\eta)
\]
(SGA 4 IX 2.14.1). 于是代替 $i:\eta\to S$ 的 Leray 谱序列的非交换低次正合列 (SGA 4 XII 3.2) 开始为
\[
 e\longrightarrow H^1(S,i_*(Z_\eta))
   \longrightarrow H^1(\eta,Z_\eta)\longrightarrow\cdots,
\]
并表明
\[
 H^1(S,Z_S)\longrightarrow H^1(\eta,Z_\eta)
\]
是单射. 由此归结为 $S$ 是域 $k$ 的谱的情形. 若 $\pi$ 是 $k$ 的可分闭包的 Galois 群, 则
\[
 H^1(k,Z)=\Hom_{\mathrm{cont}}(\pi,Z)/Z=e,
\]
因为 $Z$ 无挠. 证明完毕.

\medskip
\noindent\underline{推论 \SGARefTarget{sga7.sga7i.eviii.n5-2}{5.2}.} \underline{沿用 \SGARefLink{sga7.sga7i.eviii.n5-1}{5.1} 的记号, 设 $P$ 为 $S$ 上的光滑概形. 则 $Z_P$ 下的每个挠子 $E$ 都平凡, 且限制函子}
\[
 E\longmapsto E_\eta
\]
\underline{是从 $P$ 上 $Z_P$ 下的挠子范畴到 $P_\eta$ 上 $Z_{P_\eta}$ 下的挠子范畴的等价.}

事实上, $P$ 几何单支 (EGA IV 17.5.7). 此外, $P$ 的每个不可约分支都支配 $S$, 因为 $P$ 在 $S$ 上平坦, 而 $P$ 的不可约分支集局部有限. 因此 $P$ 是一族不可约且几何单支的概形 $P_i$ 的不交并. 现在对各个 $P_i$ 及其一般纤维 $P_{i\eta}$ 应用 \SGARefLink{sga7.sga7i.eviii.n5-1}{5.1}, 立即得到断言; 这里须注意 $P_i$ 的一般点 $p_i$ 也是 $P_{i\eta}$ 的一般点.

\medskip
\noindent \SGARefTarget{sga7.sga7i.eviii.n5-2-1}{5.2.1}. 前述推论立即给出 $P$ 上 $(Z_P,Z_P)$ 下的双挠子的类似命题: 限制到一般纤维, 诱导从其范畴到 $P_\eta$ 上 $(Z_{P_\eta},Z_{P_\eta})$ 下的双挠子范畴的等价. 事实上, 只须将 $E$ 上的双挠子结构表示为如下数据: 它作为 $Z_P$ 下的右挠子的结构, 以及从 $Z_P$ 到 $Z_P$ 作用于 $E$ 的中心化子群的一个同构. 由于 $E$ 必然平凡, 该中心化子群本身同构于 $Z_P$. 因此在 $P$ 的每个不可约分支 $X$ 上, 归结为以下直接的观察: 若 $X$ 为非空连通概形, 则从集合到 $X$-概形的函子 $Z\mapsto Z_X$ 为全忠实函子.

\medskip
\noindent\underline{命题 \SGARefTarget{sga7.sga7i.eviii.n5-3}{5.3}.} \underline{设 $S$ 和 $Z$ 如 \SGARefLink{sga7.sga7i.eviii.n5-1}{5.1} 所述, $P$ 为 $S$ 上的光滑群概形. 考虑 $P$ 以 $Z_S$ 为核的群概形扩张 $E$. 则限制到一般纤维}
\[
 E\longmapsto E_\eta
\]
\underline{给出从这些扩张的范畴到 $P_\eta$ 以 $Z_\eta$ 为核的扩张范畴的等价. 若 $P$ 和 $Z$ 交换, 它还诱导相应的交换扩张范畴之间的等价.}

首先注意, 由 \SGARefLink{sga7.sga7i.eviii.n5-2}{5.2}, $P$ 以 $Z_S$ 为核的每个群层扩张 (类似地, $P_\eta$ 以 $Z_\eta$ 为核的扩张) 给出平凡挠子, 从而尤其是可表挠子. 因而 \SGARefLink{sga7.sga7i.eviii.n5-3}{5.3} 中的范畴也可以解释为群层扩张的范畴. VII \SGARefLink{sga7.sga7i.evii.n1-1-6}{1.1.6} 和 \SGARefLink{sga7.sga7i.evii.n1-2}{1.2} 中的描述, 将它们表示为 $P$、$P\times P$ 和 $P\times P\times P$ (分别地, $P_\eta$、$P_\eta\times P_\eta$ 和 $P_\eta\times P_\eta\times P_\eta$) 上的双挠子所给的数据, 其对应的群对由 $(Z,Z)$ 经基变换得到. 对光滑 $S$-概形 $P$、$P\times P$ 和 $P\times P\times P$, 在交换情形应用 \SGARefLink{sga7.sga7i.eviii.n5-2}{5.2}, 在一般情形应用 \SGARefLink{sga7.sga7i.eviii.n5-2-1}{5.2.1}, 即得所述结论. 同样的论证给出下列结果.

\medskip
\noindent\underline{推论 \SGARefTarget{sga7.sga7i.eviii.n5-4}{5.4}.} \underline{设 $S,Z$ 如 \SGARefLink{sga7.sga7i.eviii.n5-1}{5.1} 所述, $P,Q$ 为 $S$ 上的两个光滑群概形. 假设 $P,Q,$ 和 $Z$ 交换, 并考虑 $(P,Q)$ 以 $Z_S$ 为核的双扩张 $E$. 则限制到一般纤维}
\[
 E\longmapsto E_\eta
\]
\underline{给出从这些双扩张的范畴到 $(P_\eta,Q_\eta)$ 以 $Z_\eta$ 为核的双扩张范畴的等价.}

上述结果将以无挠常值群为核的扩张或双扩张范畴的确定, 归结为底概形为域的谱的情形. 在该情形中得到以下结论.

\medskip
\noindent\underline{命题 \SGARefTarget{sga7.sga7i.eviii.n5-5}{5.5}.} \underline{设 $k$ 为域, $Z$ 为无挠群, $P$ 为局部有限型的 $k$-群概形, $P^\circ$ 为其单位连通分支, $\Phi=P/P^\circ$ 为其连通分支群 (SGA 3 $\mathrm{VI}_A$ 4). 则:}

\smallskip
\noindent (i) \underline{从作为 $\Phi$ 以 $Z_k$ 为核的扩张 (分别地, 交换扩张) 的群概形范畴, 到作为 $P$ 以 $Z_k$ 为核的扩张 (分别地, 交换扩张) 的群概形范畴的自然函子, 是一个等价. 特别地, 若 $P$ 连通, 则 $P$ 以 $Z_k$ 为核的每个扩张都平凡, 且这类扩张的范畴是刚性的.}

\smallskip
\noindent (ii) \underline{假设 $\Phi$ 为挠群, 且 $P$ 和 $Z$ 交换. 则 $P$ 以 $Z_k$ 为核的扩张范畴是刚性的: 此范畴每个对象的每个自同构都平凡. 此外, 若存在整数 $n>0$ 使 $n\id_\Phi=0$ (例如, 若 $\Phi$ 有限; 当 $P$ 为有限型时就属此情形), 则 $P$ 以 $Z_k$ 为核的交换扩张的类群由典范同构}
\begin{equation}\tag{\SGARefTarget{sga7.sga7i.eviii.n5-5-1}{5.5.1}}
 \Ext^1(P,Z_k)
 \xleftarrow{\ \sim\ }
 \Hom\bigl(\Phi,(Z\otimes_{\mathbb Z}\mathbb Q)/Z\bigr),
\end{equation}
\underline{给出, 这一同构由与正合列}
\begin{equation}\tag{\SGARefTarget{sga7.sga7i.eviii.n5-5-2}{5.5.2}}
 0\longrightarrow Z\longrightarrow Z\otimes_{\mathbb Z}\mathbb Q
 \longrightarrow (Z\otimes_{\mathbb Z}\mathbb Q)/Z
 \longrightarrow 0.
\end{equation}
\underline{相联系的 $\Ext$ 群正合列中的连接同态定义.}

典范态射
\[
 P\longrightarrow\Phi
\]
平坦且纤维不可约, 因而显然在不可约分支集上诱导同构. 根据 \SGARefLink{sga7.sga7i.eviii.n5-1}{5.1}, 从 $Z_\Phi$ 下的挠子到 $Z_P$ 下的挠子的诱导函子因此为等价. (回顾 $k$ 上局部有限型的群概形是几何单支的.) 同样的结论也适用于双挠子以及态射
\[
 P\times P\longrightarrow\Phi\times\Phi,
 \qquad
 P\times P\times P\longrightarrow\Phi\times\Phi\times\Phi.
\]
于是 VII \SGARefLink{sga7.sga7i.evii.n1-1-6}{1.1.6} 对 $P$ 以 $Z_k$ 为核的扩张及 $\Phi$ 以 $Z_k$ 为核的扩张的描述 (交换情形下则用 VII \SGARefLink{sga7.sga7i.evii.n1-3}{1.3}), 给出 (i) 中的断言.

现在只考虑交换扩张. $P$ 以 $Z_k$ 为核的扩张的自同构群为
\[
 \Hom(P,Z_k)\simeq\Hom(\Phi,Z_k),
\]
当 $\Phi$ 为挠群时, 它显然为零. 类似地, 若 $V=Z\otimes_{\mathbb Z}\mathbb Q$ 且 $\Phi$ 被整数 $n>0$ 零化, 则
\[
 \Ext^i(\Phi,V_k)=0\qquad\text{对每个 }i\in\mathbb Z,
\]
因为所讨论的群都是被 $n$ 零化的 $\mathbb Q$-向量空间. 同构 (\SGARefLink{sga7.sga7i.eviii.n5-5-1}{5.5.1}) 立刻随之得出.

同样的方法本质上证明了以下结论.

\medskip
\noindent\underline{推论 \SGARefTarget{sga7.sga7i.eviii.n5-6}{5.6}.} \underline{设 $k$ 为域, $Z$ 为无挠交换群, $P$ 和 $Q$ 为 $k$ 上的两个局部有限型交换群概形, $\Phi=P/P^\circ$ 和 $\Psi=Q/Q^\circ$ 为其连通分支群. 则:}

\smallskip
\noindent (i) \underline{从 $(\Phi,\Psi)$ 以 $Z_k$ 为核的双扩张范畴到 $(P,Q)$ 以 $Z_k$ 为核的双扩张范畴的自然函子是等价. 特别地, 若 $P$ 或 $Q$ 连通, 则 $(P,Q)$ 以 $Z_k$ 为核的双扩张范畴等价于终范畴.}

\smallskip
\noindent (ii) \underline{假设 $\Phi$ 或 $\Psi$ 为挠群. 则 $(P,Q)$ 以 $Z_k$ 为核的双扩张范畴是刚性的. 此外, 若存在整数 $n>0$ 使 $n\id_\Phi=0$ 或 $n\id_\Psi=0$, 则 $(P,Q)$ 以 $Z_k$ 为核的双扩张的类群由典范同构}
\begin{equation}\tag{\SGARefTarget{sga7.sga7i.eviii.n5-6-1}{5.6.1}}
 \Biext^1(P,Q;Z_k)
 \xleftarrow{\ \sim\ }
 \Hom\bigl(\Phi\otimes\Psi,(Z\otimes_{\mathbb Z}\mathbb Q)/Z\bigr),
\end{equation}
\underline{给出, 这一同构由与 (\SGARefLink{sga7.sga7i.eviii.n5-5-2}{5.5.2}) 相联系的 $\Biext^i$ 群 $(i=0,1)$ 的正合列中的连接同态定义.}


% END INCLUDED SGA7I_English_current_buildable_source_20260823/components/114_expose_VIII_section_5_opening_through_5_6.tex

% BEGIN INCLUDED SGA7I_English_current_buildable_source_20260823/components/115_expose_VIII_sections_5_7_through_5_10.tex
% SOURCE AUTHORITY:
% expose_VIII_body.tex lines 1917--1974; scan indices 295--298;
% source folios 284--287.
% Direct scan comparison: physical PDF pages 296--299 in overlapping
% 1100-dpi bands.  The formulas, subscripts, and arrow directions are fully
% legible at that resolution; no higher-resolution escalation was needed.
% Transparent source dispositions:
% - the reference printed as (2.7.2) in 5.7 is the local formula (5.7.2);
% - stray G_{S,T} occurrences in 5.8 and 5.9 are corrected to the group
%   Z_{S,T} fixed by the surrounding definitions and formulas.
% Hard stop before section 6, source line 1976 / scan index 298.

\medskip
\noindent \SGARefTarget{sga7.sga7i.eviii.n5-7}{5.7}. 设 $S$ 为概形, $T$ 为闭子概形, $Z$ 为集合. \underline{在 $S$ 上沿 $T$ 截断、取值为 $Z$ 的常值概形}, 记为 $Z_{S,T}$, 是如下得到的概形: 对每个 $z\in Z$ 取一个 $S$ 的副本 $S_z$, 将它们求和, 并沿开子概形 $U=S\setminus T$ 将所有副本粘合. 因而
\begin{equation}\tag{\SGARefTarget{sga7.sga7i.eviii.n5-7-1}{5.7.1}}
 Z_{S,T}\times_S U\simeq U,
 \qquad
 Z_{S,T}\times_S T\simeq Z_T,
\end{equation}
且 $Z_{S,T}$ 在 $S$ 上平展. 对于可变的 $S$-概形 $S'$, 有对 $S'$ 具有函子性的同构
\begin{equation}\tag{\SGARefTarget{sga7.sga7i.eviii.n5-7-2}{5.7.2}}
 \Hom_S(S',Z_{S,T})\simeq\Hom_T(S'_T,Z_T).
\end{equation}
映射是利用 (\SGARefLink{sga7.sga7i.eviii.n5-7-1}{5.7.1}) 中第二个同构作到 $S'_T$ 的限制; 其为双射的验证留给读者. 若用
\[
 i:T\longrightarrow S
\]
表示包含映射, 则函子性同构 (\SGARefLink{sga7.sga7i.eviii.n5-7-2}{5.7.2}) 也可以写成层的同构
\begin{equation}\tag{\SGARefTarget{sga7.sga7i.eviii.n5-7-3}{5.7.3}}
 Z_{S,T}\simeq i_*(Z_T).
\end{equation}
由这个公式或直接由定义可知, $Z$ 变化时, $Z_{S,T}$ 函子性地依赖于 $Z$, 而且此函子与有限投射极限交换. 特别地, 它将群映为群概形, 将交换群映为交换群概形; 由 $Z_{S,T}$ 所表示的函子的形式 (\SGARefLink{sga7.sga7i.eviii.n5-7-2}{5.7.2}) 也容易看出这一点.

由定义或 (\SGARefLink{sga7.sga7i.eviii.n5-7-3}{5.7.3}) 还可立即看出, $Z_{S,T}$ 的形成与每个基变换 $f:S'\to S$ 交换. 这给出对 $Z$ 具有函子性的同构
\begin{equation}\tag{\SGARefTarget{sga7.sga7i.eviii.n5-7-4}{5.7.4}}
 Z_{S,T}\times_S S'\simeq Z_{S',T'},
 \qquad T'=T\times_S S'=T_{S'}.
\end{equation}

\medskip
\noindent \SGARefTarget{sga7.sga7i.eviii.n5-8}{5.8}. 下面研究形如 $Z_{S,T}$ 的群概形下的挠子 $E$, 其中 $Z$ 为群, 为简单起见在这里假设它交换. 有自然函子
\begin{equation}\tag{\SGARefTarget{sga7.sga7i.eviii.n5-8-0}{5.8.0}}
 E\longmapsto E_T
\end{equation}
从 $Z_{S,T}$ 下的挠子范畴到 $Z_T$ 下的挠子范畴. 我们断言这个函子是\underline{范畴等价}. 它的全忠实性意味着, 对 $Z_{S,T}$ 下的两个挠子 $E,F$, 映射
\[
 \Hom(E,F)\longrightarrow\Hom(E_T,F_T)
\]
为双射. 由 fpqc 下降, 这立刻归结为 $E$ 和 $F$ 平凡的情形, 从而归结为 $E=F=Z_{S,T}$; 此时断言恰是 $S'=S$ 时 (\SGARefLink{sga7.sga7i.eviii.n5-7-2}{5.7.2}) 的双射性.

还须证明 (\SGARefLink{sga7.sga7i.eviii.n5-8-0}{5.8.0}) 本质满. 设 $E_T$ 为 $Z_T$ 下的挠子. 我们已在 \SGARefLink{sga7.sga7i.eviii.n5-1}{5.1} 的证明中指出, $E_T$ 可表, 且在平展拓扑下局部平凡. 因而只须证明
\[
 H^1(S_{\mathrm{\acute{e}t}},Z_{S,T})
 \longrightarrow H^1(T_{\mathrm{\acute{e}t}},Z_T)
\]
为双射. 这由 (\SGARefLink{sga7.sga7i.eviii.n5-7-3}{5.7.3}) 及关系
\[
 R^1i_{\mathrm{\acute{e}t},*}(Z_T)=0,
\]
得出, 后者就是 SGA 4 VIII 5.6 的特殊情形.

\medskip
\noindent \SGARefTarget{sga7.sga7i.eviii.n5-8-1}{5.8.1}. 现在设 $P$ 为 $S$ 上的交换群概形. 我们希望确定作为交换群层的 $P$ 以 $Z_{S,T}$ 为核的扩张范畴. 为此使用 VII \SGARefLink{sga7.sga7i.evii.n1-1-6}{1.1.6} 和 \SGARefLink{sga7.sga7i.evii.n1-3}{1.3} 中的描述, 即用 $P$、$P\times P$ 和 $P\times P\times P$ 上以 $Z_{S,T}$ 的逆像为作用群的挠子来描述. 由 (\SGARefLink{sga7.sga7i.eviii.n5-7-4}{5.7.4}), 这些逆像仍具有与 $Z_{S,T}$ 相同的形式. 对这些挠子应用上述结果, 可知它们对应于 $P_T$、$P_T\times_TP_T$ 和 $P_T\times_TP_T\times_TP_T$ 上以 $Z_T$ 的逆像为作用群的挠子. 因此得到下列结论.

\medskip
\noindent\underline{命题 \SGARefTarget{sga7.sga7i.eviii.n5-9}{5.9}.} \underline{设 $S$ 为概形, $T$ 为闭子概形, $Z$ 为交换群, $Z_{S,T}$ 为 \SGARefLink{sga7.sga7i.eviii.n5-7}{5.7} 中定义的相应的 $S$ 上的截断常值群概形, $P$ 为 $S$ 上的交换群概形. 则限制到 $T$, 诱导如下两个范畴之间的等价: 作为 $P$ 以 $Z_{S,T}$ 为核的扩张的交换群层范畴, 和作为 $P_T$ 以 $Z_T$ 为核的扩张的交换群层范畴.}

本质相同的论证给出下列结论.

\medskip
\noindent\underline{推论 \SGARefTarget{sga7.sga7i.eviii.n5-10}{5.10}.} \underline{沿用 \SGARefLink{sga7.sga7i.eviii.n5-9}{5.9} 的记号, 设 $Q$ 为 $S$ 上另一个交换群概形. 则限制到 $T$, 诱导从 $(P,Q)$ 以 $Z_{S,T}$ 为核的双扩张范畴到 $(P_T,Q_T)$ 以 $Z_T$ 为核的双扩张范畴的等价.}


% END INCLUDED SGA7I_English_current_buildable_source_20260823/components/115_expose_VIII_sections_5_7_through_5_10.tex

% BEGIN INCLUDED SGA7I_English_current_buildable_source_20260823/components/116_expose_VIII_section_6_1_and_lemma_6_2.tex
% SOURCE AUTHORITY:
% expose_VIII_body.tex lines 1976--2003; scan indices 298--300;
% source folios 287--289.
% Direct scan comparison: physical PDF pages 299--301 in overlapping
% 1100-dpi bands.  The prose, formulas, and subscripts are fully legible at
% that resolution; no higher-resolution escalation was needed in this unit.
% Transparent source dispositions:
% - the stray \mathbb Z_Z in the paragraph preceding Lemma 6.2 is corrected
%   to \mathbb Z_S, the free group scheme over the stated base S;
% - the corrupt insertion between "over S" and \overline G in Lemma 6.2 is
%   omitted; the authority clearly defines alpha:S->\mathbb Z_{S,T} and then
%   pulls back the extension \overline G along alpha.
% Hard stop before 6.2.1, source line 2004 / scan index 300.

\medskip
\noindent\underline{6. 以 $G_m$ 的 N\'eron 模型为核的扩张与双扩张.}

\medskip
\noindent \SGARefTarget{sga7.sga7i.eviii.n6-1}{6.1}. 设 $S$ 为概形, $T$ 为闭子概形, $Z$ 为集合, 由此得到 $S$ 上的截断常值概形 $Z_{S,T}$ (\SGARefLink{sga7.sga7i.eviii.n5-7}{5.7}). 有一个对 $Z$ 具有函子性的 $S$ 上平展概形的典范同态
\begin{equation}\tag{\SGARefTarget{sga7.sga7i.eviii.n6-1-1}{6.1.1}}
 Z_S\longrightarrow Z_{S,T},
\end{equation}
由 (\SGARefLink{sga7.sga7i.eviii.n5-7-1}{5.7.1}) 立刻可知, 它是\underline{满}态射. 它对 $Z$ 的函子性表明, 当 $Z$ 为群时, (\SGARefLink{sga7.sga7i.eviii.n6-1-1}{6.1.1}) 是 $S$ 上平展群概形的同态. 因此其核是 $S$ 上的平展群概形, 记为 $Z_{S,U}$, 并有对 $Z$ 具有函子性的典范正合列
\begin{equation}\tag{\SGARefTarget{sga7.sga7i.eviii.n6-1-2}{6.1.2}}
 0\longrightarrow Z_{S,U}\longrightarrow Z_S\longrightarrow Z_{S,T}
 \longrightarrow 0.
\end{equation}
由 (\SGARefLink{sga7.sga7i.eviii.n5-7-1}{5.7.1}) 立即得出, $Z_{S,U}$ 是一族 $S$-概形 $S_n$ ($n\in Z$) 的和, 其中当 $n=0$ 时有 $S_n=S$, 当 $n\ne0$ 时有 $S_n=U$. 由 $Z_{S,U}$ 的这一描述容易看出, 对 $S$ 上的每个群 $F$, 限制到 $U$ 诱导双射
\begin{equation}\tag{\SGARefTarget{sga7.sga7i.eviii.n6-1-3}{6.1.3}}
 \mathrm{Hom}_{\mathrm{gr}}(Z_{S,U},F)
 \xrightarrow{\ \sim\ }
 \mathrm{Hom}_{\mathrm{gr}}(Z_U,F|U).
\end{equation}

\medskip
\noindent \SGARefTarget{sga7.sga7i.eviii.n6-1-4}{6.1.4}. 此后假设 $Z$ 交换, 并研究 $Z_{S,T}$ 以 $S$ 上某个给定交换群 $G$ 为核的交换群扩张. 由正合列 (\SGARefLink{sga7.sga7i.eviii.n6-1-2}{6.1.2}), 这些扩张的范畴等价于如下范畴: $Z_S$ 以 $G$ 为核的扩张 $F$, 配以该扩张沿 $Z_{S,U}\to Z_S$ 的逆像的一个平凡化. 由 (\SGARefLink{sga7.sga7i.eviii.n6-1-3}{6.1.3}), 给出这样的平凡化, 等价于给出 $Z_U$ 以 $G|U$ 为核的扩张 $F|U$ 的一个平凡化.

在后面的讨论中, 取 $Z$ 为整数群 $\mathbb Z$. 因而 $\mathbb Z_S$ 是单生成元自由群, VII \SGARefLink{sga7.sga7i.evii.n1-4}{1.4} 给出从 $\mathbb Z_S$ 以 $G$ 为核的扩张范畴到 $G$ 下的挠子范畴的等价. 对 $\mathbb Z_U$ 和 $G|U$ 应用同一结果, 得到下列结论.

\medskip
\noindent\underline{引理 \SGARefTarget{sga7.sga7i.eviii.n6-2}{6.2}.} \underline{设 $S$ 为概形, $T$ 为闭子概形, $U=S\setminus T$, $G$ 为 $S$ 上的交换群. $S$ 上作为 $\mathbb Z_{S,T}$ (定义于 (\SGARefLink{sga7.sga7i.eviii.n5-7}{5.7})) 以 $G$ 为核的扩张的交换群所成的范畴, 等价于由二元组 $(F,s)$ 构成的范畴, 其中 $F$ 为 $G$ 下的挠子, $s$ 为 $F|U$ 的截面. 若 $\alpha:S\to\mathbb Z_{S,T}$ 是 $\mathbb Z_{S,T}$ 的一个截面, 它是 $\mathbb Z_S$ 在 $S$ 上的截面 $1$ 的像, 则与 $\mathbb Z_{S,T}$ 以 $G$ 为核的扩张 $\overline G$ 对应的挠子为 $G$ 下的 $F=\alpha^*(\overline G)$, 配以由 $\overline G$ 的单位截面限制到 $U$ 所给出的 $F|U$ 的截面.}


% END INCLUDED SGA7I_English_current_buildable_source_20260823/components/116_expose_VIII_section_6_1_and_lemma_6_2.tex

% BEGIN INCLUDED SGA7I_English_current_buildable_source_20260823/components/117_expose_VIII_sections_6_2_1_through_6_2_4.tex
% SOURCE AUTHORITY:
% expose_VIII_body.tex lines 2004--2029; scan indices 300--302;
% source folios 289--291.
% Direct scan comparison: physical PDF pages 301--303 in overlapping
% 1100-dpi bands.  The mathematical text is fully legible at that resolution.
% Transparent source disposition: the printed "description 3.2" at the end
% of 6.2.1 is corrected to the immediately preceding description in 6.2.
% Hard stop before 6.3, source line 2031 / scan index 302.

\medskip
\noindent \SGARefTarget{sga7.sga7i.eviii.n6-2-1}{6.2.1}. 若 $\overline G$ 对应于二元组 $(F,s)$, 则 $\mathbb Z_{S,T}$ 的截面 $f_n$ (由 $\mathbb Z_S$ 的值为 $n\in\mathbb Z$ 的常值截面所定义) 在 $\overline G$ 中的逆像 $f_n^*(\overline G)$, 作为 $G$ 下的挠子, 典范地与 $G$-挠子 $F$ 的 $n$ 次幂 $F^{(n)}$ 等同. 由 (\SGARefLink{sga7.sga7i.eviii.n5-7-1}{5.7.1}) 的第一个关系所定义的 $G_U$-挠子的典范同构
\[
 f_n^*(\overline G)|U\simeq G_U,
\]
由 $F_U^{(n)}$ 的截面 $s^{(n)}$ 所确定, 后者是给定的 $F_U$ 的截面 $s$ 的 $n$ 次幂. 这给出 $\overline G$ 在 $S$ 上的截面的如下具体描述. $\overline G$ 在 $S$ 上的截面 $\overline g$ 由以下数据确定:

\smallskip
\noindent a) $\mathbb Z_{S,T}$ 在 $S$ 上的截面 $\underline n$, 可将它解释为 $\mathbb Z_T$ 的截面, 即 $T$ 上取值于 $\mathbb Z$ 的局部常值函数; 等价地, 它给出分割
\[
 T=\coprod_{i\in\mathbb Z}T_i
\]
将 $T$ 写成以 $i\in\mathbb Z$ 为指标的开子集的不交并;

\smallskip
\noindent b) 一族截面
\[
 \begin{gathered}
  \overline g_i\in\Gamma(F^{(i)}/S_i'),\qquad i\in\mathbb Z,\\
  S_i'=S\setminus T_i',\qquad
  T_i'=\displaystyle\bigcup_{\substack{j\in\mathbb Z\\j\ne i}}T_j,
 \end{gathered}
\]
满足下列相容条件: 对 $i,j\in\mathbb Z$, 利用截面 $s^{(i)}$ 和 $s^{(j)}$ 分别将 $F^{(i)}|U$ 和 $F^{(j)}|U$ 与 $G_U$ 等同后, $\overline g_i$ 和 $\overline g_j$ 在 $S_i'\cap S_j'=U$ 上相同. 换言之, 必须存在一个必然唯一的截面 $h_{ij}\in\Gamma(G/U)$, 使得
\[
 \overline g_i=h_{ij}s^{(i)},\qquad
 \overline g_j=h_{ij}s^{(j)}.
\]
由于 \SGARefLink{sga7.sga7i.eviii.n6-2}{6.2} 对 $\mathbb Z_{S,T}$ 以 $G$ 为核的扩张 $\overline G$ 的描述显然与每个基变换 $S'\to S$ 相容, 上述讨论就完整地描述了层 $\overline G$ 作为变元 $S'\in\operatorname{Ob}(\mathrm{Sch}_{/S})$ 的函子.

\medskip
\noindent \SGARefTarget{sga7.sga7i.eviii.n6-2-2}{6.2.2}. 若概形 $T$ 是一族非空连通概形 $T_i$ ($i\in I$) 的和, 则 \SGARefLink{sga7.sga7i.eviii.n6-2-1}{6.2.1} 中对 $\overline G$ 在 $S$ 上的截面的描述可以简化. 这样的截面由以 $I$ 为指标的一族整数 $(n_i)_{i\in I}$ 和一族截面
\[
 \overline g_i\in\Gamma(F^{(n_i)}/S_i'),\qquad i\in I,
\]
确定, 其中
\[
 S_i'=S\setminus T_i',
 \qquad
 T_i'=\coprod_{\substack{j\in I\\j\ne i}}T_j.
\]
这些 $\overline g_i$ 满足如下相容条件: 对每一对 $i,j\in I$, 利用这些挠子的截面 $s^{(n_i)}$ 和 $s^{(n_j)}$ 分别将 $F^{(n_i)}|U$ 和 $F^{(n_j)}|U$ 等同后, $\overline g_i$ 和 $\overline g_j$ 在 $S_i'\cap S_j'=U$ 上的限制相同.

\medskip
\noindent \SGARefTarget{sga7.sga7i.eviii.n6-2-3}{6.2.3}. 最后, 若假设 $T$ 非空连通, 则给出 $\overline G$ 在 $S$ 上的截面 $\overline g$, 就等于给出整数 $n$ 和 $F^{(n)}$ 在 $S$ 上的截面 $\overline g$. 此时不再需要附加任何相容条件, 因而 $F|U$ 的截面 $s$ 不出现在 $\overline G$ 在 $S$ 上的截面群的描述中. (当然, 一旦要描述 $\overline G$ 在一般 $S$-概形 $S'$ 中取值的点, $s$ 又会出现.)

\medskip
\noindent \SGARefTarget{sga7.sga7i.eviii.n6-2-4}{6.2.4}. 因为 $\overline G$ 是 \SGARefLink{sga7.sga7i.eviii.n6-2-1}{6.2.1} 所述的“开子概形” $f_n^*(\overline G)$ ($n\in\mathbb Z$) 的并, 可知 $\overline G$ 可表当且仅当各挠子 $F^{(n)}$ 可表. 例如, 根据 SGA 1 VIII 2.1 的下降理论, 只要 $G$ 可表且在 $S$ 上仿射就足够了.

% END INCLUDED SGA7I_English_current_buildable_source_20260823/components/117_expose_VIII_sections_6_2_1_through_6_2_4.tex

% BEGIN INCLUDED SGA7I_English_current_buildable_source_20260823/components/118_expose_VIII_sections_6_3_and_6_4.tex
% SOURCE AUTHORITY:
% expose_VIII_body.tex lines 2031--2063; scan indices 302--304;
% source folios 291--293.
% Direct scan comparison: physical PDF pages 303--305 in overlapping
% 1100-dpi bands.  The mathematical text and displayed formulas are fully
% legible at that resolution.
% Transparent source disposition: formula (6.3.2) prints s_D next to
% O_S(-D), while the immediately following defining sentence correctly
% identifies the canonical rational section as s_{-D}; the English formula
% uses s_{-D}.
% Hard stop before Theorem 6.5, source line 2064 / scan index 304.

\medskip
\noindent \SGARefTarget{sga7.sga7i.eviii.n6-3}{6.3}. 后面假设 $G=\underline G_{mS}$. 于是 $\mathbb Z_{S,T}$ 以 $G=\underline G_{mS}$ 为核的每个扩张 $\overline G$ 都可表 (\SGARefLink{sga7.sga7i.eviii.n6-2-4}{6.2.4}), 且当然在 $S$ 上光滑, 因为它是平展群概形 $\mathbb Z_{S,T}$ 以光滑群概形 $\underline G_{mS}$ 为核的扩张. 此外, 给出 $\underline G_m$ 下的挠子 $F$ 等价于给出 $S$ 上的可逆模 $\mathbb L$ (给后者配以 $\mathbb L$ 的可逆截面构成的挠子). 因而 $\overline G$ 由二元组 $(\mathbb L,s)$ 描述, 其中 $\mathbb L$ 是 $S$ 上的可逆模, $s$ 是 $\mathbb L$ 在 $U=S\setminus T$ 上的平凡化, 即 $\mathbb L|U$ 的可逆截面.

当 $U$ 在 $S$ 中概形稠密, 且 $S$ 局部诺特, 或者约化且不可约分支集局部有限时, 给出这样的二元组等价于给出一个 \underline{$S$ 上的除子 $D$, 使得}
\begin{equation}\tag{\SGARefTarget{sga7.sga7i.eviii.n6-3-1}{6.3.1}}
 \operatorname{supp}D\subset T,
\end{equation}
对应的方法是给这样的除子配以二元组
\begin{equation}\tag{\SGARefTarget{sga7.sga7i.eviii.n6-3-2}{6.3.2}}
 \bigl(\underline O_S(-D),s_{-D}|U\bigr)
\end{equation}
它由 $-D$ 所对应的可逆模 $\underline O_S(-D)$ (EGA IV 21.2.8.1), 以及它的典范有理截面 $s_{-D}$ (EGA IV 21.2.11 和 20.2.11(ii)) 在 $U$ 上的限制组成. 无论如何, 不对 $S$ 作任何假设, $S$ 上满足 (\SGARefLink{sga7.sga7i.eviii.n6-3-1}{6.3.1}) 的除子 $D$ 都通过 (\SGARefLink{sga7.sga7i.eviii.n6-3-2}{6.3.2}) 定义二元组 $(\mathbb L,s)$, 从而定义 $\mathbb Z_{S,T}$ 以 $\underline G_m$ 为核的扩张 $\overline G$. 称之为与除子 $D$ 及 $S$ 的满足 (\SGARefLink{sga7.sga7i.eviii.n6-3-1}{6.3.1}) 的闭子概形 $T$ 相联系的 \underline{$\underline G_m$ 的 N\'eron--Raynaud 模型}. 若 $T=\operatorname{supp}D$, 就在名称中省略 $T$.

立即可见, 这个扩张 $\overline G$ 的形成与每个使逆像除子 $f^*(D)$ 有定义的基变换 $f:S'\to S$ 交换 (EGA IV 21.4.2); 特别地, 与每个\underline{平坦}基变换交换.

\medskip
\noindent \SGARefTarget{sga7.sga7i.eviii.n6-3-3}{6.3.3}. \SGARefLink{sga7.sga7i.eviii.n6-2-1}{6.2.1} 中对 $\overline G$ 的截面的一般描述立刻表明, 它们可与 $S$ 上正则亚纯截面层 $\mathcal M_S^*$ (EGA IV 20.1.8) 中满足如下条件的截面 $f$ 等同: 其除子 $\operatorname{div}(f)$ 在 Zariski 拓扑下局部为 $D$ 的整数倍. 若 $V$ 为 $S$ 的开子集, 使 $\operatorname{div}(f)$ 在其上的限制等于 $nD$ 的限制, 其中 $n\in\mathbb Z$, 则 $f$ 所定义的 $\mathbb Z_{S,T}$ 的截面恰是 $\mathbb Z_S$ 的值为 $n$ 的常值截面在 (\SGARefLink{sga7.sga7i.eviii.n6-1-2}{6.1.2}) 中的像. 上面对 $\overline G$ 在 $S$ 上的截面群的描述, 在 $S$ 的每个开子集上仍然成立; 更一般地, 对 $\overline G$ 在任意平坦 $S$-概形 $S'$ 中取值的点群也成立. 因为 $\overline G$ 本身在 $S$ 上平坦 (实际上光滑), 这一描述再次完整地刻画了由 $D$ 和 $T\supset\operatorname{supp}D$ 定义的 $\underline G_m$ 的 N\'eron--Raynaud 模型.

\medskip
\noindent \SGARefTarget{sga7.sga7i.eviii.n6-4}{6.4}. 现在假设 $S$ \underline{正规}, 其不可约分支集局部有限 (从而 $S$ 是一族正规整概形的和), $D$ 为各重数均等于 $1$ 的正除子 (EGA IV 21.6), $T=\operatorname{supp}D$, 且 $T$ 局部不可约, 即其不可约分支两两不交. 则 $S$ 上每个在 $U=S\setminus T$ 上正则的有理函数 $f$, 其除子 $\operatorname{div}(f)$ 局部为 $D$ 的倍数. 因而此时有群的典范同构
\begin{equation}\tag{\SGARefTarget{sga7.sga7i.eviii.n6-4-1}{6.4.1}}
 \Gamma(S,\overline G)\xrightarrow{\ \sim\ }
 \Gamma(U,\underline G_{mU}).
\end{equation}
进一步假设 $T$ 几何单支且局部整, 这一假设在光滑基变换下稳定 (EGA IV 17.5.7). 由于所有假设此时均在基变换下稳定, (\SGARefLink{sga7.sga7i.eviii.n6-4-1}{6.4.1}) 蕴含: \underline{对每个在 $S$ 上光滑的 $S'$, 有对 $S'$ 具有函子性的群的典范同构}
\begin{equation}\tag{\SGARefTarget{sga7.sga7i.eviii.n6-4-2}{6.4.2}}
 \mathrm{Hom}_S(S',\overline G)\xrightarrow{\ \sim\ }
 \mathrm{Hom}_U(S'_U,\underline G_{mU})
 \simeq\Gamma(S'_U,\underline O_{S'_U}^*).
\end{equation}
因为 $\overline G$ 本身在 $S$ 上光滑, 这个同构在唯一同构意义下刻画了 $S$ 上的群概形 $\overline G$. 它还表明, $\overline G$ 起到了 $S$ 的开子集 $U$ 上的群概形 $\underline G_{mU}$ 的“N\'eron 模型”的作用. 再加上它最初由 M.~Raynaud 引入和使用这一事实, 就说明了 \SGARefLink{sga7.sga7i.eviii.n6-3}{6.3} 中给 $\overline G$ 所取名称的由来.


% END INCLUDED SGA7I_English_current_buildable_source_20260823/components/118_expose_VIII_sections_6_3_and_6_4.tex

% BEGIN INCLUDED SGA7I_English_current_buildable_source_20260823/components/119_expose_VIII_theorem_6_5_through_corollary_6_7.tex
% SOURCE AUTHORITY:
% expose_VIII_body.tex lines 2064--2102; scan indices 304--307;
% source folios 293--296.
% Direct scan comparison: physical PDF pages 305--308 in overlapping
% 1100-dpi bands.  The theorem statements, formulas, and proof are fully
% legible at that resolution.
% Transparent source disposition: the proof of essential surjectivity in
% Theorem 6.5 prints "local rings of X" once; the stated base is S, and the
% English corrects X to S.
% Hard stop before Remarks 6.8, source line 2103 / scan index 307.

\medskip
\noindent\underline{定理 \SGARefTarget{sga7.sga7i.eviii.n6-5}{6.5}.} \underline{设 $S$ 是正规局部整概形, $D$ 是 $S$ 上没有多重分支的正除子, 其支集几何单支且局部整. 设 $\overline G$ 为 $S$ 上相对于 $D$ 的 $\underline G_m$ 的 N\'eron--Raynaud 模型 (\SGARefLink{sga7.sga7i.eviii.n6-3}{6.3}), 并考虑 $S$ 上群概形 $\overline G$ 下的挠子 $E$ 的范畴. 限制函子}
\[
 E\longmapsto E_U=E|U
\]
\underline{从这个范畴到 $U$ 上群概形 $\underline G_{mU}$ 下的挠子范畴是全忠实的. 若 $S$ 正则 (更一般地, 若其局部环都是唯一分解环), 则此函子是范畴等价.}

证明所讨论的函子全忠实, 即对 $\overline G$ 下的两个挠子 $E,F$, 限制映射
\begin{equation}\tag{*}
 \mathrm{Hom}(E,F)\longrightarrow\mathrm{Hom}(E_U,F_U)
\end{equation}
是双射, 其中 $\mathrm{Hom}$ 表示\underline{挠子}同态的集合. 首先注意, 这个问题显然\underline{对于 Zariski 拓扑在 $S$ 上是局部的}. 事实上, 假设 $S$ 为一族开子集 $S_i$ 的并, 使得对 $S$ 中包含于某个 $S_i$ 的每个开子集 $S'$ (从而特别地, 对 $S'=S_i$ 或 $S_i\cap S_j$), 映射
\[
 \mathrm{Hom}(E_{S'},F_{S'})
 \longrightarrow
 \mathrm{Hom}(E_{S'\cap U},F_{S'\cap U})
\]
都是双射. 则 (*) 是双射. 另一方面, 将 \SGARefLink{sga7.sga7i.eviii.n5-8}{5.8} 用于 $S,T,\mathbb Z$, 并将 \SGARefLink{sga7.sga7i.eviii.n5-1}{5.1} 用于 $T,\mathbb Z$, 立刻可知 $\overline G$ 下的每个挠子在 Zariski 拓扑下局部平凡. 这两个观察使我们只须在 $E=F=\overline G$ 时证明 (*) 的双射性. 此时该映射恰是 (\SGARefLink{sga7.sga7i.eviii.n6-4-1}{6.4.1}), 其双射性已得证明.

现在证明, 当 $S$ 的局部环都是唯一分解环时, 函子 $E\mapsto E_U$ 本质满. 因为 $\underline G_{mS}$ 与 $\overline G$ 的一个开子群概形等同, 只须证明 $\underline G_{mU}$ 下的每个挠子同构于某个 $\underline G_{mS}$ 下的挠子在 $U$ 上的限制. 等价地, $U$ 上的每个可逆模都必须同构于 $S$ 上某个可逆模的限制. 这由关于 $S$ 的假设得出 (EGA IV 21.6.11, 其中只处理局部诺特的情形; 这也是后文唯一需要的情形). 这就证明了 \SGARefLink{sga7.sga7i.eviii.n6-5}{6.5}.

\medskip
\noindent \SGARefTarget{sga7.sga7i.eviii.n6-5-1}{6.5.1}. 现在进一步假设 $S$ 局部诺特且\underline{正则}. 于是 $S$ 局部为唯一分解的, 给出 $D$ 等价于给出一个约化、纯余维为 $1$、几何单支的闭子概形, 仍记为 $D$. 还应注意, 关于 $S$ 和 $D$ 的假设在每个光滑基变换 $S'\to S$ 下稳定, 因而 \SGARefLink{sga7.sga7i.eviii.n6-5}{6.5} 的结论适用于带有子概形 $D'=D\times_S S'$ 的 $S'$. 此时 VII \SGARefLink{sga7.sga7i.evii.n1-1-6}{1.1.6} 对扩张群的描述立即给出如下推论.

\medskip
\noindent\underline{推论 \SGARefTarget{sga7.sga7i.eviii.n6-6}{6.6}.} \underline{设 $S$ 是局部诺特正则概形, $D$ 是约化、几何单支、纯余维为 $1$ 的闭子概形 (因而在 $S$ 上定义一个除子), $\overline G$ 是相对于 $D$ 的 $\underline G_m$ 的 N\'eron--Raynaud 模型, $P$ 是 $S$ 上的光滑交换群概形. 考虑作为群层的 $P$ 以 $\overline G$ 为核的交换扩张范畴. 限制到 $U$, $E\mapsto E|U$, 诱导此范畴与 $P_U$ 以 $\underline G_{mU}$ 为核的扩张范畴之间的等价.}

\medskip
\noindent \SGARefTarget{sga7.sga7i.eviii.n6-6-1}{6.6.1}. 同样的方法给出不一定交换的扩张的类似命题 (其中不再假设 $P$ 交换). 证明的细节留给读者; 由于必须考虑双挠子而非挠子, 证明稍微复杂一些.

用与 \SGARefLink{sga7.sga7i.eviii.n6-6}{6.6} 本质相同的证明, 同样得到:

\medskip
\noindent\underline{推论 \SGARefTarget{sga7.sga7i.eviii.n6-7}{6.7}.} \underline{设 $S,D,\overline G$ 如 \SGARefLink{sga7.sga7i.eviii.n6-6}{6.6} 所述, $P,Q$ 为 $S$ 上的两个光滑交换群概形. 考虑 $(P,Q)$ 以 $\overline G$ 为核的双扩张范畴. 则从这个范畴到 $(P_U,Q_U)$ 以 $\underline G_{mU}$ 为核的双扩张范畴的、限制到 $U$, $E\mapsto E|U$ 的限制函子, 是范畴等价.}


% END INCLUDED SGA7I_English_current_buildable_source_20260823/components/119_expose_VIII_theorem_6_5_through_corollary_6_7.tex

% BEGIN INCLUDED SGA7I_English_current_buildable_source_20260823/components/120_expose_VIII_remarks_6_8_and_autocritique_6_9.tex
% SOURCE AUTHORITY:
% expose_VIII_body.tex lines 2103--2170; scan indices 307--310;
% source folios 296--299.
% Direct scan comparison: physical PDF pages 308--311 in overlapping
% 1100-dpi bands.  The prose and formulas are fully legible at that
% resolution.  The hand-corrected diagram on physical page 310 was also
% inspected directly in a 5000-dpi targeted crop.
% Transparent source dispositions:
% - 6.9(a) begins with j:V->S but consistently defines U=S-Y; V is corrected
%   to U, and j_*G_{mS} is corrected to j_*G_{mU};
% - the source diagram retains a faint obsolete right arrowhead on the Y--S
%   stroke beneath a heavy hand-added left arrowhead.  The correction gives
%   Y->S.  The mark at U is an open-immersion hook, and the source's terminal
%   period is retained.
% Hard stop before section 7, source line 2171 / scan index 310.

\medskip
\noindent\underline{注 \SGARefTarget{sga7.sga7i.eviii.n6-8}{6.8}.}\quad a) 我们将在 $S$ 为性状、$D$ 为由 $S$ 的闭点 $s$ 到 $S$ 的包含映射所定义的除子时, 应用 \SGARefLink{sga7.sga7i.eviii.n6-6}{6.6} 和 \SGARefLink{sga7.sga7i.eviii.n6-7}{6.7}. 此时 $U$ 是 $S$ 的分式域 $K$ 的谱, $\overline G$ 是 Raynaud 引入的 $K$ 上乘法群的通常的 N\'eron 模型. 注意, 即使在这一特殊情形, \SGARefLink{sga7.sga7i.eviii.n6-6}{6.6} 和 \SGARefLink{sga7.sga7i.eviii.n6-7}{6.7} 的证明也用到了比性状更一般的底上的 \SGARefLink{sga7.sga7i.eviii.n6-5}{6.5} (例如 $P$ 及其笛卡尔幂). 不过, 在这种情形下, 只须知道所讨论的特殊情形中 N\'eron 模型 $\overline G$ 的\underline{存在性}: 这就形式地推出在 $S$ 上光滑的底概形上的 \SGARefLink{sga7.sga7i.eviii.n6-5}{6.5}, 从而推出 \SGARefLink{sga7.sga7i.eviii.n6-6}{6.6} 和 \SGARefLink{sga7.sga7i.eviii.n6-7}{6.7}. 对我们而言, 引入 $\overline G$ 主要是为了给下一节研究以乘法群为核的扩张与双扩张的延拓问题提供一个方便的中间步骤.

\smallskip
\noindent b) 在比上述条件更一般的条件下, 也可以定义 $\underline G_{mU}$ 的 N\'eron--Raynaud 模型 $\overline G$, 使它对每个在 $S$ 上光滑的 $S'$ 都满足基本条件 (\SGARefLink{sga7.sga7i.eviii.n6-4-2}{6.4.2}); 此时只要求 $U$ 是正规局部诺特概形 $S$ 的一个概形稠密开子集. 这时 $\overline G$ 是 fppf \underline{层}, 未必可表. 用 $\underline D$ 表示 $S$ 的小平展景上由 $S$ 上支集包含于 $T=S\setminus U$ 的除子构成的层, 并稍滥用术语, 用同一个字母表示它在 $(\mathrm{Sch})_{/S}$ 上的层的典范延拓. 它在相对概形 $f:S'\to S$ 上的截面就是 $f_{\mathrm{\acute{e}t}}^*(\underline D)$ 在 $S'$ 上的截面. 所求模型为群的扩张
\begin{equation}\tag{\SGARefTarget{sga7.sga7i.eviii.n6-8-1}{6.8.1}}
 0\longrightarrow\underline G_{mS}\longrightarrow\overline G
 \longrightarrow\underline D\longrightarrow0.
\end{equation}
利用一般描述 VII \SGARefLink{sga7.sga7i.evii.n1-1-6}{1.1.6}, 不对 $S$ 作任何正规性假设, 就可以毫无困难地构造典范扩张 (\SGARefLink{sga7.sga7i.eviii.n6-8-1}{6.8.1}); 同时也得到对 $S$ 上平展的 $S'$ 成立的 (\SGARefLink{sga7.sga7i.eviii.n6-4-2}{6.4.2}). 要使同一关系对每个在 $S$ 上光滑的 $S'$ 都成立, 等价于验证层 $\underline D$ 的形成与每个光滑基变换 $f:S'\to S$ 交换. 对正规 $S$, 这由 EGA $\mathrm{Err}_{\mathrm{IV}}$ 53, 21.4.9 得出.

\medskip
\noindent \SGARefTarget{sga7.sga7i.eviii.n6-9}{6.9}.\quad “\underline{自我批评}” (P.~Deligne 的意见).\footnote{本篇作者既已认输, 如有必要, 请读者按照所附指引重写第 5、6 节.}

第 5、6 节确实颇为乏味. 不时会令人觉得, 是在“手工”重做 $i_*,i^*$ 的形式体系, 却看不出哪些是几何内容, 哪些是“一般性形式论证”.

我建议抽出其中的“几何”命题 a)、b)、c):

\smallskip
\noindent a) (参见 \SGARefLink{sga7.sga7i.eviii.n6-8}{6.8} b)) 设 $j:U\hookrightarrow S$ 是正规诺特概形 $S$ 的概形稠密开子集. 则平展层 $\underline G_{mS}$ 可与平展层 $j_*\underline G_{mU}$ 的一个子层等同. 设 $\underline D$ 为正合列
\[
 0\longrightarrow\underline G_{mS}\longrightarrow
 j_*\underline G_{mU}\longrightarrow\underline D\longrightarrow0.
\]
中的商平展层. \underline{则} $\underline D$ 的形成与每个光滑基变换交换 (EGA $\mathrm{Err}_{\mathrm{IV}}$ 53, 21.4.9); \underline{此外}, 若 $S$ 的严格局部环是唯一分解环, 则 $R^1j_*\underline G_m=0$ (定理 90).

\smallskip
\noindent b) 若 $U=S\setminus Y$, 其中 $Y$ 为几何单支除子, 且 $i$ 为 $Y$ 到 $S$ 的包含映射, 则
\[
 \underline D=i_*\mathbb Z_Y.
\]

\smallskip
% SGA-SRC-001273
\noindent c) 若 $Y$ 为概形 $X$ 的闭子概形, 闭浸入为 $i:Y\hookrightarrow X$, 且 $E$ 为集合, 则平展层 $i_*\underline E_Y$ 由 $X$ 上的一个平展概形表示, 其形成与每个基变换交换.

\medskip
\noindent\underline{推论.} \underline{设 $S$ 为正规诺特概形, $Y$ 为几何单支除子, 并设 $U=S\setminus Y$:}
\begin{center}
\begin{tikzcd}[column sep=huge]
 U \arrow[r, hook, "j"'] & S & Y \arrow[l, "i"']
\end{tikzcd}
\qquad .
\end{center}
\underline{存在唯一的 $S$ 上的光滑群概形 $\mathcal G_m$, 使得在每个光滑基变换后都有 $\mathcal G_m|S_{\mathrm{\acute{e}t}}=j_*(\underline G_m)_U$.}

事实上, 上述要求将 $\mathcal G_m$ 确定为 $S$ 上光滑概形所成范畴上的平展层. 根据 a) 和 b), 正合列
\[
 0\longrightarrow\underline G_m\longrightarrow\mathcal G_m
 \xrightarrow{\ q\ }i_*\mathbb Z_Y\longrightarrow0
\]
是正合的. 根据 c), $i_*\mathbb Z_Y$ 可表. 最后, 态射 $q$ 可表.

使用上述“小光滑景”, 应当可以由 a) 形式地推出第 6 节的命题.

% END INCLUDED SGA7I_English_current_buildable_source_20260823/components/120_expose_VIII_remarks_6_8_and_autocritique_6_9.tex

% BEGIN INCLUDED SGA7I_English_current_buildable_source_20260823/components/121_expose_VIII_section_7_opening_and_theorem_7_1_statement.tex
% SOURCE AUTHORITY:
% expose_VIII_body.tex lines 2171--2203; scan indices 310--311;
% source folios 299--300.
% Direct scan comparison: physical PDF pages 311--312 in overlapping
% 1100-dpi bands.  The prose, component-group notation, and obstruction
% formulas are fully legible; no higher-resolution escalation was needed.
% Transparent source dispositions:
% - typewritten letter-o special-fibre subscripts are normalized to zero;
% - in 7.1(b), the source twice drops the eta from E_eta in the extension
%   criterion; the English retains the stated generic-fibre object E_eta;
% - the coefficient is typeset coherently as \underline G_{mS}.
% Hard stop before the proof, source line 2204 / scan index 311.

\medskip
\noindent\underline{7. 以 $\underline G_m$ 为核的扩张与双扩张的典范延拓.}

\medskip
\noindent \SGARefTarget{sga7.sga7i.eviii.n7-0}{7.0}. 现在收获第 5、6 节准备工作的成果. 首先设底概形 $S$ 为一个\underline{性状}, 即离散赋值环的谱; 设 $\eta$ 为其一般点, $s$ 为其闭点, 且 $k=k(s)$. 设 $P$ 为 $S$ 上的光滑交换群概形, 用 $P_0=P_s$ 表示其特殊纤维. 本节只考虑交换群; 特别地, 所考虑的一切群扩张都是交换的.

\medskip
\noindent\underline{定理 \SGARefTarget{sga7.sga7i.eviii.n7-1}{7.1}.} \underline{沿用 \SGARefLink{sga7.sga7i.eviii.n7-0}{7.0} 中回顾的记号.}

\smallskip
\noindent a) \underline{假设 $P$ 的特殊纤维的连通分支群}
\[
 \Phi=P_0/P_0^{\circ},
\]
\underline{是挠群; 特别地, 当 $P$ 在 $S$ 上为有限型时就属此情形. 则从 $P$ 以 $\underline G_{mS}$ 为核的扩张范畴到 $P_\eta$ 以 $\underline G_{m\eta}$ 为核的扩张范畴的函子 $E\mapsto E_\eta$ 为全忠实函子. 假设存在整数 $n>0$ 使 $n\Phi=0$. 则 $P_\eta$ 以 $\underline G_{m\eta}$ 为核的每个扩张 $E_\eta$ 都对应于一个典范同态}
\begin{equation}\tag{\SGARefTarget{sga7.sga7i.eviii.n7-1-1}{7.1.1}}
 d(E_\eta):\Phi\longrightarrow(\mathbb Q/\mathbb Z)_k,
\end{equation}
\underline{它为零是 $E_\eta$ 同构于某个 $P$ 以 $\underline G_{mS}$ 为核的扩张的一般纤维的充分必要条件. 特别地, 若 $P_0$ 连通, 则上述函子 $E\mapsto E_\eta$ 是范畴等价.}

\smallskip
\noindent b) \underline{设 $P$ 和 $Q$ 是 $S$ 上的两个光滑群概形, 令}
\[
 \Phi=P_0/P_0^{\circ},\qquad \Psi=Q_0/Q_0^{\circ},
\]
\underline{并假设 $\Phi$ 或 $\Psi$ 为挠群 (例如, $P$ 或 $Q$ 在 $S$ 上为有限型). 则从 $(P,Q)$ 以 $\underline G_{mS}$ 为核的双扩张范畴到 $(P_\eta,Q_\eta)$ 以 $\underline G_{m\eta}$ 为核的双扩张范畴的函子 $E\mapsto E_\eta$ 为全忠实函子. 假设存在整数 $n>0$ 使 $n\Phi=0$ 或 $n\Psi=0$. 则 $(P_\eta,Q_\eta)$ 以 $\underline G_{m\eta}$ 为核的每个双扩张 $E_\eta$ 都对应于一个典范配对}
\begin{equation}\tag{\SGARefTarget{sga7.sga7i.eviii.n7-1-2}{7.1.2}}
 d(E_\eta):\Phi\otimes_{\mathbb Z}\Psi
 \longrightarrow(\mathbb Q/\mathbb Z)_k,
\end{equation}
\underline{它为零是 $E_\eta$ 同构于某个 $(P,Q)$ 以 $\underline G_{mS}$ 为核的双扩张 $E$ 的一般纤维的充分必要条件. 特别地, 若 $P_0$ 或 $Q_0$ 连通, 则上述函子 $E\mapsto E_\eta$ 是范畴等价.}


% END INCLUDED SGA7I_English_current_buildable_source_20260823/components/121_expose_VIII_section_7_opening_and_theorem_7_1_statement.tex

% BEGIN INCLUDED SGA7I_English_current_buildable_source_20260823/components/122_expose_VIII_theorem_7_1_proof.tex
% SOURCE AUTHORITY:
% expose_VIII_body.tex lines 2204--2242; scan indices 311--313;
% source folios 300--302.
% Direct scan comparison: physical PDF pages 312--314 in overlapping
% 1100-dpi bands.  Every displayed Ext/Biext formula and tensor product is
% fully legible at that resolution.
% Transparent source dispositions:
% - part (a) prints "biextensions by \overline G" where Theorem 6.6 and the
%   entire argument concern extensions; the English uses extensions;
% - the final target printed Q/Z is the same (Q/Z)_k coefficient used in
%   (7.1.2), 5.6, and the preceding Ext calculation, so the English restores
%   (\mathbb Q/\mathbb Z)_k.
% Hard stop before Remark 7.2, source line 2243 / scan index 313.

\medskip
\noindent\underline{证明.}\quad a) 设 $\overline G$ 为 $\underline G_m$ 的 N\'eron--Raynaud 模型. 所讨论的函子是“核的扩张 $\underline G_{mS}\to\overline G$”函子与定义在以 $\overline G$ 为核的扩张范畴上的“限制到 $\eta$”函子的复合. 由 \SGARefLink{sga7.sga7i.eviii.n6-6}{6.6}, 后一个函子是范畴等价, 因而只须研究第一个函子. 为此使用群概形的正合列
\begin{equation}\tag{\SGARefTarget{sga7.sga7i.eviii.n7-1-3}{7.1.3}}
 0\longrightarrow\underline G_{mS}\longrightarrow\overline G
 \longrightarrow\mathbb Z_{S,s}\longrightarrow0,
\end{equation}
其中 $\mathbb Z_{S,s}$ 的记号取自 \SGARefLink{sga7.sga7i.eviii.n5-7}{5.7}. 与 (\SGARefLink{sga7.sga7i.eviii.n7-1-3}{7.1.3}) 相联系的 $\mathrm{Ext}^i$ ($i=0,1$) 正合列表明, 所讨论的函子全忠实当且仅当
\begin{equation}\tag{\SGARefTarget{sga7.sga7i.eviii.n7-1-4}{7.1.4}}
 \mathrm{Hom}(P,\mathbb Z_{S,s})=0,
\end{equation}
并且 $P$ 以 $\overline G$ 为核的扩张 $\overline E$ 来自 $P$ 以 $\underline G_{mS}$ 为核的扩张 $E$ 的障碍, 落在
\begin{equation}\tag{\SGARefTarget{sga7.sga7i.eviii.n7-1-5}{7.1.5}}
 \mathrm{Ext}^1(P,\mathbb Z_{S,s}).
\end{equation}
中. 由 (\SGARefLink{sga7.sga7i.eviii.n5-7-2}{5.7.2}), (\SGARefLink{sga7.sga7i.eviii.n7-1-4}{7.1.4}) 中第一个群同构于 $\mathrm{Hom}(P_0,\mathbb Z_k)$, 后者又同构于 $\mathrm{Hom}(\Phi,\mathbb Z_k)$. 因而 (\SGARefLink{sga7.sga7i.eviii.n7-1-4}{7.1.4}) 等价于
\begin{equation}\tag{7.1.4 bis}
 \mathrm{Hom}(\Phi,\mathbb Z_k)=0.
\end{equation}
如果 $\Phi$ 为挠群, 这个关系当然成立. 此外, 若 $n\Phi=0$, 则 \SGARefLink{sga7.sga7i.eviii.n5-9}{5.9} 和 \SGARefLink{sga7.sga7i.eviii.n5-5}{5.5} 给出典范同构
\begin{equation}\tag{\SGARefTarget{sga7.sga7i.eviii.n7-1-6}{7.1.6}}
 \mathrm{Ext}^1(P,\mathbb Z_{S,s})
 \simeq\mathrm{Ext}^1(P_0,\mathbb Z_k)
 \simeq\mathrm{Hom}\bigl(\Phi,(\mathbb Q/\mathbb Z)_k\bigr),
\end{equation}
这就证明了 \SGARefLink{sga7.sga7i.eviii.n7-1}{7.1}(a).

\smallskip
\noindent b) 证明本质上相同. 由 \SGARefLink{sga7.sga7i.eviii.n6-7}{6.7}, 归结为研究“核的扩张 $\underline G_{mS}\to\overline G$”函子, 它从 $(P,Q)$ 以 $\underline G_{mS}$ 为核的双扩张范畴映到 $(P,Q)$ 以 $\overline G$ 为核的双扩张范畴. 使用正合列 (\SGARefLink{sga7.sga7i.eviii.n7-1-3}{7.1.3}) 及相联系的 $\mathrm{Biext}^i$ ($i=0,1$) 正合列 (VII \SGARefLink{sga7.sga7i.evii.n3-7-5}{3.7.5}). 函子的全忠实性等价于
\begin{equation}\tag{\SGARefTarget{sga7.sga7i.eviii.n7-1-7}{7.1.7}}
 \mathrm{Hom}(P\otimes Q,\mathbb Z_{S,s})=0.
\end{equation}
由 (\SGARefLink{sga7.sga7i.eviii.n5-7-2}{5.7.2}), 左端同构于 $\mathrm{Hom}(P_0\otimes Q_0,\mathbb Z_k)$, 从而同构于 $\mathrm{Hom}(\Phi\otimes\Psi,\mathbb Z_k)$; 这个群当 $\Phi$ 或 $\Psi$ 为挠群时的确为零. 另一方面, 对给定的 $(P,Q)$ 以 $\overline G$ 为核的双扩张 $\overline E$, 它来自 $(P,Q)$ 以 $\underline G_{mS}$ 为核的双扩张 $E$ 的障碍, 落在
\[
 \mathrm{Biext}^1(P,Q;\mathbb Z_{S,s}).
\]
中. 由 \SGARefLink{sga7.sga7i.eviii.n5-10}{5.10}, 这同构于 $\mathrm{Biext}^1(P_0,Q_0;\mathbb Z_k)$, 当 $n\Phi=0$ 或 $n\Psi=0$ 时, 后者又由 \SGARefLink{sga7.sga7i.eviii.n5-6}{5.6} 同构于
\[
 \mathrm{Hom}\bigl(\Phi\otimes\Psi,(\mathbb Q/\mathbb Z)_k\bigr).
\]
这就证明了 \SGARefLink{sga7.sga7i.eviii.n7-1}{7.1}.


% END INCLUDED SGA7I_English_current_buildable_source_20260823/components/122_expose_VIII_theorem_7_1_proof.tex

% BEGIN INCLUDED SGA7I_English_current_buildable_source_20260823/components/123_expose_VIII_remark_7_2.tex
% SOURCE AUTHORITY:
% expose_VIII_body.tex lines 2243--2261; scan indices 313--314;
% source folios 302--303.
% Direct scan comparison: physical PDF pages 314--315 in overlapping
% 1100-dpi bands.  The obstruction map, codimension bound, and footnote are
% fully legible; no higher-resolution escalation was needed.
% Transparent source disposition: spacing/typographical slips in "for every"
% and the products P x P x P are normalized without changing content.
% Hard stop before 7.3, source line 2262 / scan index 314.

\medskip
\noindent\underline{注 \SGARefTarget{sga7.sga7i.eviii.n7-2}{7.2}.} 给定一般点为 $\eta$ 的连通正则局部诺特概形 $S$ 上的光滑群概形 $P$, 可以研究从 $P$ 以 $\underline G_{mS}$ 为核的扩张范畴到 $P_\eta$ 以 $\underline G_{m\eta}$ 为核的扩张范畴的函子 $E\mapsto E_\eta$. 这个问题容易归结为 \SGARefLink{sga7.sga7i.eviii.n7-1}{7.1}(a) 所处理的问题. 为简单起见, 假设 $P$ 在 $S$ 上为有限型. 由 \SGARefLink{sga7.sga7i.eviii.n7-1}{7.1} 容易得出这个函子全忠实.\footnote{为此, $S$ 正规就已足够, 不必要求正则.} 容易推出, 若 $E_\eta$ 是 $P_\eta$ 以 $\underline G_{m\eta}$ 为核的扩张, 则它同构于某个 $P$ 以 $\underline G_{mS}$ 为核的扩张 $E$ 的一般纤维, 当且仅当对每个满足
\[
 \dim\operatorname{Spec}(\underline O_{S,s})=1,
\]
的 $s\in S$, 通过基变换 $\operatorname{Spec}(\underline O_{S,s})\to S$ 所得到的类似延拓问题均有解. 根据 \SGARefLink{sga7.sga7i.eviii.n7-1}{7.1}(a), 这表现为某个同态
\begin{equation}\tag{\SGARefTarget{sga7.sga7i.eviii.n7-2-1}{7.2.1}}
 d(E_\eta,s):P_s/P_s^{\circ}
 \longrightarrow(\mathbb Q/\mathbb Z)_{k(s)}.
\end{equation}
为零.

上述断言显然归结为以下断言. 设 $U$ 为 $S$ 的开子集, 其补集 $T$ 满足
\[
 \operatorname{codim}(T,S)\geq2.
\]
则从 $P$ 以 $\underline G_{mS}$ 为核的扩张范畴到 $P_U$ 以 $\underline G_{mU}$ 为核的扩张范畴的函子 $E\mapsto E|U$ 是范畴等价. 用 VII \SGARefLink{sga7.sga7i.evii.n1-1-6}{1.1.6} 的描述, 这个断言本质上平凡: 对每个光滑 $S$-概形 $X$ (例如 $P$、$P\times P$、$P\times P\times P$ 等), $X$ 局部为唯一分解的, 且 $X_T=X\setminus X_U$ 在 $X$ 中的余维至少为 $2$. 因此, 从 $\underline G_{mX}$ 下的挠子范畴到 $\underline G_{mX_U}$ 下的挠子范畴的函子 $F\mapsto F|X_U$ 为范畴等价.

对于双扩张的延拓问题有完全类似的结果; 其表述留给读者.


% END INCLUDED SGA7I_English_current_buildable_source_20260823/components/123_expose_VIII_remark_7_2.tex

% BEGIN INCLUDED SGA7I_English_current_buildable_source_20260823/components/124_expose_VIII_section_7_3_1.tex
% SOURCE AUTHORITY:
% expose_VIII_body.tex lines 2262--2314; scan indices 314--315;
% source folios 303--304.
% Direct scan comparison: physical PDF pages 315--316 in overlapping
% 1100-dpi bands. Both triangular diagrams were checked node by node and
% arrow by arrow; their directions, label sides, and terminal punctuation
% are clear at this resolution, so no higher-resolution escalation was needed.
% Transparent source dispositions:
% - typewritten letter-o special-fibre subscripts are normalized to zero;
% - the coefficient of the pulled-back extension in 7.3.1 is restored from
%   the printed bare G_m to the contextually required G_{m\eta};
% - the malformed phrase "si on des morphismes" is rendered "if we have".
% Hard stop before 7.3.2, source line 2316 / scan index 315.

\medskip
\noindent \SGARefTarget{sga7.sga7i.eviii.n7-3}{7.3}. 下面具体说明 (\SGARefLink{sga7.sga7i.eviii.n7-1-1}{7.1.1}) 和 (\SGARefLink{sga7.sga7i.eviii.n7-1-2}{7.1.2}) 中定义的障碍在 $P,Q,S$ 可变时的某些函子性, 以及它们之间的某些关系. 以下假设使这些障碍有定义的条件均成立. 验证没有困难, 留给读者.

\medskip
\noindent \SGARefTarget{sga7.sga7i.eviii.n7-3-1}{7.3.1}. 设
\[
 u:P'\longrightarrow P
\]
为 $S$ 上的光滑群概形的态射, $E_\eta$ 为 $P_\eta$ 以 $\underline G_{m\eta}$ 为核的扩张. 沿 $u_\eta:P'_\eta\to P_\eta$ 拉回, 得到 $P'_\eta$ 以 $\underline G_{m\eta}$ 为核的扩张 $E'_\eta$. 沿用这些记号, 下图交换:
\begin{equation}\tag{\SGARefTarget{sga7.sga7i.eviii.n7-3-1-1}{7.3.1.1}}
\begin{tikzcd}[row sep=large,column sep=huge]
 P'_0/{P'_0}^{\circ}
   \arrow[dd,"u_0"']
   \arrow[rd,"d(E'_\eta)"] & \\
 & (\mathbb Q/\mathbb Z)_k \\
 P_0/P_0^{\circ}
   \arrow[ru,"d(E_\eta)"'] &
\end{tikzcd}
\qquad .
\end{equation}

类似地, 若有 $S$ 上的光滑群概形的态射
\[
 u:P'\longrightarrow P,\qquad v:Q'\longrightarrow Q,
\]
以及 $(P_\eta,Q_\eta)$ 以 $\underline G_{m\eta}$ 为核的双扩张 $E_\eta$, 它沿 $(u_\eta,v_\eta)$ 的拉回是 $(P'_\eta,Q'_\eta)$ 以 $\underline G_{m\eta}$ 为核的双扩张 $E'_\eta$, 则下图交换:
\begin{equation}\tag{\SGARefTarget{sga7.sga7i.eviii.n7-3-1-2}{7.3.1.2}}
\begin{tikzcd}[row sep=large,column sep=huge]
 P'_0/{P'_0}^{\circ}\otimes Q'_0/{Q'_0}^{\circ}
   \arrow[dd,"u_0\otimes v_0"']
   \arrow[rd,"d(E'_\eta)"] & \\
 & (\mathbb Q/\mathbb Z)_k \\
 P_0/P_0^{\circ}\otimes Q_0/Q_0^{\circ}
   \arrow[ru,"d(E_\eta)"'] &
\end{tikzcd}
\end{equation}
这些相容性的验证本质上是机械的, 留给读者.

% END INCLUDED SGA7I_English_current_buildable_source_20260823/components/124_expose_VIII_section_7_3_1.tex

% BEGIN INCLUDED SGA7I_English_current_buildable_source_20260823/components/125_expose_VIII_sections_7_3_2_and_7_3_3.tex
% R3 explicit revisions R3-VIII-09, R3-VIII-10; see qa/FIDELITY-R3-CHANGES.json.
% SOURCE AUTHORITY:
% expose_VIII_body.tex lines 2316--2361; scan indices 315--317;
% source folios 304--306.
% Direct scan comparison: physical PDF pages 316--318 in overlapping
% 1100-dpi bands. The kernel formulas and the symmetry diagram were checked
% directly; all nodes, arrows, label sides, and the closing comma are clear.
% Transparent source dispositions:
% - typewritten letter-o special-fibre subscripts are normalized to zero;
% - the biextension denoted once by E rather than E_eta in 7.3.2 is rendered
%   consistently as the stated generic-fibre biextension E_eta;
% - "biextension symetrique" in the construction of {}^sE is rendered
%   "transposed biextension" to distinguish that operation from the later
%   property E_eta \simeq {}^sE_eta of being symmetric.
% Hard stop before 7.3.4, source line 2367 / scan index 317.

\medskip
\noindent \SGARefTarget{sga7.sga7i.eviii.n7-3-2}{7.3.2}. 由相容性 (\SGARefLink{sga7.sga7i.eviii.n7-3-1-1}{7.3.1.1}) 得出, 若 $E_\eta$ 为 $P_\eta$ 以 $\underline G_{m\eta}$ 为核的扩张, 则存在 $P$ 的最大群子概形 $V$, 使 $E_\eta$ 在 $V_\eta$ 上的限制可延拓为 $V$ 以 $\underline G_{mS}$ 为核的扩张. 它是 $P$ 中包含 $P_\eta$ 并由条件
\begin{equation}\tag{\SGARefTarget{sga7.sga7i.eviii.n7-3-2-1}{7.3.2.1}}
 V_0/V_0^{\circ}=\ker d(E_\eta).
\end{equation}
定义的开子群 $V$. 若此时 $u:P'\to P$ 为 $S$ 上光滑群概形的态射, 则将 $E_\eta$ 沿 $u_\eta$ 拉回得到的扩张 $E'_\eta$ 可延拓为 $P'$ 以 $\underline G_{mS}$ 为核的扩张, 当且仅当
\begin{equation}\tag{\SGARefTarget{sga7.sga7i.eviii.n7-3-2-2}{7.3.2.2}}
 u(P')\subset V.
\end{equation}

在 $(P_\eta,Q_\eta)$ 以 $\underline G_{m\eta}$ 为核的双扩张 $E_\eta$ 的情形, 同样存在 $P$ 的最大群子概形 $V$, 使 $E_\eta$ 在 $(V_\eta,Q_\eta)$ 上的限制可延拓为 $(V,Q)$ 以 $\underline G_{mS}$ 为核的双扩张. 它是 $P$ 中包含一般纤维 $P_\eta$ 的开子群, 其刻画为: $V_0/V_0^{\circ}$ 是配对 (\SGARefLink{sga7.sga7i.eviii.n7-1-2}{7.1.2}) 在 $\Phi$ 中的左零化子, 即相应的从 $\Phi$ 到 $\underline{\mathrm{Hom}}(\Psi,(\mathbb Q/\mathbb Z)_k)$ 的同态的核:
\begin{equation}\tag{\SGARefTarget{sga7.sga7i.eviii.n7-3-2-3}{7.3.2.3}}
 V_0/V_0^{\circ}
 =\ker\!\left(
   \Phi\xrightarrow{\ d(E_\eta)\ }
   \underline{\mathrm{Hom}}(\Psi,(\mathbb Q/\mathbb Z)_k)
 \right).
\end{equation}
若 $u:P'\to P$ 为 $S$ 上光滑群概形的态射, 则将 $E_\eta$ 沿 $u_\eta$ 拉回得到的 $(P'_\eta,Q_\eta)$ 以 $\underline G_{m\eta}$ 为核的双扩张, 可延拓为 $(P',Q)$ 以 $\underline G_{mS}$ 为核的双扩张, 当且仅当 (\SGARefLink{sga7.sga7i.eviii.n7-3-2-2}{7.3.2.2}) 成立. 交换 $P$ 和 $Q$ 的地位, 得到对称的结果.

当然, 一般不存在最大的子群\underline{对} $P'\to P$、$Q'\to Q$, 使 $E_\eta$ 在 $(P'_\eta,Q'_\eta)$ 以 $\underline G_{m\eta}$ 为核的双扩张上的限制 $E'_\eta$, 可延拓为 $(P',Q')$ 以 $\underline G_{mS}$ 为核的双扩张. 只能说, 对给定的一对子群 $(P',Q')$, 或更一般地, 对给定的一对光滑群概形态射 $P'\to P$、$Q'\to Q$, 所求延拓存在当且仅当形式 $d(E_\eta)$ 在
\[
 P'_0/{P'_0}^{\circ}\otimes Q'_0/{Q'_0}^{\circ}
\]
上的限制为零.

\medskip
\noindent \SGARefTarget{sga7.sga7i.eviii.n7-3-3}{7.3.3}. \underline{对称性}. 设 $E_\eta$ 为 $(P_\eta,Q_\eta)$ 以 $\underline G_{m\eta}$ 为核的双扩张, 考虑转置双扩张 ${}^{s}E_\eta$ (VII \SGARefLink{sga7.sga7i.evii.n2-7}{2.7}), 它是 $(Q_\eta,P_\eta)$ 以 $\underline G_{m\eta}$ 为核的双扩张. 下图交换:
\begin{equation}\tag{\SGARefTarget{sga7.sga7i.eviii.n7-3-3-1}{7.3.3.1}}
\begin{tikzcd}[row sep=large,column sep=huge]
 \Phi\otimes\Psi
   \arrow[dd,"\mathrm{sym}"']
   \arrow[rd,"d(E_\eta)"] & \\
 & (\mathbb Q/\mathbb Z)_k \\
 \Psi\otimes\Phi
   \arrow[ru,"d({}^{s}E_\eta)"'] &
\end{tikzcd}
\qquad ,
\end{equation}
其中竖直箭头为典范对称映射. (将此命题与 \SGARefLink{sga7.sga7i.eviii.n2-2-11}{2.2.11} 的\underline{反对称性}命题比较.) 特别地可知, 若 $P=Q$, 且 $E_\eta$ 为 $(P_\eta,Q_\eta)$ 以 $\underline G_{m\eta}$ 为核的\underline{对称}双扩张, 即 $E_\eta$ 同构于 ${}^{s}E_\eta$, 则 $\Phi\times\Phi$ 上取值于 $(\mathbb Q/\mathbb Z)_k$ 的双线性形式 $d(E_\eta)$ 是\underline{对称的} (而非像 \SGARefLink{sga7.sga7i.eviii.n2-2-12}{2.2.12} 中那样是交错的!). 若它为零, 即若 $E_\eta$ 延拓为 $(P,P)$ 以 $\underline G_{mS}$ 为核的双扩张 $E$, 则后者也必对称: 根据 \SGARefLink{sga7.sga7i.eviii.n7-1}{7.1}(b), 每个同构 ${}^{s}E_\eta\to E_\eta$ 都唯一延拓为同构 ${}^{s}E\to E$.

% END INCLUDED SGA7I_English_current_buildable_source_20260823/components/125_expose_VIII_sections_7_3_2_and_7_3_3.tex

% BEGIN INCLUDED SGA7I_English_current_buildable_source_20260823/components/126_expose_VIII_sections_7_3_4_and_7_3_5.tex
% R3 explicit revisions R3-VIII-11; see qa/FIDELITY-R3-CHANGES.json.
% SOURCE AUTHORITY:
% expose_VIII_body.tex lines 2367--2446; scan indices 317--321;
% source folios 306--310.
% Direct scan comparison: physical PDF pages 318--322 in overlapping
% 1100-dpi bands. Formulas, primes, residue-field subscripts, and the
% ramification index are clear at this resolution; no higher escalation was
% needed.
% Transparent source dispositions:
% - formula (7.3.4.1) is restored from the transcription's impossible E_*
%   to the source scan's q_eta^*(E_eta);
% - the family P_{eta Q_eta} is written explicitly as
%   P_eta x_eta Q_eta;
% - typewritten letter-o special-fibre subscripts are normalized to zero;
% - the malformed phrase "a l'aise du changement de base" is rendered
%   "by base change".
% Hard stop before Proposition 7.4, source line 2449 / scan index 321.

\medskip
\noindent \SGARefTarget{sga7.sga7i.eviii.n7-3-4}{7.3.4}. \underline{将 (\SGARefLink{sga7.sga7i.eviii.n7-1-2}{7.1.2}) 型障碍归结为 (\SGARefLink{sga7.sga7i.eviii.n7-1-1}{7.1.1}) 型障碍.}

为明确起见, 假设存在整数 $n>0$ 使 $n\Phi=0$, 并设
\[
 q\in Q(S)
\]
是 $Q$ 的截面. 稍滥用记号, 用
\[
 q_0\in\Psi(k)
\]
表示 $q$ 在 $s$ 的值在 $\Psi(k)$ 中的像. 将 $E_\eta$ 视为 $Q_\eta$ 上 $P_\eta\times_\eta Q_\eta$ 以 $\underline G_{mQ_\eta}$ 为核的扩张, 拉回得到扩张
\[
 q_\eta^*(E_\eta),\qquad
 \text{$P_\eta$ 以 $\underline G_{m\eta}$ 为核的扩张},
\]
从而得到障碍
\begin{equation}\tag{\SGARefTarget{sga7.sga7i.eviii.n7-3-4-1}{7.3.4.1}}
 d_q\overset{\mathrm{dfn}}{=}
 d\bigl(q_\eta^*(E_\eta)\bigr):
 \Phi\longrightarrow(\mathbb Q/\mathbb Z)_k.
\end{equation}
记 $d=d(E_\eta)$, 则有关系
\begin{equation}\tag{\SGARefTarget{sga7.sga7i.eviii.n7-3-4-2}{7.3.4.2}}
 d_q(x)=d(x,q_0),
\end{equation}
它对 $\Phi$ 在 $k$-概形 $T$ 中取值的每个点 $x$ 都成立.

\smallskip
\noindent \SGARefTarget{sga7.sga7i.eviii.n7-3-4-3}{7.3.4.3}. 对每个 $q\in Q(S)$, 此公式用配对 $d$ 来确定 $d_q$; 反过来, 至少当 $\Psi$ 的每个点都来自 $Q$ 的截面时, 它显然可从不变量 $d_q$ 恢复 $d$. 根据 Hensel 引理, 当 $S$ 严格局部时此条件自动成立. 配对 $d$ 只须在 $k$ 的可分闭包上已知, 就完全确定, 而由下面的 \SGARefLink{sga7.sga7i.eviii.n7-3-5-3}{7.3.5.3}, 它的形成与 $S$ 上的平展局部化交换. 因此, 必要时将 $S$ 换为其严格局部化, $d$ 的确定就归结为不变量 $d_q$ 的确定.

特别地, 当 $\Psi$ 的每个点都来自 $Q$ 的截面 $q$ 时, 配对 $d$ 为零 (即 $E_\eta$ 延拓为 $(P,Q)$ 以 $\underline G_m$ 为核的双扩张 $E$), 当且仅当对 $Q$ 的每个截面 $q$, $d_q$ 都为零, 也即当且仅当 $q_\eta^*(E_\eta)$ 延拓为 $P$ 以 $\underline G_m$ 为核的扩张. 不对 $\Psi$ 作任何假设时, 只要允许 $q$ 为在局部平展 $S$-概形 $S'$ 中取值的点, 类似的判别法仍然成立.

\smallskip
\noindent \SGARefTarget{sga7.sga7i.eviii.n7-3-4-4}{7.3.4.4}. 反过来, 从 \SGARefLink{sga7.sga7i.eviii.n7-1}{7.1}(a) 的情形出发, 可将障碍 (\SGARefLink{sga7.sga7i.eviii.n7-1-1}{7.1.1}) 解释为 (\SGARefLink{sga7.sga7i.eviii.n7-1-2}{7.1.2}) 型障碍: 取 $Q=\mathbb Z_S$ 以及从 $E_\eta$ 推出的 $(P_\eta,\mathbb Z_\eta)$ 的双扩张. 事实上, 只须在上述构造中取 $q$ 等于 $\mathbb Z_S$ 的截面 $1$ 即可.

\medskip
\noindent \SGARefTarget{sga7.sga7i.eviii.n7-3-5}{7.3.5}. \underline{基变换 $S'\to S$ 的影响.} 假设给定性状的支配局部态射
\[
 f:S'\longrightarrow S,
\]
用 $\eta'$ 表示 $S'$ 的一般点, $s'$ 表示其闭点, 并令 $k'=k(s')$. 若 $S=\operatorname{Spec}(V)$ 且 $S'=\operatorname{Spec}(V')$, 则 $V'$ 是支配离散赋值环 $V$ 的离散赋值环. 令
\begin{equation}\tag{\SGARefTarget{sga7.sga7i.eviii.n7-3-5-1}{7.3.5.1}}
 e(S'/S)=e(V'/V)=\operatorname{length}_{V'}(V'/\mathfrak m V'),
\end{equation}
其中 $\mathfrak m$ 为 $V$ 的极大理想.

现在设 $P$ 如 \SGARefLink{sga7.sga7i.eviii.n7-1}{7.1}(a) 所述, $E_\eta$ 为 $P_\eta$ 以 $\underline G_{m\eta}$ 为核的扩张, 令 $P'=P\times_S S'$, 并设 $E'_{\eta'}$ 为拉回 $E_\eta$ 得到的 $P'_{\eta'}$ 以 $\underline G_{m\eta'}$ 为核的扩张. 沿用 \SGARefLink{sga7.sga7i.eviii.n7-1}{7.1}(a) 的记号, 有典范同构
\[
 \Phi'\overset{\mathrm{dfn}}{=}P'_0/{P'_0}^{\circ}
 \simeq\Phi\otimes_k k'\simeq\Phi_{k'},
\]
因而可以考虑
\[
 d(E_\eta)_{k'}:\Phi'\longrightarrow(\mathbb Q/\mathbb Z)_{k'}.
\]
则
\begin{equation}\tag{\SGARefTarget{sga7.sga7i.eviii.n7-3-5-2}{7.3.5.2}}
 d(E'_{\eta'})=e(S'/S)\,d(E_\eta)_{k'}.
\end{equation}

类似地, 在 \SGARefLink{sga7.sga7i.eviii.n7-1}{7.1}(b) 的假设下, 使用显然的记号, (\SGARefLink{sga7.sga7i.eviii.n7-3-5-2}{7.3.5.2}) 仍然成立; 此时 $d(E_\eta)$ 表示配对 (\SGARefLink{sga7.sga7i.eviii.n7-1-2}{7.1.2}), $d(E_\eta)_{k'}$ 表示由它通过基变换 $k\to k'$ 得到的配对
\[
 \Phi'\otimes\Psi'\longrightarrow(\mathbb Q/\mathbb Z)_{k'}
\]
而 $E'_{\eta'}$ 表示从 $(P_\eta,Q_\eta)$ 以 $\underline G_{m\eta}$ 为核的双扩张 $E_\eta$, 经基变换 $\eta'\to\eta$ 得到的 $(P'_{\eta'},Q'_{\eta'})$ 以 $\underline G_{m\eta'}$ 为核的双扩张.

\smallskip
\noindent \SGARefTarget{sga7.sga7i.eviii.n7-3-5-3}{7.3.5.3}. 一个有趣的情形是 $e(S'/S)=1$, 即
\[
 \mathfrak m'=\mathfrak m V',
\]
其中 $\mathfrak m'$ 表示 $V'$ 的极大理想. 特别地, 若 $V'$ 在 $V$ 上平展, 或者若 $V'$ 是 $V$ 的 Hensel 化、严格 Hensel 化或完备化, 就有此情形. 这时 (\SGARefLink{sga7.sga7i.eviii.n7-3-5-2}{7.3.5.2}) 只是说, \underline{障碍 (\SGARefLink{sga7.sga7i.eviii.n7-1-1}{7.1.1}) 或 (\SGARefLink{sga7.sga7i.eviii.n7-1-2}{7.1.2}) 的形成与所讨论的基变换 $S'\to S$ 交换}.

\smallskip
\noindent \SGARefTarget{sga7.sga7i.eviii.n7-3-5-4}{7.3.5.4}. 另一方面, 一般情形的 (\SGARefLink{sga7.sga7i.eviii.n7-3-5-2}{7.3.5.2}) 表明, 对每个如 \SGARefLink{sga7.sga7i.eviii.n7-1}{7.1}(a) 中的扩张 $E_\eta$ (分别地, 如 \SGARefLink{sga7.sga7i.eviii.n7-1}{7.1}(b) 中的双扩张 $E_\eta$), 存在性状的一个\underline{有限}态射 $S'\to S$, 使 $k(\eta')$ 为 $k(\eta)$ 的可分扩张, 并使从 $E_\eta$ 经基变换 $\eta'\to\eta$ 得到的扩张 (分别地, 双扩张) $E'_{\eta'}$, 延拓为整个 $S'$ 上的扩张 (分别地, 双扩张) $E'$. 事实上, 只须选取 $S'/S$ 使 $e(S'/S)$ 为 \SGARefLink{sga7.sga7i.eviii.n7-1}{7.1} 中出现的整数 $n>0$ 的倍数, 而总能找到这样的扩张使 $k(\eta')$ 在 $k(\eta)$ 上可分.

% END INCLUDED SGA7I_English_current_buildable_source_20260823/components/126_expose_VIII_sections_7_3_4_and_7_3_5.tex

% BEGIN INCLUDED SGA7I_English_current_buildable_source_20260823/components/127_expose_VIII_proposition_7_4_corollary_7_5_and_bibliography.tex
% SOURCE AUTHORITY:
% expose_VIII_body.tex lines 2449--2495; scan indices 321--323;
% source folios 310--312.
% Direct scan comparison: physical PDF pages 322--324 in overlapping
% 1100-dpi bands. The category names, product bases, superscripts, theorem
% underlining, and bibliography were all clear; no higher-resolution
% escalation was needed.
% Transparent source dispositions:
% - source spacing around Z-algebra is normalized without changing meaning;
% - bibliographic accents and punctuation follow the printed authority where
%   the transcription contained obvious keyboard errors.
% This component closes Expose VIII at the end of its bibliography.

\medskip
\noindent\underline{命题 \SGARefTarget{sga7.sga7i.eviii.n7-4}{7.4}. 设 $S$ 为约化概形.}

\smallskip
\noindent a) \underline{设 $P$ 为 $S$ 上的光滑群概形, 其在 $S$ 的极大点上的纤维连通. 则从 $P$ 以 $\underline G_{mS}$ 为核的扩张范畴, 到沿 $P$ 的单位截面刚化的 $\underline G_{mP}$ 下的挠子范畴的自然函子}
\begin{equation}\tag{\SGARefTarget{sga7.sga7i.eviii.n7-4-1}{7.4.1}}
 \mathrm{EXT}(P,\underline G_{mS})
 \longrightarrow
 \mathrm{TORSRIG}(P,e_P;\underline G_{mS})
\end{equation}
\underline{是全忠实的. 若 $S$ 正规且 $P$ 的纤维连通, 则 (\SGARefLink{sga7.sga7i.eviii.n7-4-1}{7.4.1}) 的目标范畴中的对象 $E$ 属于此函子的本质像, 当且仅当对 $S$ 的每个极大点 $\eta$, $P_\eta$ 上的刚化挠子 $E_\eta$ 来自 $P_\eta$ 以 $\underline G_{m\eta}$ 为核的扩张.}

\smallskip
\noindent b) \underline{设 $P,Q$ 是 $S$ 上的两个光滑群概形, 其在 $S$ 的极大点上的纤维连通. 则从 $(P,Q)$ 以 $\underline G_{mS}$ 为核的双扩张范畴, 到沿 $P$ 和 $Q$ 的单位截面对作双刚化的 $\underline G_{m,P\times_S Q}$ 下的挠子范畴的自然函子}
\begin{equation}\tag{\SGARefTarget{sga7.sga7i.eviii.n7-4-2}{7.4.2}}
 \mathrm{BIEXT}(P,Q;\underline G_{mS})
 \longrightarrow
 \mathrm{TORSBIRIG}(P,Q;\underline G_{mS})
\end{equation}
\underline{是全忠实的. 若 $S$ 正规, 且 $P$ 或 $Q$ 的纤维连通, 则 (\SGARefLink{sga7.sga7i.eviii.n7-4-2}{7.4.2}) 的目标范畴中的对象 $E$ 属于此函子的本质像, 当且仅当对 $S$ 的每个极大点 $\eta$, $E_\eta$ 来自 $(P_\eta,Q_\eta)$ 以 $\underline G_{m\eta}$ 为核的双扩张.}

我们证明 (a); (b) 的证明本质上相同, 留给读者. 关于全忠实性的断言, 注意如果 $X$ 和 $Y$ 是约化概形 $S$ 上的平坦概形, 且其在极大点上的纤维满足 \SGARefLink{sga7.sga7i.eviii.n4-1}{4.1} 所述条件, 则 \SGARefLink{sga7.sga7i.eviii.n4-1}{4.1} 仍然成立. 事实上, $X\times_S Y$ (它在约化概形 $S$ 上平坦) 的极大纤维族在 $X\times_S Y$ 中概形稠密. 因而从 $X\times_S Y$ 到分离 $S$-概形 $\underline G_{mS}$ 的态射, 只要在极大纤维上平凡, 就是平凡的. 于是与 \SGARefLink{sga7.sga7i.eviii.n4-2}{4.2} 中一样证明了 (\SGARefLink{sga7.sga7i.eviii.n7-4-1}{7.4.1}) 的全忠实性, 并且那里给出的对本质像的刻画也仍然成立.

还须证明: 假设 $S$ 正规且 $P$ 的纤维连通, 若 $P$ 上的刚化挠子 $E$ 满足每个 $E_\eta$ 均来自扩张, 则 $E$ 也来自扩张. 后一个条件蕴含 $P$ 在 $S$ 上有限表现 (SGA 3 $\mathrm{VI}_B$ 5.5). 全忠实性还表明, 这个问题在 $S$ 上是局部的, 因而可以假设 $S$ 仿射且局部. 利用 Nagata 关于有限生成约化 $\mathbb Z$-代数的整闭包有限性的定理 (EGA IV 7.8), 通常的极限论证进一步将问题归结为 $S$ 为诺特概形的情形.

此时 \SGARefLink{sga7.sga7i.eviii.n7-1}{7.1}(a) 表明, 存在一个补集余维至少为 $2$ 的开子集 $U\subset S$, 使 $E_U$ 来自 $P_U$ 以 $\underline G_{mU}$ 为核的扩张. 换言之, 所求同构
\[
 \pi^*(E)\simeq
 \operatorname{pr}_1^*(E)\operatorname{pr}_2^*(E)
\]
可以在 $(P\times_S P)_U$ 上构造. 这是正规概形 $P\times_S P$ 的开子集, 其补集余维至少为 $2$. 因而 $(P\times_S P)_U$ 上的同构可延拓到整个 $P\times_S P$, 证明完毕.

\medskip
\noindent\underline{推论 \SGARefTarget{sga7.sga7i.eviii.n7-5}{7.5}. 设 $S$ 为正规概形, $P$ 和 $Q$ 为 $S$ 上的两个光滑群概形, 其在极大点上的纤维连通, 且 $P$ 或 $Q$ 的纤维连通. 再假设对 $S$ 的每个极大点 $\eta$, $P_\eta$ 或 $Q_\eta$ 是阿贝尔概形以环面为核的扩张 (或更一般地, 具有因子为阿贝尔概形、环面或群 $\underline G_a$ 的合成列). 则函子 (\SGARefLink{sga7.sga7i.eviii.n7-4-2}{7.4.2}) 为范畴等价.}

这是 \SGARefLink{sga7.sga7i.eviii.n7-4}{7.4}(b) 和 \SGARefLink{sga7.sga7i.eviii.n4-7}{4.7} 的直接推论.

\bigskip
\noindent\underline{参考文献}

\begin{enumerate}
\item[{[1]}] M. Demazure, J. Giraud, M. Raynaud,
  \emph{Sch\'emas ab\'eliens}, S\'eminaire Orsay 1967/68, 即将出版.
\item[{[2]}] M. Raynaud, \emph{Faisceaux amples sur les sch\'emas en groupes
  et les espaces homog\`enes}, Paris 学位论文, 1968 (即将在 Springer 的 Lecture
  Notes 丛书出版).
\item[{[3]}] M. Rosenlicht, \emph{Some basic theorems on algebraic groups},
  Amer. Jour. Math. \textbf{78}, 第~2 期 (1956), 第~401--443 页.
\item[{[3bis]}] M. Rosenlicht, \emph{Some rationality questions on algebraic
  groups}, Annali di Matematica, 第 4 辑, 第~43 期 (1957), 第~25--50 页.
\item[{[4]}] J.-P. Serre, \emph{Corps locaux}, Act. Sci. Ind. 1296 (1962),
  Hermann, Paris.
\end{enumerate}

% END INCLUDED SGA7I_English_current_buildable_source_20260823/components/127_expose_VIII_proposition_7_4_corollary_7_5_and_bibliography.tex
